% !TEX program = xelatex
% !TEX encoding = UTF-8 Unicode
% Generated by scripts/convert_mdf_master.py from the MDF 1.9a source; do not edit.
\documentclass[10pt,openany]{book}
\usepackage{fontspec}
\usepackage{geometry}
\usepackage{xcolor}
\usepackage{fancyvrb}
\usepackage{fvextra}
\usepackage{hyperref}
\geometry{margin=25mm}
\definecolor{cflink}{RGB}{0,110,0}
\definecolor{cfmissing}{RGB}{170,0,0}
\definecolor{exlabel}{RGB}{110,110,110}
\IfFontExistsTF{Noto Serif}{\setmainfont{Noto Serif}}{%
\IfFontExistsTF{Gentium Book Plus}{\setmainfont{Gentium Book Plus}}{%
\IfFontExistsTF{Gentium Plus}{\setmainfont{Gentium Plus}}{%
\IfFontExistsTF{Gentium}{\setmainfont{Gentium}}{\setmainfont{Latin Modern Roman}}}}}
\IfFontExistsTF{Noto Sans Mono}{\setmonofont{Noto Sans Mono}}{%
\IfFontExistsTF{DejaVu Sans Mono}{\setmonofont{DejaVu Sans Mono}}{\setmonofont{Latin Modern Mono}}}
\setcounter{secnumdepth}{-2}
\setcounter{tocdepth}{0}
\pagestyle{headings}
\hypersetup{hidelinks,bookmarksnumbered=false,bookmarksdepth=2,
pdftitle={MDF Lexical Fields},pdfsubject={MDF 1.9a field documentation}}
\begin{document}
\frontmatter
\begin{titlepage}
\centering
{\Huge MDF Lexical Fields\par}
\vspace{1.5em}
{\large Version v3.0\par}
\vspace{0.5em}
{\large 2026-09-30\par}
\vspace{2em}
{\small Generated from \texttt{docs/mdf/MDFields19a\_UTF8.txt} (MDF 1.9a field documentation).\par}
\end{titlepage}
\clearpage
\phantomsection
\addcontentsline{toc}{chapter}{Foreword}
\chapter*{Foreword}
This reference presents the MDF 1.9a field documentation in a printable, browsable form. The source document is a Toolbox/Shoebox help file — a flat collection of records designed for on-screen keyword search rather than linear reading. The source identifies itself as draft version 1.9a (May 12, 2006); the original database was constructed by David Coward and the revisions by Karen Buseman.

This edition was assembled with the assistance of AI. An automated conversion pipeline (a Python parser, XeLaTeX typesetting, and static HTML generation), designed and implemented by an AI agent under human direction and review, transformed the source into the present formats. All field content is preserved from the source without alteration; the AI's role was structural and typographic, not editorial. No content was invented, paraphrased, or omitted.

Changes made relative to the source:
\begin{enumerate}
  \item Table of contents. The source has no print-style table of contents; its navigation is a cross-reference table on the home record plus keyword search. This edition adds a proper ToC. The chapter order follows the source's own recommended reading order: the topics listed on the home record, in the order it presents them.
  \item Unified field reference. The field-marker records, which the source interleaves alphabetically for on-screen lookup, are collected into a single Field Marker Reference section, grouped by the source's own field classifications — the record marker, Basic fields, Reserved fields, Optional fields (as defined in the source's Introduction), and Discontinued markers — alphabetical within each group. One obsolete-marker record is placed with the discussion of old and changed markers, as its own cross-reference directs.
  \item Cross-references (``See also'') render as live links, set in green following the source's own stated convention.
  \item Formatting and printing examples are reproduced verbatim in monospaced blocks.
  \item Typography, page layout, bookmarks, pagination, and typographic quotation marks are new to this edition; the text itself is unchanged. Formatting examples keep the source's literal straight quotes, as they depict exact database input.
\end{enumerate}

Typeset with XeLaTeX (Noto Serif, with Gentium). HTML edition generated from the same parsed source.

Conversion of 2026-10-06. Assembled with AI assistance — OpenCode (huggingface/zai-org/GLM-5.3-Flash), directed by Michael Conrad.

\tableofcontents
\mainmatter
\chapter{Helps Database for MDF Marker Set}\label{key:aa}

Most records in this file are keyed to the two- or three-letter field marker codes, standardized by the Multi-Dictionary Formatter (MDF) program for lexical databases. To read about a particular field, if you know the marker, simply search for the letter code of that marker, e.g. for the English gloss field (\textbackslash{}ge), search for `ge'. If a field is listed in a discussion, right-clicking on the field marker (option-click on the Macintosh) will jump you to that marker's topic. If you don't know the marker, read the organized grouping of markers listed in the Introduction. You can get there by searching or by dong a right click on the green word in the next line.


{\color{cflink} \hyperref[key:Introduction]{Introduction}}


Other pertinent topics are covered as well. Look for the green text (or the cf Cross Reference field) To read these discussions simply search for them (or jump to by doing a right-click) using the following labels:


{\color{cflink} \hyperref[key:Introduction]{Introduction} \hyperref[key:Range_Sets]{Range\_Sets}}


{\color{cflink} \hyperref[key:Order_of_Fields]{Order\_of\_Fields} \hyperref[key:Character_Style_Codes]{Character\_Style\_Codes}}


{\color{cflink} \hyperref[key:Summary_of_Fields]{Summary\_of\_Fields} \hyperref[key:Punctuation_and_Special_Codes]{Punctuation\_and\_Special\_Codes}}


{\color{cflink} \hyperref[key:Sections_in_a_Lexical_Entry]{Sections\_in\_a\_Lexical\_Entry} \hyperref[key:Printed_Field_Labels]{Printed\_Field\_Labels}}


{\color{cflink} \hyperref[key:Using_Subentries_or_Lexical_Entries]{Using\_Subentries\_or\_Lexical\_Entries} \hyperref[key:Unknown_Fields]{Unknown\_Fields}}


{\color{cflink} \hyperref[key:Alternate_Hierarchy]{Alternate\_Hierarchy} \hyperref[key:Formatting_and_Printing]{Formatting\_and\_Printing}}


{\color{cflink} \hyperref[key:Free-form_Fields]{Free-form\_Fields} \hyperref[key:References]{References}}


{\color{cflink} \hyperref[key:The_MDF_Documentation]{The\_MDF\_Documentation} \hyperref[key:Old_and_Changed_Markers]{Old\_and\_Changed\_Markers}}


{\color{cflink} \hyperref[key:When_MDF_Fails_to_Meet_Your_Requirements]{When\_MDF\_Fails\_to\_Meet\_Your\_Requirements}}


(This is draft version 1.9a, in preparation for ``version 2'' of the MDFields Helps database. The original database was constructed by David Coward. The revisions are being done by Karen Buseman. May 12, 2006)


\chapter{Standard Lexical Database Field Markers (for MDF)}\label{key:Introduction}

\section{Introduction}

MDF is a set of markers designed for multi-lingual dictionaries. It was created by Grimes and Coward and documented in, ``Making Dictionaries''. MDF is also a printing (export) process. Because this printing process was built into the Windows version of Shoebox (now Toolbox), both have spread -- each feeding the popularity of the other. The MDF marker set is now used by linguists worldwide. For further information on the MDF manual itself, see


{\color{cflink} \hyperref[key:The_MDF_Documentation]{The\_MDF\_Documentation}}


This file is designed to be an MDF help file available in Toolbox while you work on your lexicon. Most of the information in this file describes the standardized MDF field markers and their contents (i.e. what kind of information the markers mark).


To look up the description of a field marker, use the Toolbox Search command and type the field marker code, (without the backslash  \textbackslash{} character), e.g. for information on the \textbackslash{}ps field, search for `ps' (no quotes).


Another way to look up a marker is to click the Browse button and scroll down the list of markers until you see what you're interested in. Browse is also one way to see the titles of information available on other topics. Remember, too, that the ``aa'' record contains a full list of other topics available.


MDF is not perfect and does not fit everyone's needs. If you need fields that aren't available, or if the printout is not what you want, read the following:


{\color{cflink} \hyperref[key:When_MDF_Fails_to_Meet_Your_Requirements]{When\_MDF\_Fails\_to\_Meet\_Your\_Requirements}}


\section{Organizing the Markers}

One problem faced by users of MDF is that of ``too many markers''. The following summary is an attempt to make the number seem less threatening.


\subsection{... by Language}

A large number of markers come in sets of three or four. MDF attemps to generalize  the need for multiple languages by recognizing four possible languages and providing markers for them. Most of the MDF markers are two characters long. The second character is often (not always) a language character.


\_e   indicates English language markers


\_n   indicates National language markers


\_r   indicates Regional language markers


\_v   indicates Vernacular language markers


Some markers have the language assumed. So the lexeme (lx) and subentries (se) are always vernacular language, even though they don't end with ``v'' (and even though subentry ends with ``e''). Some of these, like the variant (va), have related markers (comments in this case) for the other three languages.


\subsection{... by Function}

In a real sense, everything in the entry relates to the headword. But some fields are more tightly related than others. Some are more related to the definition, some are your personal notes on different topics, etc. This is an attempt to organize the approximately 100 MDF markers.


\begin{quote}\small\itshape Note: The following does not give the order in which these fields normally appear in a dictionary.\end{quote}

\begin{quote}\small\itshape Note: A marker like de* means that there is a set of  language markers beginning with d: dv, dn, dr, de. A right click on such a marker will take you to the one, and the others will be listed there.\end{quote}

\subsubsection{Information Relating Directly to the Headword}

\begin{description}
\item[{\hyperref[key:lx]{\textcolor{cflink}{lx}}}] lexeme
\item[{\hyperref[key:hm]{\textcolor{cflink}{hm}}}] homonym number
\item[{\hyperref[key:lc]{\textcolor{cflink}{lc}}}] lexical citation form
\item[{\hyperref[key:ph]{\textcolor{cflink}{ph}}}] phonetic (pronunciation)
\item[{\hyperref[key:ps]{\textcolor{cflink}{ps}}}] part of speech
\item[{\hyperref[key:pn]{\textcolor{cflink}{pn}}}] National part of speech
\item[{\hyperref[key:sn]{\textcolor{cflink}{sn}}}] sense number
\end{description}


\subsubsection{Information about the Headword}

\begin{description}
\item[{\hyperref[key:va]{\textcolor{cflink}{va}}}] variant form (also \hyperref[key:ve]{\textcolor{cflink}{ve*}} comments)
\item[{\hyperref[key:mr]{\textcolor{cflink}{mr}}}] morphology
\item[{\hyperref[key:lt]{\textcolor{cflink}{lt}}}] literally
\item[{\hyperref[key:sc]{\textcolor{cflink}{sc}}}] scientific name
\item[{\hyperref[key:sd]{\textcolor{cflink}{sd}}}] semantic domain
\item[{\hyperref[key:is]{\textcolor{cflink}{is}}}] index of semantics
\item[{\hyperref[key:th]{\textcolor{cflink}{th}}}] thesaurus
\item[{\hyperref[key:pd]{\textcolor{cflink}{pd}}}] paradigm class
\item[{\hyperref[key:pdl]{\textcolor{cflink}{pdl}}}] paradigm label (also \hyperref[key:pdv]{\textcolor{cflink}{pdv*}} paradigm form \& glosses)
\end{description}


\subsubsection{Other forms of the Headword}

\begin{description}
\item[{\hyperref[key:va]{\textcolor{cflink}{va}}}] variant form (also \hyperref[key:ve]{\textcolor{cflink}{ve*}} comments)
\item[{\hyperref[key:mr]{\textcolor{cflink}{mr}}}] morphology
\item[{\hyperref[key:pd]{\textcolor{cflink}{pd}}}] paradigm class
\item[{\hyperref[key:pdl]{\textcolor{cflink}{pdl}}}] paradigm label (also \hyperref[key:pdv]{\textcolor{cflink}{pdv*}} paradigm form \& glosses)
\end{description}


\subsubsection{Related entries}

\begin{description}
\item[{\hyperref[key:cf]{\textcolor{cflink}{cf}}}] cross reference
\item[{\hyperref[key:mn]{\textcolor{cflink}{mn}}}] main-entry cross reference
\item[{\hyperref[key:lf]{\textcolor{cflink}{lf}}}] lexical function (like synonym, causal, generic, deverbal noun, locative, etc.) (also \hyperref[key:lv]{\textcolor{cflink}{lv*}} lexical function form and glosses)
\end{description}


\subsubsection{Origins}

\begin{description}
\item[{\hyperref[key:bw]{\textcolor{cflink}{bw}}}] borrowed word
\item[{\hyperref[key:et]{\textcolor{cflink}{et}}}] etymology (also \hyperref[key:eg]{\textcolor{cflink}{eg}} gloss, \hyperref[key:es]{\textcolor{cflink}{es}} source, \hyperref[key:ec]{\textcolor{cflink}{ec}} comment)
\end{description}


\subsubsection{Definition and Range of Meaning}

\begin{description}
\item[{\hyperref[key:de]{\textcolor{cflink}{de*}}}] definition
\item[{\hyperref[key:ee]{\textcolor{cflink}{ee*}}}] encyclopedic information (expanded definition)
\item[{\hyperref[key:ue]{\textcolor{cflink}{ue*}}}] usage
\item[{\hyperref[key:oe]{\textcolor{cflink}{oe*}}}] ``only'' (restrictions)
\item[{\hyperref[key:ge]{\textcolor{cflink}{ge*}}}] gloss (for interlinear)
\item[{\hyperref[key:re]{\textcolor{cflink}{re*}}}] reversal (for finder lists)
\item[{\hyperref[key:we]{\textcolor{cflink}{we*}}}] word-level gloss
\end{description}


\subsubsection{Examples}

\begin{description}
\item[{\hyperref[key:rf]{\textcolor{cflink}{rf}}}] source of example (reference)
\item[{\hyperref[key:xv]{\textcolor{cflink}{xv*}}}] examples and translations
\item[{\hyperref[key:pc]{\textcolor{cflink}{pc}}}] picture snd sound vid video
\end{description}


\subsubsection{Dialects}

\begin{quote}\small\itshape There are no specifically ``dialect'' markers. Various alternative approaches can be used. The revision of Chapter 2 discusses the following as possibilities.\end{quote}

\begin{description}
\item[{\hyperref[key:va]{\textcolor{cflink}{va*}}}] (variant)
\item[{\hyperref[key:ue]{\textcolor{cflink}{ue*}}}] (usage)
\item[{\hyperref[key:oe]{\textcolor{cflink}{oe*}}}] (only / restrictions)
\item[{\hyperref[key:lf]{\textcolor{cflink}{lf}}}] SynD (lexical function: synonym dialect)
\end{description}


\subsubsection{Miscellany}

\begin{description}
\item[{\hyperref[key:tb]{\textcolor{cflink}{tb}}}] table (caption)
\item[{\hyperref[key:bb]{\textcolor{cflink}{bb}}}] bibliography
\item[{\hyperref[key:nt]{\textcolor{cflink}{nt}}}] notes (various kinds)
\item[{\hyperref[key:so]{\textcolor{cflink}{so}}}] source of data
\item[{\hyperref[key:st]{\textcolor{cflink}{st}}}] status
\item[{\hyperref[key:dt]{\textcolor{cflink}{dt}}}] date last edited
\end{description}


If you are unsure which marker to use, first look under `Summary\_of\_Fields' for the code you want, and then search for that field marker.


\section{Three Types of Field Markers}

In this database, each field marker is describe by of the following categories:


{\small <Basic>}\par

This marks all of the ``basic'' field markers. These are the field markers which will be useful to the most people. Consider including all of the basic field markers in your Toolbox lexical database template, so they will be inserted automatically into each of your new records. Leaving a basic field empty is okay.


{\small <Reserved>}\par

This marks the ``reserved'' fields, meaning, these fields serve a crucial function in the MDF formatting process and so should be used only if needed. These are the \textbackslash{}lc, \textbackslash{}hm, \textbackslash{}se, and  \textbackslash{}sn fields. Empty reserved fields should not be included in your records, i.e. if present, they should contain data.


{\small <Optional>}\par

This marks the ``optional'' fields. These are the more technical fields given to cover just about every topic. You may want to include some of these in your database template as well. And no doubt, some of these optional fields will be of no use to you. Like with basic fields, leaving an optional field empty is okay.


\section{Multiple Fields}

The number of occurrences of any given field is not restricted. But in printing through the standard MDF export process, multiple occurrences of a given field will be concatenated within a hierarchical section, e.g. with a sense, or if there is no sense, then within a part of speech section, or subentry section, etc. Multiple occurrences within the same record but in separate sections (sense, part of speech, or subentry) are kept separate.


To continue this discussion, jump to:


{\color{cflink} \hyperref[key:Order_of_Fields]{Order\_of\_Fields}}


\chapter{Range Sets}\label{key:Range_Sets}

Through the Range Set feature, Toolbox provides you with a powerful way to maintain consistent labeling on any field where you are using a relatively small set of categorizing codes. The most common fields for this are:


\begin{description}
\item[{\hyperref[key:ps]{\textcolor{cflink}{\textbackslash{}ps}}}] part of speech
\item[{\hyperref[key:pn]{\textcolor{cflink}{\textbackslash{}pn}}}] part of speech (national)
\item[{\hyperref[key:pd]{\textcolor{cflink}{\textbackslash{}pd}}}] paradigm set
\item[{\hyperref[key:sd]{\textcolor{cflink}{\textbackslash{}sd}}}] semantic domain
\item[{\hyperref[key:is]{\textcolor{cflink}{\textbackslash{}is}}}] index of semantics
\item[{\hyperref[key:th]{\textcolor{cflink}{\textbackslash{}th}}}] thesaurus
\end{description}


You are strongly encouraged to use Range Sets for these fields also:


\begin{description}
\item[{\hyperref[key:lf]{\textcolor{cflink}{\textbackslash{}lf}}}] lexical function
\item[{\hyperref[key:pdl]{\textcolor{cflink}{\textbackslash{}pdl}}}] paradigm form label
\end{description}


Here are other fields which could benefit from the Range Set feature:


\begin{description}
\item[{\hyperref[key:hm]{\textcolor{cflink}{\textbackslash{}hm}}}] homonym number
\item[{\hyperref[key:sn]{\textcolor{cflink}{\textbackslash{}sn}}}] sense number
\item[{\hyperref[key:bw]{\textcolor{cflink}{\textbackslash{}bw}}}] borrowed word (source of a loan word)
\item[{\hyperref[key:es]{\textcolor{cflink}{\textbackslash{}es}}}] etymology source
\item[{\hyperref[key:so]{\textcolor{cflink}{\textbackslash{}so}}}] source of data/information
\item[{\hyperref[key:st]{\textcolor{cflink}{\textbackslash{}st}}}] status of the lexical entry
\end{description}


The easiest way to start a Range Set for a particular field, is to Right-Click on the field marker (option-click on the Macintosh) and select the ``Range Set'' tab from the Marker Properties dialog box that pops up. You can enter the labels directly (separated by a <space>) or use the  Build Range Set feature. This will build a Range Set on the current field marker using the contents of all occurrences of this field in all of the currently opened databases of this type. This is a great way to start a new Range Set on an existing database or to compile a new list at any time of what is currently in use (to see how consistent you are being, e.g. am I using : or \_ to concatenate complex labels, did I add a new label when I already had one that is spelled almost the same, etc.).


\begin{quote}\small\itshape Do not build a Range Set on a field which does not have a closed set of possible contents.\end{quote}

\section{Changing Range Set Labels (for the \textbackslash{}ps, \textbackslash{}lf, and \textbackslash{}pdl fields)}

MDF standard printing provides a way (for certain fields) to list English and National Audience substitutes for each label in the Range Sets. (The purpose of the English substitute allows users to consistently change the formatted output labels from the labels they use in their database, e.g. one-letter codes in the paradigm label fields could be changed to more readable abbreviations).


\begin{quote}\small\itshape This ability to make Range Set label substitutions is limited to the \textbackslash{}ps, \textbackslash{}lf, and \textbackslash{}pdl fields. (If more fields need to be covered, contact Toolbox@sil.org.\end{quote}

{\small\color{exlabel}(1)}\par
\noindent
\begin{Verbatim}[fontsize=\footnotesize,breaklines,breakanywhere]
Note: with this feature, you no longer need to use the \pn field to output national
parts of speech. (But if you ever want to sort or browse on the national part of speech, 
you will still need to use \pn.)

\end{Verbatim}

These Range Set label substitutions are handled by the same CC tables that change the default MDF field labels. For more information on these CC tables, see:


{\color{cflink} \hyperref[key:Printed_Field_Labels]{Printed\_Field\_Labels}}


The areas in these CC tables that need ``tweaking'' are the special ``group'' sections, e.g. group(gPartOfSpeech), group(gParadigmLabel), etc. Be sure to follow the instructions given in these groups carefully. It is best to include a CC rule for every label in the field's Range Set. To make sure you don't miss any:


{\small\color{exlabel}(2)}\par
\noindent
\begin{Verbatim}[fontsize=\footnotesize,breaklines,breakanywhere]
   1. In Toolbox, Right-click on the field marker, to go to its 
      properties sheet.
\end{Verbatim}

{\small\color{exlabel}(3)}\par
\noindent
\begin{Verbatim}[fontsize=\footnotesize,breaklines,breakanywhere]
   2. Click on the "Range Sets..." button
\end{Verbatim}

{\small\color{exlabel}(4)}\par
\noindent
\begin{Verbatim}[fontsize=\footnotesize,breaklines,breakanywhere]
   3. Click on "Build Range Set"
\end{Verbatim}

{\small\color{exlabel}(5)}\par
\noindent
\begin{Verbatim}[fontsize=\footnotesize,breaklines,breakanywhere]
   4. Use the mouse to select all of the Range Set labels in the 
      edit window on the left.
\end{Verbatim}

{\small\color{exlabel}(6)}\par
\noindent
\begin{Verbatim}[fontsize=\footnotesize,breaklines,breakanywhere]
   5. Copy them to the Clipboard (<Ctrl-C>)
\end{Verbatim}

{\small\color{exlabel}(7)}\par
\noindent
\begin{Verbatim}[fontsize=\footnotesize,breaklines,breakanywhere]
   6. Open Notepad (or any text editor) and Paste in the labels
      (that is <Ctrl-V>).
\end{Verbatim}

{\small\color{exlabel}(8)}\par
\noindent
\begin{Verbatim}[fontsize=\footnotesize,breaklines,breakanywhere]
   7. Now you can print the labels or save them to a file.
\end{Verbatim}

{\small\color{exlabel}(9)}\par
\noindent
\begin{Verbatim}[fontsize=\footnotesize,breaklines,breakanywhere]
   8. Use Notepad (or whatever) to edit the appropriate CC table
      and add a CC rule for each label in your Range Set.

\end{Verbatim}

\begin{quote}\small\itshape For more information on Range Sets, see the Shoebox helps file.\end{quote}

To continue with the general discussion, jump to:


{\color{cflink} \hyperref[key:Character_Style_Codes]{Character\_Style\_Codes}}


\chapter{Recommended Field Order}\label{key:Order_of_Fields}

The following are all of the fields in their recommended order for a lexical entry. The order shown here assumes you are using the standard hierarchy (the alternate hierarchy only varies in the hierarchical order of the \textbackslash{}sn, \textbackslash{}se, and \textbackslash{}ps fields). For more information on the alternate hierarchy, see:


{\color{cflink} \hyperref[key:Alternate_Hierarchy]{Alternate\_Hierarchy}}


The field order shown here is nearly identical to the order in which MDF outputs the fields when formatting your lexicon. This needs to be stressed that when MDF formats your lexicon, the fields are rearranged to (nearly) the following order. This is to provide a consistent field sequence to each and every lexical entry (a much desired goal of any published work).


The main differences between the recommended order given here and the order in which MDF outputs the fields are as follows:


1)  the \textbackslash{}lx and \textbackslash{}lc fields -- in the database, the \textbackslash{}lx field must come first,


but for printing, the citation (\textbackslash{}lc) field is the first field out (if it exists).


2)  the gloss and definition fields -- the gloss field prints only if there


is no definition field.


3)  the reverse fields and the word gloss fields -- these generally do not


print; but if you request them, they print grouped together after the


definitions.


It is recommended that you use the database template feature of Toolbox to help you maintain a consistent ordering of your fields.


\begin{quote}\small\itshape Deviating from  this order (especially for the \textbackslash{}se, \textbackslash{}ps,  \textbackslash{}pn, and \textbackslash{}sn fields) could foul MDF formatting, since these fields also serve to mark section boundaries. For more on this, see:\end{quote}

{\color{cflink} \hyperref[key:Sections_in_a_Lexical_Entry]{Sections\_in\_a\_Lexical\_Entry}}


In the following list, the following codes are used:


{\small\color{exlabel}(10)}\par
\noindent
\begin{Verbatim}[fontsize=\footnotesize,breaklines,breakanywhere]
      'B' marks 'Basic' fields (17 total)
\end{Verbatim}

{\small\color{exlabel}(11)}\par
\noindent
\begin{Verbatim}[fontsize=\footnotesize,breaklines,breakanywhere]
      'R' marks 'Reserved' fields (4 total)
\end{Verbatim}

{\small\color{exlabel}(12)}\par
\noindent
\begin{Verbatim}[fontsize=\footnotesize,breaklines,breakanywhere]
      (all other fields are 'Optional')

\end{Verbatim}

{\small\color{exlabel}(13)}\par
\noindent
\begin{Verbatim}[fontsize=\footnotesize,breaklines,breakanywhere]
---------------------------------------------------------
\end{Verbatim}

{\small\color{exlabel}(14)}\par
\noindent
\begin{Verbatim}[fontsize=\footnotesize,breaklines,breakanywhere]
  \lx  (lexical entry)         ONLY ONE PER LEXICAL ENTRY
\end{Verbatim}

{\small\color{exlabel}(15)}\par
\noindent
\begin{Verbatim}[fontsize=\footnotesize,breaklines,breakanywhere]
R \hm  (homonym number)
\end{Verbatim}

{\small\color{exlabel}(16)}\par
\noindent
\begin{Verbatim}[fontsize=\footnotesize,breaklines,breakanywhere]
R \lc  (lexical citation)
\end{Verbatim}

{\small\color{exlabel}(17)}\par
\noindent
\begin{Verbatim}[fontsize=\footnotesize,breaklines,breakanywhere]
  \ph  (phonetic/phonemic form)

\end{Verbatim}

{\small\color{exlabel}(18)}\par
\noindent
\begin{Verbatim}[fontsize=\footnotesize,breaklines,breakanywhere]
R \se  (subentry)
\end{Verbatim}

{\small\color{exlabel}(19)}\par
\noindent
\begin{Verbatim}[fontsize=\footnotesize,breaklines,breakanywhere]
B \ps  (part of speech)
\end{Verbatim}

{\small\color{exlabel}(20)}\par
\noindent
\begin{Verbatim}[fontsize=\footnotesize,breaklines,breakanywhere]
B \pn  (part of speech-national)
\end{Verbatim}

{\small\color{exlabel}(21)}\par
\noindent
\begin{Verbatim}[fontsize=\footnotesize,breaklines,breakanywhere]
R \sn  (sense number)

\end{Verbatim}

{\small\color{exlabel}(22)}\par
\noindent
\begin{Verbatim}[fontsize=\footnotesize,breaklines,breakanywhere]
  \gv  (gloss-vernacular)
\end{Verbatim}

{\small\color{exlabel}(23)}\par
\noindent
\begin{Verbatim}[fontsize=\footnotesize,breaklines,breakanywhere]
  \dv  (definition-vernacular)
\end{Verbatim}

{\small\color{exlabel}(24)}\par
\noindent
\begin{Verbatim}[fontsize=\footnotesize,breaklines,breakanywhere]
B \ge  (gloss-English)
\end{Verbatim}

{\small\color{exlabel}(25)}\par
\noindent
\begin{Verbatim}[fontsize=\footnotesize,breaklines,breakanywhere]
B \re  (reversal form-English)  [default set not to print]
\end{Verbatim}

{\small\color{exlabel}(26)}\par
\noindent
\begin{Verbatim}[fontsize=\footnotesize,breaklines,breakanywhere]
  \we  (word gloss-English)     [default set not to print]
\end{Verbatim}

{\small\color{exlabel}(27)}\par
\noindent
\begin{Verbatim}[fontsize=\footnotesize,breaklines,breakanywhere]
B \de  (definition-English)
\end{Verbatim}

{\small\color{exlabel}(28)}\par
\noindent
\begin{Verbatim}[fontsize=\footnotesize,breaklines,breakanywhere]
B \gn  (gloss-national)
\end{Verbatim}

{\small\color{exlabel}(29)}\par
\noindent
\begin{Verbatim}[fontsize=\footnotesize,breaklines,breakanywhere]
B \rn  (reversal form-national) [default set not to print]
\end{Verbatim}

{\small\color{exlabel}(30)}\par
\noindent
\begin{Verbatim}[fontsize=\footnotesize,breaklines,breakanywhere]
  \wn  (word gloss-national)    [default set not to print]
\end{Verbatim}

{\small\color{exlabel}(31)}\par
\noindent
\begin{Verbatim}[fontsize=\footnotesize,breaklines,breakanywhere]
B \dn  (definition-national)
\end{Verbatim}

{\small\color{exlabel}(32)}\par
\noindent
\begin{Verbatim}[fontsize=\footnotesize,breaklines,breakanywhere]
  \gr  (gloss-regional)
\end{Verbatim}

{\small\color{exlabel}(33)}\par
\noindent
\begin{Verbatim}[fontsize=\footnotesize,breaklines,breakanywhere]
  \rr  (reversal form-regional) [default set not to print]
\end{Verbatim}

{\small\color{exlabel}(34)}\par
\noindent
\begin{Verbatim}[fontsize=\footnotesize,breaklines,breakanywhere]
  \wr  (word gloss-regional)    [default set not to print]
\end{Verbatim}

{\small\color{exlabel}(35)}\par
\noindent
\begin{Verbatim}[fontsize=\footnotesize,breaklines,breakanywhere]
  \dr  (definition-regional) 

\end{Verbatim}

{\small\color{exlabel}(36)}\par
\noindent
\begin{Verbatim}[fontsize=\footnotesize,breaklines,breakanywhere]
  \lt  (literal meaning)
\end{Verbatim}

{\small\color{exlabel}(37)}\par
\noindent
\begin{Verbatim}[fontsize=\footnotesize,breaklines,breakanywhere]
  \sc  (scientific name) 

\end{Verbatim}

{\small\color{exlabel}(38)}\par
\noindent
\begin{Verbatim}[fontsize=\footnotesize,breaklines,breakanywhere]
B \rf  (reference to notebooks, etc.)
\end{Verbatim}

{\small\color{exlabel}(39)}\par
\noindent
\begin{Verbatim}[fontsize=\footnotesize,breaklines,breakanywhere]
B \xv  (example-vernacular)
\end{Verbatim}

{\small\color{exlabel}(40)}\par
\noindent
\begin{Verbatim}[fontsize=\footnotesize,breaklines,breakanywhere]
B \xe  (example-English)
\end{Verbatim}

{\small\color{exlabel}(41)}\par
\noindent
\begin{Verbatim}[fontsize=\footnotesize,breaklines,breakanywhere]
B \xn  (example-national)
\end{Verbatim}

{\small\color{exlabel}(42)}\par
\noindent
\begin{Verbatim}[fontsize=\footnotesize,breaklines,breakanywhere]
  \xr  (example-regional)

\end{Verbatim}

{\small\color{exlabel}(43)}\par
\noindent
\begin{Verbatim}[fontsize=\footnotesize,breaklines,breakanywhere]
  \uv  (usage-vernacular)
\end{Verbatim}

{\small\color{exlabel}(44)}\par
\noindent
\begin{Verbatim}[fontsize=\footnotesize,breaklines,breakanywhere]
  \ue  (usage-English)
\end{Verbatim}

{\small\color{exlabel}(45)}\par
\noindent
\begin{Verbatim}[fontsize=\footnotesize,breaklines,breakanywhere]
  \un  (usage-national)
\end{Verbatim}

{\small\color{exlabel}(46)}\par
\noindent
\begin{Verbatim}[fontsize=\footnotesize,breaklines,breakanywhere]
  \ur  (usage-regional)
\end{Verbatim}

{\small\color{exlabel}(47)}\par
\noindent
\begin{Verbatim}[fontsize=\footnotesize,breaklines,breakanywhere]
  \ev  (encyclopedic info.-vernacular)
\end{Verbatim}

{\small\color{exlabel}(48)}\par
\noindent
\begin{Verbatim}[fontsize=\footnotesize,breaklines,breakanywhere]
  \ee  (encyclopedic info.-English)
\end{Verbatim}

{\small\color{exlabel}(49)}\par
\noindent
\begin{Verbatim}[fontsize=\footnotesize,breaklines,breakanywhere]
  \en  (encyclopedic info.-national)
\end{Verbatim}

{\small\color{exlabel}(50)}\par
\noindent
\begin{Verbatim}[fontsize=\footnotesize,breaklines,breakanywhere]
  \er  (encyclopedic info.-regional)
\end{Verbatim}

{\small\color{exlabel}(51)}\par
\noindent
\begin{Verbatim}[fontsize=\footnotesize,breaklines,breakanywhere]
  \ov  (only [restrictions]-vernacular)
\end{Verbatim}

{\small\color{exlabel}(52)}\par
\noindent
\begin{Verbatim}[fontsize=\footnotesize,breaklines,breakanywhere]
  \oe  (only [restrictions]-English)
\end{Verbatim}

{\small\color{exlabel}(53)}\par
\noindent
\begin{Verbatim}[fontsize=\footnotesize,breaklines,breakanywhere]
  \on  (only [restrictions]-national)
\end{Verbatim}

{\small\color{exlabel}(54)}\par
\noindent
\begin{Verbatim}[fontsize=\footnotesize,breaklines,breakanywhere]
  \or  (only [restrictions]-regional)

\end{Verbatim}

{\small\color{exlabel}(55)}\par
\noindent
\begin{Verbatim}[fontsize=\footnotesize,breaklines,breakanywhere]
  \lf  (lexical function label)
\end{Verbatim}

{\small\color{exlabel}(56)}\par
\noindent
\begin{Verbatim}[fontsize=\footnotesize,breaklines,breakanywhere]
  \lv  (vernacular lexeme referenced by the lexical function)
\end{Verbatim}

{\small\color{exlabel}(57)}\par
\noindent
\begin{Verbatim}[fontsize=\footnotesize,breaklines,breakanywhere]
  \le  (lexical function-English gloss of referenced lexeme)
\end{Verbatim}

{\small\color{exlabel}(58)}\par
\noindent
\begin{Verbatim}[fontsize=\footnotesize,breaklines,breakanywhere]
  \ln  (lexical function-national gloss of referenced lexeme)
\end{Verbatim}

{\small\color{exlabel}(59)}\par
\noindent
\begin{Verbatim}[fontsize=\footnotesize,breaklines,breakanywhere]
  \lr  (lexical function-regional gloss of referenced lexeme)
\end{Verbatim}

{\small\color{exlabel}(60)}\par
\noindent
\begin{Verbatim}[fontsize=\footnotesize,breaklines,breakanywhere]
  \sy  (synonyms)
\end{Verbatim}

{\small\color{exlabel}(61)}\par
\noindent
\begin{Verbatim}[fontsize=\footnotesize,breaklines,breakanywhere]
  \an  (antonyms)
\end{Verbatim}

{\small\color{exlabel}(62)}\par
\noindent
\begin{Verbatim}[fontsize=\footnotesize,breaklines,breakanywhere]
  \mr  (morphemic representation)
\end{Verbatim}

{\small\color{exlabel}(63)}\par
\noindent
\begin{Verbatim}[fontsize=\footnotesize,breaklines,breakanywhere]
B \cf  (confer/cross-reference)
\end{Verbatim}

{\small\color{exlabel}(64)}\par
\noindent
\begin{Verbatim}[fontsize=\footnotesize,breaklines,breakanywhere]
B \ce  (cross-reference-English gloss)
\end{Verbatim}

{\small\color{exlabel}(65)}\par
\noindent
\begin{Verbatim}[fontsize=\footnotesize,breaklines,breakanywhere]
B \cn  (cross-reference-national gloss)
\end{Verbatim}

{\small\color{exlabel}(66)}\par
\noindent
\begin{Verbatim}[fontsize=\footnotesize,breaklines,breakanywhere]
  \cr  (cross-reference-regional gloss)
\end{Verbatim}

{\small\color{exlabel}(67)}\par
\noindent
\begin{Verbatim}[fontsize=\footnotesize,breaklines,breakanywhere]
  \mn  (main entry form)
\end{Verbatim}

{\small\color{exlabel}(68)}\par
\noindent
\begin{Verbatim}[fontsize=\footnotesize,breaklines,breakanywhere]
  \va  (variant forms, e.g. dialect, etc.)
\end{Verbatim}

{\small\color{exlabel}(69)}\par
\noindent
\begin{Verbatim}[fontsize=\footnotesize,breaklines,breakanywhere]
  \ve  (variant forms comment-English)
\end{Verbatim}

{\small\color{exlabel}(70)}\par
\noindent
\begin{Verbatim}[fontsize=\footnotesize,breaklines,breakanywhere]
  \vn  (variant forms comment-national)
\end{Verbatim}

{\small\color{exlabel}(71)}\par
\noindent
\begin{Verbatim}[fontsize=\footnotesize,breaklines,breakanywhere]
  \vr  (variant forms comment-regional)
\end{Verbatim}

{\small\color{exlabel}(72)}\par
\noindent
\begin{Verbatim}[fontsize=\footnotesize,breaklines,breakanywhere]
  \bw  (borrowed word)
\end{Verbatim}

{\small\color{exlabel}(73)}\par
\noindent
\begin{Verbatim}[fontsize=\footnotesize,breaklines,breakanywhere]
  \et  (etymology)
\end{Verbatim}

{\small\color{exlabel}(74)}\par
\noindent
\begin{Verbatim}[fontsize=\footnotesize,breaklines,breakanywhere]
  \eg  (etymology-gloss)
\end{Verbatim}

{\small\color{exlabel}(75)}\par
\noindent
\begin{Verbatim}[fontsize=\footnotesize,breaklines,breakanywhere]
  \es  (etymology-source)   [default set not to print]
\end{Verbatim}

{\small\color{exlabel}(76)}\par
\noindent
\begin{Verbatim}[fontsize=\footnotesize,breaklines,breakanywhere]
  \ec  (etymology-comment)  [default set not to print]

\end{Verbatim}

{\small\color{exlabel}(77)}\par
\noindent
\begin{Verbatim}[fontsize=\footnotesize,breaklines,breakanywhere]
  \pd  (paradigm set)
\end{Verbatim}

{\small\color{exlabel}(78)}\par
\noindent
\begin{Verbatim}[fontsize=\footnotesize,breaklines,breakanywhere]
  \pdl (paradigm label)
\end{Verbatim}

{\small\color{exlabel}(79)}\par
\noindent
\begin{Verbatim}[fontsize=\footnotesize,breaklines,breakanywhere]
  \pdv (paradigm vernacular form)
\end{Verbatim}

{\small\color{exlabel}(80)}\par
\noindent
\begin{Verbatim}[fontsize=\footnotesize,breaklines,breakanywhere]
  \pde (paradigm form-English gloss)
\end{Verbatim}

{\small\color{exlabel}(81)}\par
\noindent
\begin{Verbatim}[fontsize=\footnotesize,breaklines,breakanywhere]
  \pdn (paradigm form-national gloss)
\end{Verbatim}

{\small\color{exlabel}(82)}\par
\noindent
\begin{Verbatim}[fontsize=\footnotesize,breaklines,breakanywhere]
  \pdr (paradigm form-regional gloss)
\end{Verbatim}

{\small\color{exlabel}(83)}\par
\noindent
\begin{Verbatim}[fontsize=\footnotesize,breaklines,breakanywhere]
  \sg  (singular noun form)
\end{Verbatim}

{\small\color{exlabel}(84)}\par
\noindent
\begin{Verbatim}[fontsize=\footnotesize,breaklines,breakanywhere]
  \pl  (plural noun form)
\end{Verbatim}

{\small\color{exlabel}(85)}\par
\noindent
\begin{Verbatim}[fontsize=\footnotesize,breaklines,breakanywhere]
  \rd  (reduplication forms)
\end{Verbatim}

{\small\color{exlabel}(86)}\par
\noindent
\begin{Verbatim}[fontsize=\footnotesize,breaklines,breakanywhere]
  \1s  (1st person singular verb form)
\end{Verbatim}

{\small\color{exlabel}(87)}\par
\noindent
\begin{Verbatim}[fontsize=\footnotesize,breaklines,breakanywhere]
  \2s  (2nd person singular verb form)
\end{Verbatim}

{\small\color{exlabel}(88)}\par
\noindent
\begin{Verbatim}[fontsize=\footnotesize,breaklines,breakanywhere]
  \3s  (3rd person singular verb form)
\end{Verbatim}

{\small\color{exlabel}(89)}\par
\noindent
\begin{Verbatim}[fontsize=\footnotesize,breaklines,breakanywhere]
  \4s  (singular non-human/non-animate verb form)
\end{Verbatim}

{\small\color{exlabel}(90)}\par
\noindent
\begin{Verbatim}[fontsize=\footnotesize,breaklines,breakanywhere]
  \1d  (1st person dual verb form)
\end{Verbatim}

{\small\color{exlabel}(91)}\par
\noindent
\begin{Verbatim}[fontsize=\footnotesize,breaklines,breakanywhere]
  \2d  (2nd person dual verb form)
\end{Verbatim}

{\small\color{exlabel}(92)}\par
\noindent
\begin{Verbatim}[fontsize=\footnotesize,breaklines,breakanywhere]
  \3d  (3rd person dual verb form)
\end{Verbatim}

{\small\color{exlabel}(93)}\par
\noindent
\begin{Verbatim}[fontsize=\footnotesize,breaklines,breakanywhere]
  \4d  (dual non-human/non-animate verb form)
\end{Verbatim}

{\small\color{exlabel}(94)}\par
\noindent
\begin{Verbatim}[fontsize=\footnotesize,breaklines,breakanywhere]
  \1p  (1st person plural-generic verb form)
\end{Verbatim}

{\small\color{exlabel}(95)}\par
\noindent
\begin{Verbatim}[fontsize=\footnotesize,breaklines,breakanywhere]
  \1e  (1st person plural-exclusive verb form)
\end{Verbatim}

{\small\color{exlabel}(96)}\par
\noindent
\begin{Verbatim}[fontsize=\footnotesize,breaklines,breakanywhere]
  \1i  (1st person plural-inclusive verb form)
\end{Verbatim}

{\small\color{exlabel}(97)}\par
\noindent
\begin{Verbatim}[fontsize=\footnotesize,breaklines,breakanywhere]
  \2p  (2nd person plural verb form)
\end{Verbatim}

{\small\color{exlabel}(98)}\par
\noindent
\begin{Verbatim}[fontsize=\footnotesize,breaklines,breakanywhere]
  \3p  (3rd person plural verb form)
\end{Verbatim}

{\small\color{exlabel}(99)}\par
\noindent
\begin{Verbatim}[fontsize=\footnotesize,breaklines,breakanywhere]
  \4p  (plural non-human/non-animate verb form)
 
\end{Verbatim}

{\small\color{exlabel}(100)}\par
\noindent
\begin{Verbatim}[fontsize=\footnotesize,breaklines,breakanywhere]
  \tb  (table)

\end{Verbatim}

{\small\color{exlabel}(101)}\par
\noindent
\begin{Verbatim}[fontsize=\footnotesize,breaklines,breakanywhere]
  \sd  (semantic domain)    [default set not to print]
\end{Verbatim}

{\small\color{exlabel}(102)}\par
\noindent
\begin{Verbatim}[fontsize=\footnotesize,breaklines,breakanywhere]
  \is  (index of semantics) [default set not to print]
\end{Verbatim}

{\small\color{exlabel}(103)}\par
\noindent
\begin{Verbatim}[fontsize=\footnotesize,breaklines,breakanywhere]
  \th  (thesaurus)          [default set not to print]

\end{Verbatim}

{\small\color{exlabel}(104)}\par
\noindent
\begin{Verbatim}[fontsize=\footnotesize,breaklines,breakanywhere]
  \bb  (bibliographic reference)
\end{Verbatim}

{\small\color{exlabel}(105)}\par
\noindent
\begin{Verbatim}[fontsize=\footnotesize,breaklines,breakanywhere]
  \pc  (picture reference)

\end{Verbatim}

{\small\color{exlabel}(106)}\par
\noindent
\begin{Verbatim}[fontsize=\footnotesize,breaklines,breakanywhere]
B \nt  (notes, etc.)
\end{Verbatim}

{\small\color{exlabel}(107)}\par
\noindent
\begin{Verbatim}[fontsize=\footnotesize,breaklines,breakanywhere]
  \np  (notes on phonology)
\end{Verbatim}

{\small\color{exlabel}(108)}\par
\noindent
\begin{Verbatim}[fontsize=\footnotesize,breaklines,breakanywhere]
  \ng  (notes on grammar)
\end{Verbatim}

{\small\color{exlabel}(109)}\par
\noindent
\begin{Verbatim}[fontsize=\footnotesize,breaklines,breakanywhere]
  \nd  (notes on discourse)
\end{Verbatim}

{\small\color{exlabel}(110)}\par
\noindent
\begin{Verbatim}[fontsize=\footnotesize,breaklines,breakanywhere]
  \na  (notes on anthropology)
\end{Verbatim}

{\small\color{exlabel}(111)}\par
\noindent
\begin{Verbatim}[fontsize=\footnotesize,breaklines,breakanywhere]
  \ns  (notes on sociolinguistics)
\end{Verbatim}

{\small\color{exlabel}(112)}\par
\noindent
\begin{Verbatim}[fontsize=\footnotesize,breaklines,breakanywhere]
  \nq  (questions)

\end{Verbatim}

{\small\color{exlabel}(113)}\par
\noindent
\begin{Verbatim}[fontsize=\footnotesize,breaklines,breakanywhere]
  \so  (source of data)     [default set not to print]
\end{Verbatim}

{\small\color{exlabel}(114)}\par
\noindent
\begin{Verbatim}[fontsize=\footnotesize,breaklines,breakanywhere]
  \st  (status)             [default set not to print]
\end{Verbatim}

{\small\color{exlabel}(115)}\par
\noindent
\begin{Verbatim}[fontsize=\footnotesize,breaklines,breakanywhere]
B \dt  (datestamp)          [default set not to print]
\end{Verbatim}

{\small\color{exlabel}(116)}\par
\noindent
\begin{Verbatim}[fontsize=\footnotesize,breaklines,breakanywhere]
--------------------end of example record--------------------

\end{Verbatim}

\begin{quote}\small\itshape (To see the description of any particular field, simply right-click on its marker.)\end{quote}

To continue this discussion, jump to:


{\color{cflink} \hyperref[key:Summary_of_Fields]{Summary\_of\_Fields}}


\chapter{Character Style Codes}\label{key:Character_Style_Codes}

Character style codes can be used in any field (usually the information or discussion type fields) to uniquely mark the specific language of the text that you are discussing. For example, `fv:' marks the following word as vernacular text:


{\small\color{exlabel}(117)}\par
\noindent
\begin{Verbatim}[fontsize=\footnotesize,breaklines,breakanywhere]
  \ue This usage is a bit blunt.  Most people don't say 
      this directly, rather they refer to where they are 
      going, e.g. fv:keta to refer to 'outside the village',
      but fv:-dai is used in familiar settings.

\end{Verbatim}

Character style codes are applied only to the following word. A space or any punctuation (except the `-') terminates the style. The style codes must be in lower case, and there must not be any space between the colon and the following word, e.g.


{\small\color{exlabel}(118)}\par
\noindent
\begin{Verbatim}[fontsize=\footnotesize,breaklines,breakanywhere]
       fv:sabun       is okay
\end{Verbatim}

{\small\color{exlabel}(119)}\par
\noindent
\begin{Verbatim}[fontsize=\footnotesize,breaklines,breakanywhere]
       fv: sabun      won't work (the fv: is treated as normal text)

\end{Verbatim}

Use the underline ( \_ ) character to continue the style and link together words in a phrase, e.g.


{\small\color{exlabel}(120)}\par
\noindent
\begin{Verbatim}[fontsize=\footnotesize,breaklines,breakanywhere]
  \ue In the phrase fv:Mala_sai_desy_de? there is a sense of 
      anger or rudeness.
\end{Verbatim}

MDF will format this whole phrase with the vernacular character style (the underlined characters are converted to spaces). The question mark will not be included in the vernacular style.


Character codes should be placed with the word, inside parentheses, quotes, brackets, etc., otherwise the opening brackets will receive the character style along with the word.


\section{Basic Language and Character Codes}

{\small\color{exlabel}(121)}\par
\noindent
\begin{Verbatim}[fontsize=\footnotesize,breaklines,breakanywhere]
     fv: (font--vernacular)
\end{Verbatim}

{\small\color{exlabel}(122)}\par
\noindent
\begin{Verbatim}[fontsize=\footnotesize,breaklines,breakanywhere]
     fe: (font--English)
\end{Verbatim}

{\small\color{exlabel}(123)}\par
\noindent
\begin{Verbatim}[fontsize=\footnotesize,breaklines,breakanywhere]
     fn: (font--national language)
\end{Verbatim}

{\small\color{exlabel}(124)}\par
\noindent
\begin{Verbatim}[fontsize=\footnotesize,breaklines,breakanywhere]
     fr: (font--regional language)

\end{Verbatim}

{\small\color{exlabel}(125)}\par
\noindent
\begin{Verbatim}[fontsize=\footnotesize,breaklines,breakanywhere]
     fs: (font--standard)
\end{Verbatim}

{\small\color{exlabel}(126)}\par
\noindent
\begin{Verbatim}[fontsize=\footnotesize,breaklines,breakanywhere]
     fb: (font--bold)
\end{Verbatim}

{\small\color{exlabel}(127)}\par
\noindent
\begin{Verbatim}[fontsize=\footnotesize,breaklines,breakanywhere]
     fi: (font--italic)

\end{Verbatim}

{\small\color{exlabel}(128)}\par
\noindent
\begin{Verbatim}[fontsize=\footnotesize,breaklines,breakanywhere]
     uc: (underline character)
\end{Verbatim}

{\small\color{exlabel}(129)}\par
\noindent
\begin{Verbatim}[fontsize=\footnotesize,breaklines,breakanywhere]
     ub: (underline--bold)
\end{Verbatim}

{\small\color{exlabel}(130)}\par
\noindent
\begin{Verbatim}[fontsize=\footnotesize,breaklines,breakanywhere]
     ui: (underline--italic)

\end{Verbatim}

\section{Other Character Codes}

{\small\color{exlabel}(131)}\par
\noindent
\begin{Verbatim}[fontsize=\footnotesize,breaklines,breakanywhere]
     sc: (underline a scientific name--not 
               required in the \sc field)

\end{Verbatim}

The uc: code is able to detect which type of field it is used in. If the field is a vernacular field uc: will underline with bold characters (following the vernacular character style); if the field uses the national language, uc: will underline italic characters; and if the field uses English, uc: will underline normal characters (to follow the character style for English). If you want to specifically control the underlining character style, use ui: and ub:.


If you need to tweak your formatted dictionary or reversed finderlists and apply a character style where you had forgotten to place a code, simply use the same character style that MS-Word uses to format other such text (e.g. the vernacular font in MS-Word uses the ``f\_vernacular'' style). Just be sure to also add the missing character style code in your lexical database too, otherwise you'll have to ``tweak'' the output every time.


\section{Explicit Font Codes}

To explicitly control the character style or font of any phrase, word, or part of word, use the underlying bar codes used by MDF. For example, to underline only a part of a word, use the |u bar code:


{\small\color{exlabel}(132)}\par
\noindent
\begin{Verbatim}[fontsize=\footnotesize,breaklines,breakanywhere]
  \xv Aulopo|u{a} au lae weidu.

\end{Verbatim}

The |u\{\} underlines the bracketed character (`a') only. Note that the curly braces \{\} can contain any string of letters, numbers, spaces, or punctuation.


The complete list of explicit font codes follows:


{\small\color{exlabel}(133)}\par
\noindent
\begin{Verbatim}[fontsize=\footnotesize,breaklines,breakanywhere]
     |fv{ }      Vernacular font
\end{Verbatim}

{\small\color{exlabel}(134)}\par
\noindent
\begin{Verbatim}[fontsize=\footnotesize,breaklines,breakanywhere]
     |fe{ }      English font
\end{Verbatim}

{\small\color{exlabel}(135)}\par
\noindent
\begin{Verbatim}[fontsize=\footnotesize,breaklines,breakanywhere]
     |fn{ }      National font
\end{Verbatim}

{\small\color{exlabel}(136)}\par
\noindent
\begin{Verbatim}[fontsize=\footnotesize,breaklines,breakanywhere]
     |fr{ }      Regional font
\end{Verbatim}

{\small\color{exlabel}(137)}\par
\noindent
\begin{Verbatim}[fontsize=\footnotesize,breaklines,breakanywhere]
     |fs{ }      Standard font (for general punctuation)
\end{Verbatim}

{\small\color{exlabel}(138)}\par
\noindent
\begin{Verbatim}[fontsize=\footnotesize,breaklines,breakanywhere]
     |fl{ }      The font for field labels

\end{Verbatim}

{\small\color{exlabel}(139)}\par
\noindent
\begin{Verbatim}[fontsize=\footnotesize,breaklines,breakanywhere]
     |fb{ }      Bold font 
\end{Verbatim}

{\small\color{exlabel}(140)}\par
\noindent
\begin{Verbatim}[fontsize=\footnotesize,breaklines,breakanywhere]
     |fi{ }      Italic font

\end{Verbatim}

Still supported, but not encouraged, are the older (``Manuscripter'') style codes.


Underlining is more complex and so provides the following:


{\small\color{exlabel}(141)}\par
\noindent
\begin{Verbatim}[fontsize=\footnotesize,breaklines,breakanywhere]
     |u{ }       General underlined character font (as shown above) 
\end{Verbatim}

{\small\color{exlabel}(142)}\par
\noindent
\begin{Verbatim}[fontsize=\footnotesize,breaklines,breakanywhere]
     |uc{ }      Specific regular underlined character font 
\end{Verbatim}

{\small\color{exlabel}(143)}\par
\noindent
\begin{Verbatim}[fontsize=\footnotesize,breaklines,breakanywhere]
     |ub{ }      Specific underlined, bold font 
\end{Verbatim}

{\small\color{exlabel}(144)}\par
\noindent
\begin{Verbatim}[fontsize=\footnotesize,breaklines,breakanywhere]
     |ui{ }      Specific underlined, italic font

\end{Verbatim}

To format a homonym number in a discussion field you can use the explicit code:


{\small\color{exlabel}(145)}\par
\noindent
\begin{Verbatim}[fontsize=\footnotesize,breaklines,breakanywhere]
     |hm{ }      Homonym number font  

\end{Verbatim}

For example:


{\small\color{exlabel}(146)}\par
\noindent
\begin{Verbatim}[fontsize=\footnotesize,breaklines,breakanywhere]
  \ec This lexeme is now a homonym with fv:asw|hm{1} 'dog' 
      because proto *l ~ s in Selaru resulting in a loss 
      of contrast.

\end{Verbatim}

\begin{quote}\small\itshape If you choose to use these explicit style codes, you must not forget the closing brace ` \} `. If you do, the rest of the field will be considered a part of the specified font!\end{quote}

\begin{quote}\small\itshape Actually, any field marker can become a character style code and transmit its formatting to the data enclosed in the curly brackets. This is mostly useful for the fields mentioned already (sc, hm) but any field code can be used.\end{quote}

To continue with the general discussion, jump to:


{\color{cflink} \hyperref[key:Punctuation_and_Special_Codes]{Punctuation\_and\_Special\_Codes}}


\chapter{Summary of Fields (not ordered)}\label{key:Summary_of_Fields}

The following are all of the fields MDF recognizes, listed  in their suggested groups:


\begin{quote}\small\itshape ( ` f ` marks free-form fields)\end{quote}

{\small\color{exlabel}(147)}\par
\noindent
\begin{Verbatim}[fontsize=\footnotesize,breaklines,breakanywhere]
RECORD MARKER
 \lx  lexical entry  (only one allowed per record)

\end{Verbatim}

{\small\color{exlabel}(148)}\par
\noindent
\begin{Verbatim}[fontsize=\footnotesize,breaklines,breakanywhere]
BASIC FIELDS                   OPTIONAL FIELDS (Cont.)
  \ps  part of speech            \1i  1st plural incl. form    
  \pn  part of speech (natnl)    \2p  2nd plural form 
  \ge  gloss (English)           \3p  3rd plural form
  \gn  gloss (national)          \4p  pl. non-human form
  \re  reversal (English)        \ph  phonetic form
  \rn  reversal (national)       \cr  cross ref. (regional)
  \de  definition (English)      \mr  morphemic form 
  \dn  definition (national)     \rd  reduplication form 
f \rf  reference to notebooks    \va  variant form
f \xv  example (vern.)           \ve  variant (Engish)
f \xe  example (English)         \vn  variant (national)
f \xn  example (national)        \vr  variant (regional)
  \cf  cross reference           \mn  main entry form
  \ce  cross ref. (English)      \lf  lexical function
  \cn  cross ref. (national)     \lv  lexeme ref'd by lexical fnct
f \nt  notes, etc.               \le  lexical fnct (English)
  \dt  datestamp                 \ln  lexical fnct (national)
                                 \lr  lexical fnct (regional)
RESERVED FIELDS                  \sy  synonyms
  \hm  homonym number            \an  antonyms 
  \lc  lexical citation        f \uv  usage (vernacular)
  \se  subentry                f \ue  usage (English)
  \sn  sense number            f \un  usage (national)
                               f \ur  usage (regional)
OPTIONAL FIELDS                f \ov  only (vernacular)
  \gv  gloss (vernacular)      f \oe  only (English)
  \gr  gloss (regional)        f \on  only (national)
  \rr  reversal (Regn)         f \or  only (regional)
  \we  word-gloss (Engl)       f \ev  encyclo. (vern)
  \wn  word-gloss (Natn)       f \ee  encyclo. (Engl)
  \wr  word-gloss (Regn)       f \en  encyclo. (Natnl)
  \dv  definition (vern)       f \er  encyclo. (Regnl)
  \dr  definition (Regn)         \bw  borrowed word
  \lt  literal meaning           \et  etymology
f \xr  example (Regnl)           \eg  etymology (gloss)
  \pd  paradigm set              \es  etymology (source)
  \pdl paradigm label          f \ec  etymology (comment)
  \pdv paradigm form (vernac)    \sd  semantic domain
  \pde paradigm gloss (Engl)     \is  index of semantics
  \pdn paradigm gloss (Natn)     \th  thesaurus
  \pdr paradigm gloss (Regn)   f \bb  bibliographic ref.
  \sg  sing. noun form           \sc  scientific name
  \pl  plural noun form        f \tb  table/chart
  \1s  1st singular verb form  f \pc  picture
  \2s  2nd singular verb form  f \np  notes on phonology
  \3s  3rd singular verb form  f \ng  notes on grammar
  \4s  sing. non-human form    f \nd  notes on discourse
  \1d  1st dual verb form      f \na  notes on anthro.
  \2d  2nd dual verb form      f \ns  sociolinguistics
  \3d  3rd dual verb form      f \nq  questions
  \4d  dual non-human form       \st  status
  \1p  1st plural form           \so  source of data
  \1e  1st plural excl. form   

\end{Verbatim}

\begin{quote}\small\itshape (To see the description of any particular field, simply right-click on its marker.)\end{quote}

To continue this discussion, jump to:


{\color{cflink} \hyperref[key:Sections_in_a_Lexical_Entry]{Sections\_in\_a\_Lexical\_Entry}}


\chapter{Punctuation}\label{key:Punctuation_and_Special_Codes}

Generally, leave off all punctuation at the end of bookkeeping fields and basic data fields (\textbackslash{}hm, \textbackslash{}sn, \textbackslash{}ge, \textbackslash{}lf, etc.). The only places where you need to include punctuation is in and at the end of free-form (discussion type) fields (\textbackslash{}ue, \textbackslash{}ee, \textbackslash{}nt, etc.). In all other data fields, final punctuation is added by MDF automatically. For a list of free-form fields, see:


{\color{cflink} \hyperref[key:Free-form_Fields]{Free-form\_Fields}}


Toolbox allows you to hard code the following special characters directly into you lexical database (this works for both the RTF and the MDF exporting process):


{\small\color{exlabel}(149)}\par
\noindent
\begin{Verbatim}[fontsize=\footnotesize,breaklines,breakanywhere]
   |{tab}          Inserts a Tab character
\end{Verbatim}

{\small\color{exlabel}(150)}\par
\noindent
\begin{Verbatim}[fontsize=\footnotesize,breaklines,breakanywhere]
   |{endash}       Inserts an En-dash character
\end{Verbatim}

{\small\color{exlabel}(151)}\par
\noindent
\begin{Verbatim}[fontsize=\footnotesize,breaklines,breakanywhere]
   |{emdash}       Inserts an Em-dash character
\end{Verbatim}

{\small\color{exlabel}(152)}\par
\noindent
\begin{Verbatim}[fontsize=\footnotesize,breaklines,breakanywhere]
   |{~}            Inserts a non-breaking space
\end{Verbatim}

{\small\color{exlabel}(153)}\par
\noindent
\begin{Verbatim}[fontsize=\footnotesize,breaklines,breakanywhere]
   |{-}            Inserts an optional hyphen
\end{Verbatim}

{\small\color{exlabel}(154)}\par
\noindent
\begin{Verbatim}[fontsize=\footnotesize,breaklines,breakanywhere]
   |{_}            Inserts a non-breaking hyphen

\end{Verbatim}

For an example which uses the special code for Tab, see:


\begin{description}
\item[{\hyperref[key:tb]{\textcolor{cflink}{\textbackslash{}tb}}}] Tables
\end{description}


For information on special codes for embedding character styles on , see:


{\color{cflink} \hyperref[key:Character_Style_Codes]{Character\_Style\_Codes}}


To continue with the general discussion, jump to:


{\color{cflink} \hyperref[key:Printed_Field_Labels]{Printed\_Field\_Labels}}


\chapter{Sections in a Lexical Entry}\label{key:Sections_in_a_Lexical_Entry}

\section{(Understanding the hierarchical structure of an entry)}

MDF has two built-in hierarchical structures that should be flexible enough to meet most needs. (See ``Alternate\_Hierarchy'' for a discussion of MDF's other hierarchy.) The field codes that mark the boundaries to lexical subsections are \textbackslash{}lx, \textbackslash{}ps (\textbackslash{}pn), \textbackslash{}sn, \textbackslash{}se. Each of these sections or subsections can take a full set of field markers (except \textbackslash{}lc and  \textbackslash{}hm, which should only occur at the top of the record).


\section{Multiple parts of speech}

Multiple parts of speech (\textbackslash{}ps) are used to organize sections within an entry. A lexeme that fills more than one syntactic slot (as a noun, verb, etc.) should not be handled as homonyms. This is because the different syntactic functions (e.g. `shower' (n) `shower' (v) are still clearly related to each other in meaning.


{\small\color{exlabel}(155)}\par
\noindent
\begin{Verbatim}[fontsize=\footnotesize,breaklines,breakanywhere]
 \lx shower
\end{Verbatim}

{\small\color{exlabel}(156)}\par
\noindent
\begin{Verbatim}[fontsize=\footnotesize,breaklines,breakanywhere]
 \ps n
\end{Verbatim}

{\small\color{exlabel}(157)}\par
\noindent
\begin{Verbatim}[fontsize=\footnotesize,breaklines,breakanywhere]
 \de a light rain
\end{Verbatim}

{\small\color{exlabel}(158)}\par
\noindent
\begin{Verbatim}[fontsize=\footnotesize,breaklines,breakanywhere]
 \ps vt
\end{Verbatim}

{\small\color{exlabel}(159)}\par
\noindent
\begin{Verbatim}[fontsize=\footnotesize,breaklines,breakanywhere]
 \de to bestow special things on someone

\end{Verbatim}

MDF's standard printing function starts new \textbackslash{}ps fields within an entry on a new line, preceded by an em dash to show that the lexeme form has not changed but its function has.


{\small\color{exlabel}(160)}\par
\noindent
\begin{Verbatim}[fontsize=\footnotesize,breaklines,breakanywhere]
shower   n. a light rain.
\end{Verbatim}

{\small\color{exlabel}(161)}\par
\noindent
\begin{Verbatim}[fontsize=\footnotesize,breaklines,breakanywhere]
       --  vt. to bestow special things on someone.

\end{Verbatim}

\section{Multiple Senses}

MDF allows sense numbers (\textbackslash{}sn) to be used as another level of hierarchy within an entry. In the standard hierarchy, the sense number is lower than part of speech, and so multiple senses should be grouped under the relevant parts of speech to denote related but distinct meanings within a particular part of speech. (See ``Alternate\_Hierarchy'' for a discussion of MDF's other hierarchy.) Multiple senses in each separate part of speech should start with `1' (as seen in the following, more complete entry for `shower'):


{\small\color{exlabel}(162)}\par
\noindent
\begin{Verbatim}[fontsize=\footnotesize,breaklines,breakanywhere]
 \lx shower
\end{Verbatim}

{\small\color{exlabel}(163)}\par
\noindent
\begin{Verbatim}[fontsize=\footnotesize,breaklines,breakanywhere]
 \ps n
\end{Verbatim}

{\small\color{exlabel}(164)}\par
\noindent
\begin{Verbatim}[fontsize=\footnotesize,breaklines,breakanywhere]
 \sn 1
\end{Verbatim}

{\small\color{exlabel}(165)}\par
\noindent
\begin{Verbatim}[fontsize=\footnotesize,breaklines,breakanywhere]
 \de a light rain
\end{Verbatim}

{\small\color{exlabel}(166)}\par
\noindent
\begin{Verbatim}[fontsize=\footnotesize,breaklines,breakanywhere]
 \sn 2
\end{Verbatim}

{\small\color{exlabel}(167)}\par
\noindent
\begin{Verbatim}[fontsize=\footnotesize,breaklines,breakanywhere]
 \de a man-made device for dispensing water in droplets
     on a person; used for bathing
\end{Verbatim}

{\small\color{exlabel}(168)}\par
\noindent
\begin{Verbatim}[fontsize=\footnotesize,breaklines,breakanywhere]
 \sn 3
\end{Verbatim}

{\small\color{exlabel}(169)}\par
\noindent
\begin{Verbatim}[fontsize=\footnotesize,breaklines,breakanywhere]
 \de an event in which gifts are given to someone; 
     as in baby showers and wedding showers

\end{Verbatim}

{\small\color{exlabel}(170)}\par
\noindent
\begin{Verbatim}[fontsize=\footnotesize,breaklines,breakanywhere]
 \ps v
\end{Verbatim}

{\small\color{exlabel}(171)}\par
\noindent
\begin{Verbatim}[fontsize=\footnotesize,breaklines,breakanywhere]
 \sn 1
\end{Verbatim}

{\small\color{exlabel}(172)}\par
\noindent
\begin{Verbatim}[fontsize=\footnotesize,breaklines,breakanywhere]
 \de raining lightly
\end{Verbatim}

{\small\color{exlabel}(173)}\par
\noindent
\begin{Verbatim}[fontsize=\footnotesize,breaklines,breakanywhere]
 \sn 2
\end{Verbatim}

{\small\color{exlabel}(174)}\par
\noindent
\begin{Verbatim}[fontsize=\footnotesize,breaklines,breakanywhere]
 \de to bathe using a device which causes water to
     dispense in droplets on a persons head; usually
     done standing up
\end{Verbatim}

{\small\color{exlabel}(175)}\par
\noindent
\begin{Verbatim}[fontsize=\footnotesize,breaklines,breakanywhere]
 \sn 3
\end{Verbatim}

{\small\color{exlabel}(176)}\par
\noindent
\begin{Verbatim}[fontsize=\footnotesize,breaklines,breakanywhere]
 \de to bestow special things on someone

\end{Verbatim}

This complex record would print as:


{\small\color{exlabel}(177)}\par
\noindent
\begin{Verbatim}[fontsize=\footnotesize,breaklines,breakanywhere]
shower   n. 1)  a light  rain. 2)  a man-made  device for
         dispensing water  in droplets  on a person; used
         for bathing.  3) an  event in  which  gifts  are
         given to someone; as in baby showers and wedding
         showers.
     --  v. 1)  raining lightly.  2) to  bathe using  a
         device  which   causes  water   to  dispense  in
         droplets  on   a  persons   head;  usually  done
         standing up.  3) to  bestow  special  things  on
         someone.

\end{Verbatim}

Some lexicographers want to make fine distinctions between subsenses. These can be handled in MDF in the  \textbackslash{}sn field with a, b, c, etc. subcategorization.


{\small\color{exlabel}(178)}\par
\noindent
\begin{Verbatim}[fontsize=\footnotesize,breaklines,breakanywhere]
 \lx lexeme
\end{Verbatim}

{\small\color{exlabel}(179)}\par
\noindent
\begin{Verbatim}[fontsize=\footnotesize,breaklines,breakanywhere]
 \ps n
\end{Verbatim}

{\small\color{exlabel}(180)}\par
\noindent
\begin{Verbatim}[fontsize=\footnotesize,breaklines,breakanywhere]
 \sn 1a
\end{Verbatim}

{\small\color{exlabel}(181)}\par
\noindent
\begin{Verbatim}[fontsize=\footnotesize,breaklines,breakanywhere]
 \ge gloss
\end{Verbatim}

{\small\color{exlabel}(182)}\par
\noindent
\begin{Verbatim}[fontsize=\footnotesize,breaklines,breakanywhere]
 \de definition
\end{Verbatim}

{\small\color{exlabel}(183)}\par
\noindent
\begin{Verbatim}[fontsize=\footnotesize,breaklines,breakanywhere]
 \sn 1b
\end{Verbatim}

{\small\color{exlabel}(184)}\par
\noindent
\begin{Verbatim}[fontsize=\footnotesize,breaklines,breakanywhere]
 \ge gloss
\end{Verbatim}

{\small\color{exlabel}(185)}\par
\noindent
\begin{Verbatim}[fontsize=\footnotesize,breaklines,breakanywhere]
 \de definition
\end{Verbatim}

{\small\color{exlabel}(186)}\par
\noindent
\begin{Verbatim}[fontsize=\footnotesize,breaklines,breakanywhere]
 \sn 1c
\end{Verbatim}

{\small\color{exlabel}(187)}\par
\noindent
\begin{Verbatim}[fontsize=\footnotesize,breaklines,breakanywhere]
 \ge gloss
\end{Verbatim}

{\small\color{exlabel}(188)}\par
\noindent
\begin{Verbatim}[fontsize=\footnotesize,breaklines,breakanywhere]
 \de definition

\end{Verbatim}

{\small\color{exlabel}(189)}\par
\noindent
\begin{Verbatim}[fontsize=\footnotesize,breaklines,breakanywhere]
 \sn 2
\end{Verbatim}

{\small\color{exlabel}(190)}\par
\noindent
\begin{Verbatim}[fontsize=\footnotesize,breaklines,breakanywhere]
 \ge gloss
\end{Verbatim}

{\small\color{exlabel}(191)}\par
\noindent
\begin{Verbatim}[fontsize=\footnotesize,breaklines,breakanywhere]
 \de definition

\end{Verbatim}

{\small\color{exlabel}(192)}\par
\noindent
\begin{Verbatim}[fontsize=\footnotesize,breaklines,breakanywhere]
 \sn 3
\end{Verbatim}

{\small\color{exlabel}(193)}\par
\noindent
\begin{Verbatim}[fontsize=\footnotesize,breaklines,breakanywhere]
 \ge gloss
\end{Verbatim}

{\small\color{exlabel}(194)}\par
\noindent
\begin{Verbatim}[fontsize=\footnotesize,breaklines,breakanywhere]
 \de definition

\end{Verbatim}

Which would have the general printed structure of:


{\small\color{exlabel}(195)}\par
\noindent
\begin{Verbatim}[fontsize=\footnotesize,breaklines,breakanywhere]
lexeme   n. 1a) definition. 1b) definition. 1c) definition.
               2) definition. 3) definition.

\end{Verbatim}

\section{Using Subentries}

Subentries (\textbackslash{}se) provide a further level of hierarchy. These are commonly built around polymorphemic forms in a root-based dictionary (see the MDF field manual (Coward and Grimes, 1995) for an extended discussion).


{\small\color{exlabel}(196)}\par
\noindent
\begin{Verbatim}[fontsize=\footnotesize,breaklines,breakanywhere]
 \lx bren
\end{Verbatim}

{\small\color{exlabel}(197)}\par
\noindent
\begin{Verbatim}[fontsize=\footnotesize,breaklines,breakanywhere]
 \ps vi
\end{Verbatim}

{\small\color{exlabel}(198)}\par
\noindent
\begin{Verbatim}[fontsize=\footnotesize,breaklines,breakanywhere]
 \ge play
\end{Verbatim}

{\small\color{exlabel}(199)}\par
\noindent
\begin{Verbatim}[fontsize=\footnotesize,breaklines,breakanywhere]
 \ee Implies lack of focus or purpose.

\end{Verbatim}

{\small\color{exlabel}(200)}\par
\noindent
\begin{Verbatim}[fontsize=\footnotesize,breaklines,breakanywhere]
 \se brenak
\end{Verbatim}

{\small\color{exlabel}(201)}\par
\noindent
\begin{Verbatim}[fontsize=\footnotesize,breaklines,breakanywhere]
 \ps vt
\end{Verbatim}

{\small\color{exlabel}(202)}\par
\noindent
\begin{Verbatim}[fontsize=\footnotesize,breaklines,breakanywhere]
 \ge play_s.t.
\end{Verbatim}

{\small\color{exlabel}(203)}\par
\noindent
\begin{Verbatim}[fontsize=\footnotesize,breaklines,breakanywhere]
 \de play a game, or play with s.t.

\end{Verbatim}

{\small\color{exlabel}(204)}\par
\noindent
\begin{Verbatim}[fontsize=\footnotesize,breaklines,breakanywhere]
 \se inabren
\end{Verbatim}

{\small\color{exlabel}(205)}\par
\noindent
\begin{Verbatim}[fontsize=\footnotesize,breaklines,breakanywhere]
 \ps n
\end{Verbatim}

{\small\color{exlabel}(206)}\par
\noindent
\begin{Verbatim}[fontsize=\footnotesize,breaklines,breakanywhere]
 \ge recreation ; entertainment

\end{Verbatim}

{\small\color{exlabel}(207)}\par
\noindent
\begin{Verbatim}[fontsize=\footnotesize,breaklines,breakanywhere]
 \se rabrenak
\end{Verbatim}

{\small\color{exlabel}(208)}\par
\noindent
\begin{Verbatim}[fontsize=\footnotesize,breaklines,breakanywhere]
 \ps n
\end{Verbatim}

{\small\color{exlabel}(209)}\par
\noindent
\begin{Verbatim}[fontsize=\footnotesize,breaklines,breakanywhere]
 \ge toy
\end{Verbatim}

{\small\color{exlabel}(210)}\par
\noindent
\begin{Verbatim}[fontsize=\footnotesize,breaklines,breakanywhere]
 \dt 17/Jun/92

\end{Verbatim}

This would print like the following:


{\small\color{exlabel}(211)}\par
\noindent
\begin{Verbatim}[fontsize=\footnotesize,breaklines,breakanywhere]
bren     vi. play. Implies lack of focus or purpose.
     brenak    vt. play a game, or play with s.t.
     inabren    n. recreation, entertainment.
     rabrenak n. toy.

\end{Verbatim}

A more complete example using subentries:


{\small\color{exlabel}(212)}\par
\noindent
\begin{Verbatim}[fontsize=\footnotesize,breaklines,breakanywhere]
 \lx bersih
\end{Verbatim}

{\small\color{exlabel}(213)}\par
\noindent
\begin{Verbatim}[fontsize=\footnotesize,breaklines,breakanywhere]
 \ps adj
\end{Verbatim}

{\small\color{exlabel}(214)}\par
\noindent
\begin{Verbatim}[fontsize=\footnotesize,breaklines,breakanywhere]
 \sn 1
\end{Verbatim}

{\small\color{exlabel}(215)}\par
\noindent
\begin{Verbatim}[fontsize=\footnotesize,breaklines,breakanywhere]
 \ge clean
\end{Verbatim}

{\small\color{exlabel}(216)}\par
\noindent
\begin{Verbatim}[fontsize=\footnotesize,breaklines,breakanywhere]
 \de be clean, not dirty or messy
\end{Verbatim}

{\small\color{exlabel}(217)}\par
\noindent
\begin{Verbatim}[fontsize=\footnotesize,breaklines,breakanywhere]
 \sn 2
\end{Verbatim}

{\small\color{exlabel}(218)}\par
\noindent
\begin{Verbatim}[fontsize=\footnotesize,breaklines,breakanywhere]
 \ge innocent
\end{Verbatim}

{\small\color{exlabel}(219)}\par
\noindent
\begin{Verbatim}[fontsize=\footnotesize,breaklines,breakanywhere]
 \de be innocent, without fault

\end{Verbatim}

{\small\color{exlabel}(220)}\par
\noindent
\begin{Verbatim}[fontsize=\footnotesize,breaklines,breakanywhere]
 \se kebersihan
\end{Verbatim}

{\small\color{exlabel}(221)}\par
\noindent
\begin{Verbatim}[fontsize=\footnotesize,breaklines,breakanywhere]
 \ps n
\end{Verbatim}

{\small\color{exlabel}(222)}\par
\noindent
\begin{Verbatim}[fontsize=\footnotesize,breaklines,breakanywhere]
 \ge cleanliness

\end{Verbatim}

{\small\color{exlabel}(223)}\par
\noindent
\begin{Verbatim}[fontsize=\footnotesize,breaklines,breakanywhere]
 \se membersihkan
\end{Verbatim}

{\small\color{exlabel}(224)}\par
\noindent
\begin{Verbatim}[fontsize=\footnotesize,breaklines,breakanywhere]
 \ps vt
\end{Verbatim}

{\small\color{exlabel}(225)}\par
\noindent
\begin{Verbatim}[fontsize=\footnotesize,breaklines,breakanywhere]
 \sn 1
\end{Verbatim}

{\small\color{exlabel}(226)}\par
\noindent
\begin{Verbatim}[fontsize=\footnotesize,breaklines,breakanywhere]
 \ge clean_up
\end{Verbatim}

{\small\color{exlabel}(227)}\par
\noindent
\begin{Verbatim}[fontsize=\footnotesize,breaklines,breakanywhere]
 \de clean s.t. up
\end{Verbatim}

{\small\color{exlabel}(228)}\par
\noindent
\begin{Verbatim}[fontsize=\footnotesize,breaklines,breakanywhere]
 \sn 2
\end{Verbatim}

{\small\color{exlabel}(229)}\par
\noindent
\begin{Verbatim}[fontsize=\footnotesize,breaklines,breakanywhere]
 \ge purify
\end{Verbatim}

{\small\color{exlabel}(230)}\par
\noindent
\begin{Verbatim}[fontsize=\footnotesize,breaklines,breakanywhere]
 \de purify, repent or renounce immoral actions

\end{Verbatim}

{\small\color{exlabel}(231)}\par
\noindent
\begin{Verbatim}[fontsize=\footnotesize,breaklines,breakanywhere]
 \se pembersih
\end{Verbatim}

{\small\color{exlabel}(232)}\par
\noindent
\begin{Verbatim}[fontsize=\footnotesize,breaklines,breakanywhere]
 \ps n
\end{Verbatim}

{\small\color{exlabel}(233)}\par
\noindent
\begin{Verbatim}[fontsize=\footnotesize,breaklines,breakanywhere]
 \sn 1
\end{Verbatim}

{\small\color{exlabel}(234)}\par
\noindent
\begin{Verbatim}[fontsize=\footnotesize,breaklines,breakanywhere]
 \ge cleanser
\end{Verbatim}

{\small\color{exlabel}(235)}\par
\noindent
\begin{Verbatim}[fontsize=\footnotesize,breaklines,breakanywhere]
 \sn 2
\end{Verbatim}

{\small\color{exlabel}(236)}\par
\noindent
\begin{Verbatim}[fontsize=\footnotesize,breaklines,breakanywhere]
 \ge janitor
\end{Verbatim}

{\small\color{exlabel}(237)}\par
\noindent
\begin{Verbatim}[fontsize=\footnotesize,breaklines,breakanywhere]
 \dt 17/Jun/92

\end{Verbatim}

Which would print as:


{\small\color{exlabel}(238)}\par
\noindent
\begin{Verbatim}[fontsize=\footnotesize,breaklines,breakanywhere]
bersih adj. 1) be clean, not dirty or messy. 2) be innocent, 
           without fault.
    kebersihan n. cleanliness.
    membersihkan vt. 1) clean s.t. up. 2) purify, repent or
           renounce immoral actions.
    pembersih n. 1) cleanser. 2) janitor.

\end{Verbatim}

For information on  MDF's alternate hierarchy, see:


{\color{cflink} \hyperref[key:Alternate_Hierarchy]{Alternate\_Hierarchy}}


To continue with the general discussion, jump to:


{\color{cflink} \hyperref[key:Using_Subentries_or_Lexical_Entries]{Using\_Subentries\_or\_Lexical\_Entries}}


\chapter{Printed Field Labels}\label{key:Printed_Field_Labels}

MDF adds field labels to many of the fields as it formats your lexicon for printing.  These  labels are explicitly noted in the discussion of each of the field markers. MDF allows for these field labels to be changed for either English or National audiences.


By default, MDF uses the following two CC tables to make the label substitutions:


{\small\color{exlabel}(239)}\par
\noindent
\begin{Verbatim}[fontsize=\footnotesize,breaklines,breakanywhere]
   MDF_eng.cct       (for English labels)
\end{Verbatim}

{\small\color{exlabel}(240)}\par
\noindent
\begin{Verbatim}[fontsize=\footnotesize,breaklines,breakanywhere]
   MDF_inz.cct       (for national (Indonesian)language labels)

\end{Verbatim}

MDF also comes with other label tables. For example, for French and Spanish:


{\small\color{exlabel}(241)}\par
\noindent
\begin{Verbatim}[fontsize=\footnotesize,breaklines,breakanywhere]
   MDF_frn.cct       (for French labels)
\end{Verbatim}

{\small\color{exlabel}(242)}\par
\noindent
\begin{Verbatim}[fontsize=\footnotesize,breaklines,breakanywhere]
   MDF_spn.cct       (for Spanish labels)

\end{Verbatim}

These CC tables are in the MDF folder, just under the Toolbox program folder.


\begin{quote}\small\itshape For MDF to work properly, all of MDF's CC tables need to be copied to your Projects folder. (This is the folder where you keep your project (PRJ) files, langauge (LNG) files, and database type (TYP) files.)\end{quote}

Try one of these ``canned'' CC tables first, if these won't meet your needs, you can edit them so that they put on the exact labels that you need.


\begin{quote}\small\itshape These files are specified in the File-Export-MDF-Options-Audience dialog box.\end{quote}

To change a field label:


{\small\color{exlabel}(243)}\par
\noindent
\begin{Verbatim}[fontsize=\footnotesize,breaklines,breakanywhere]
   1. IMPORTANT: First, using File Manager or Explorer, to copy the 
      CCT file you want to change to a new filename.
\end{Verbatim}

{\small\color{exlabel}(244)}\par
\noindent
\begin{Verbatim}[fontsize=\footnotesize,breaklines,breakanywhere]
   2. Next open the new CCT file in Notepad.
\end{Verbatim}

{\small\color{exlabel}(245)}\par
\noindent
\begin{Verbatim}[fontsize=\footnotesize,breaklines,breakanywhere]
   3. At the end of this file you will find a section marked as 
      "group(gLabel)"; under here are all of the field labels that
      MDF adds to the beginning of fields. The following are a few
      of the label-changing lines from one of these files:

"Ant:"            > "Lawan:"                c \an Antonym
"Read:"           > "Baca:"                 c \bb Bibliography
"From:"           > "Pinjaman:"             c \bw Borrowed word (loan)

\end{Verbatim}

{\small\color{exlabel}(246)}\par
\noindent
\begin{Verbatim}[fontsize=\footnotesize,breaklines,breakanywhere]
   4. Carefully change each label in the "Lawan:" column that needs
      changing. Be careful you do NOT delete the " marks! 
\end{Verbatim}

{\small\color{exlabel}(247)}\par
\noindent
\begin{Verbatim}[fontsize=\footnotesize,breaklines,breakanywhere]
   5. Save the file to the new name (be careful that Notepad doesn't
      stick on a .TXT extension on the file --if it does, rename it).
\end{Verbatim}

{\small\color{exlabel}(248)}\par
\noindent
\begin{Verbatim}[fontsize=\footnotesize,breaklines,breakanywhere]
   6. Now go back to Toolbox and the MDF Export-Options-Audience box 
      to give the name of the new CC table you created.

\end{Verbatim}

In these CC tables  there is also options for making consistent changes to your lexical function (\textbackslash{}lf) labels, part of speech (\textbackslash{}ps) labels, and paradigm form (\textbackslash{}pdl) labels. This allows you to keep your Range Sets labels for these all in one language, but when the dictionary is printed for another audience, the labels are changed appropriately.


For a detailed description of the MDF formatting process, see:


{\color{cflink} \hyperref[key:Formatting_and_Printing]{Formatting\_and\_Printing}}


For related information on how MDF can also substitute Range Set labels, see:


{\color{cflink} \hyperref[key:Range_Sets]{Range\_Sets}}


To continue with the general discussion, jump to:


{\color{cflink} \hyperref[key:Unknown_Fields]{Unknown\_Fields}}


\chapter{Using Subentries, Lexical Entries, and Minimal Entries}\label{key:Using_Subentries_or_Lexical_Entries}

A common question that all lexicographers must face when wrestling with complex (polymorphemic) lexemes is whether to incorporate them in the lexicon as subentries under the lexical entry of their root lexeme, or to place them in the lexicon as independent lexical entries. This is not an easy question to answer, but it is addressed in greater detail in Chapter 4 (section 4.6) of the MDF field manual (Coward and Grimes, 1995). For this database, we will summarize the two options below:


\section{1. Using Subentries and Minimal Entries}

If you want to keep all related (root and polymorphemic) lexemes together, use the subentry method shown below (this is a simplified example):


{\small\color{exlabel}(249)}\par
\noindent
\begin{Verbatim}[fontsize=\footnotesize,breaklines,breakanywhere]
 \lx bren
\end{Verbatim}

{\small\color{exlabel}(250)}\par
\noindent
\begin{Verbatim}[fontsize=\footnotesize,breaklines,breakanywhere]
 \ps vi
\end{Verbatim}

{\small\color{exlabel}(251)}\par
\noindent
\begin{Verbatim}[fontsize=\footnotesize,breaklines,breakanywhere]
 \ge play
\end{Verbatim}

{\small\color{exlabel}(252)}\par
\noindent
\begin{Verbatim}[fontsize=\footnotesize,breaklines,breakanywhere]
 \ee Implies lack of focus or purpose.

\end{Verbatim}

{\small\color{exlabel}(253)}\par
\noindent
\begin{Verbatim}[fontsize=\footnotesize,breaklines,breakanywhere]
 \se brenak
\end{Verbatim}

{\small\color{exlabel}(254)}\par
\noindent
\begin{Verbatim}[fontsize=\footnotesize,breaklines,breakanywhere]
 \ps vt
\end{Verbatim}

{\small\color{exlabel}(255)}\par
\noindent
\begin{Verbatim}[fontsize=\footnotesize,breaklines,breakanywhere]
 \ge play_s.t.
\end{Verbatim}

{\small\color{exlabel}(256)}\par
\noindent
\begin{Verbatim}[fontsize=\footnotesize,breaklines,breakanywhere]
 \de play a game, or play with s.t.

\end{Verbatim}

{\small\color{exlabel}(257)}\par
\noindent
\begin{Verbatim}[fontsize=\footnotesize,breaklines,breakanywhere]
 \se inabren
\end{Verbatim}

{\small\color{exlabel}(258)}\par
\noindent
\begin{Verbatim}[fontsize=\footnotesize,breaklines,breakanywhere]
 \ps n
\end{Verbatim}

{\small\color{exlabel}(259)}\par
\noindent
\begin{Verbatim}[fontsize=\footnotesize,breaklines,breakanywhere]
 \ge recreation ; entertainment

\end{Verbatim}

{\small\color{exlabel}(260)}\par
\noindent
\begin{Verbatim}[fontsize=\footnotesize,breaklines,breakanywhere]
 \se rabrenak
\end{Verbatim}

{\small\color{exlabel}(261)}\par
\noindent
\begin{Verbatim}[fontsize=\footnotesize,breaklines,breakanywhere]
 \ps n
\end{Verbatim}

{\small\color{exlabel}(262)}\par
\noindent
\begin{Verbatim}[fontsize=\footnotesize,breaklines,breakanywhere]
 \ge toy
\end{Verbatim}

{\small\color{exlabel}(263)}\par
\noindent
\begin{Verbatim}[fontsize=\footnotesize,breaklines,breakanywhere]
 \dt 17/Jun/92

\end{Verbatim}

This would print as:


{\small\color{exlabel}(264)}\par
\noindent
\begin{Verbatim}[fontsize=\footnotesize,breaklines,breakanywhere]
bren     vi. play. Implies lack of focus or purpose.
     brenak    vt. play a game, or play with s.t.
     inabren    n. recreation, entertainment.
     rabrenak n. toy.

\end{Verbatim}

This method allows you to see all of the related pieces close together. But you have to know where to look. And this is a major problem; outsiders will not know that to find a discussion of ``inabren'' they must look under ``bren''. To make this easier for new users, we recommend you also make extensive use of ``Minimal Entries'' (essentially the same as ``minor entries'') -- these are lexical entries that have only the most basic fields (maybe a \textbackslash{}ps field and a gloss field), but they include the \textbackslash{}mn ``main entry'' field which provides the sign-post back to the lexical entry where more information can be found.


The minimal entry for ``inabren'' might be:


{\small\color{exlabel}(265)}\par
\noindent
\begin{Verbatim}[fontsize=\footnotesize,breaklines,breakanywhere]
  \lx inabren
\end{Verbatim}

{\small\color{exlabel}(266)}\par
\noindent
\begin{Verbatim}[fontsize=\footnotesize,breaklines,breakanywhere]
  \ps n
\end{Verbatim}

{\small\color{exlabel}(267)}\par
\noindent
\begin{Verbatim}[fontsize=\footnotesize,breaklines,breakanywhere]
  \ge recreation
\end{Verbatim}

{\small\color{exlabel}(268)}\par
\noindent
\begin{Verbatim}[fontsize=\footnotesize,breaklines,breakanywhere]
  \mn bren

\end{Verbatim}

When printed the \textbackslash{}mn field typically has the label ``See main entry:''.


That is all there is to this minimal entry. For this reason, it is expected that the subentry for ``inabren'' would look very much like a full lexical entry; it would contain verbose definitions, have example sentences, multiple senses (if needed), cross-referencing, and any encyclopedic information, notes, or comments that are needed to elucidate this lexeme, i.e. it would be as though a full lexical entry were stuck right inside another lexical entry. (The above example of subentry use obviously falls short of this expectation -- it is merely an example.)


But practically speaking, many find it difficult to actually put the same amount of effort into embellishing a subentry as they would into a main lexical entry. For this reason, and for the fact that such embellished records soon become enormous, you might find the alternative way better for you:


\section{2. Lexical Entries and Cross-referencing}

The alternative is nearly the exact opposite: make all lexemes (roots or polymorphemic forms) main lexical entries. But in order to see the relationships between the scattered pieces, you will need to use an extensive cross-referencing system. MDF provides the \textbackslash{}cf and the \textbackslash{}lf field bundles for this.


Again using the above example, the ``bren'' lexical entry would include the following fields:


{\small\color{exlabel}(269)}\par
\noindent
\begin{Verbatim}[fontsize=\footnotesize,breaklines,breakanywhere]
  \cf brenak
\end{Verbatim}

{\small\color{exlabel}(270)}\par
\noindent
\begin{Verbatim}[fontsize=\footnotesize,breaklines,breakanywhere]
  \ce play s.t.

\end{Verbatim}

{\small\color{exlabel}(271)}\par
\noindent
\begin{Verbatim}[fontsize=\footnotesize,breaklines,breakanywhere]
  \cf inabren
\end{Verbatim}

{\small\color{exlabel}(272)}\par
\noindent
\begin{Verbatim}[fontsize=\footnotesize,breaklines,breakanywhere]
  \ce recreation

\end{Verbatim}

{\small\color{exlabel}(273)}\par
\noindent
\begin{Verbatim}[fontsize=\footnotesize,breaklines,breakanywhere]
  \cf rabrenak
\end{Verbatim}

{\small\color{exlabel}(274)}\par
\noindent
\begin{Verbatim}[fontsize=\footnotesize,breaklines,breakanywhere]
  \ce toy

\end{Verbatim}

These \textbackslash{}cf fields reference full lexical entries for both ``inabren'' and for ``rabrenak''. If you wish to provide even further information regarding the relationship between the ``bren'' lexeme and these referenced lexemes, use the \textbackslash{}lf field (the \textbackslash{}lf labels shown here are not certain):


{\small\color{exlabel}(275)}\par
\noindent
\begin{Verbatim}[fontsize=\footnotesize,breaklines,breakanywhere]
  \lf Vgoal
\end{Verbatim}

{\small\color{exlabel}(276)}\par
\noindent
\begin{Verbatim}[fontsize=\footnotesize,breaklines,breakanywhere]
  \lv brenak
\end{Verbatim}

{\small\color{exlabel}(277)}\par
\noindent
\begin{Verbatim}[fontsize=\footnotesize,breaklines,breakanywhere]
  \le play s.t.

\end{Verbatim}

{\small\color{exlabel}(278)}\par
\noindent
\begin{Verbatim}[fontsize=\footnotesize,breaklines,breakanywhere]
  \lf Ngoal
\end{Verbatim}

{\small\color{exlabel}(279)}\par
\noindent
\begin{Verbatim}[fontsize=\footnotesize,breaklines,breakanywhere]
  \lv inabren
\end{Verbatim}

{\small\color{exlabel}(280)}\par
\noindent
\begin{Verbatim}[fontsize=\footnotesize,breaklines,breakanywhere]
  \le recreation

\end{Verbatim}

{\small\color{exlabel}(281)}\par
\noindent
\begin{Verbatim}[fontsize=\footnotesize,breaklines,breakanywhere]
  \lf Ninst
\end{Verbatim}

{\small\color{exlabel}(282)}\par
\noindent
\begin{Verbatim}[fontsize=\footnotesize,breaklines,breakanywhere]
  \lv rabrenak
\end{Verbatim}

{\small\color{exlabel}(283)}\par
\noindent
\begin{Verbatim}[fontsize=\footnotesize,breaklines,breakanywhere]
  \le toy

\end{Verbatim}

This strategy keeps the individual records smaller and still provides a way to see what all the related forms are and their general meanings and relationships. But for more detailed information about a related lexeme requires searching through the lexicon (easy  in Toolbox, but tedious when in book form).


You will need to decide which strategy will work best for you.


For more information on sections in a lexical entry, see:


{\color{cflink} \hyperref[key:Sections_in_a_Lexical_Entry]{Sections\_in\_a\_Lexical\_Entry}}


{\color{cflink} \hyperref[key:Alternate_Hierarchy]{Alternate\_Hierarchy}}


For information on related fields, see:


\begin{description}
\item[{\hyperref[key:mn]{\textcolor{cflink}{\textbackslash{}mn}}}] main entry form
\item[{\hyperref[key:cf]{\textcolor{cflink}{\textbackslash{}cf}}}] general cross-reference
\item[{\hyperref[key:lf]{\textcolor{cflink}{\textbackslash{}lf}}}] lexical function cross-reference
\end{description}


To continue with the general discussion, jump to:


{\color{cflink} \hyperref[key:Free-form_Fields]{Free-form\_Fields}}


\chapter{Unknown Fields}\label{key:Unknown_Fields}

Any field marker in a lexical database that is not recognized by MDF is considered an unknown or Non-MDF Field.


\begin{quote}\small\itshape Normally MDF discards unknown fields from the formatted output.\end{quote}

While it is always best to use available fields where possible (for the sake of functionality and compatibility), if you should need to create a field not covered by the 102 MDF fields (say, for some special bookkeeping or analysis chore), you may do so without worry. As a general rule, any field marker of three or four letters is a non-MDF field. The only exceptions to this rule are the new paradigm fields: \textbackslash{}pdl, \textbackslash{}pdv, \textbackslash{}pde, \textbackslash{}pdn, and \textbackslash{}pdr.


\begin{quote}\small\itshape If you want MDF to include non-MDF fields in the formatted output, simply check ``Include non-MDF fields'' under Options in the MDF Export dialog box.\end{quote}

All unknown fields are collected together and placed at the end of the section they occur in. They are bracketed with [ ] and marked with the label: ??.


To continue with the general discussion, jump to:


{\color{cflink} \hyperref[key:Formatting_and_Printing]{Formatting\_and\_Printing}}


\chapter{A Sense-Oriented Hierarchy}\label{key:Alternate_Hierarchy}

The alternative hierarchy that MDF supports differs from the standard hierarchy in its ordering of these three crucial field markers: \textbackslash{}sn \textbackslash{}se \textbackslash{}ps. (The standard hierarchy expects the order \textbackslash{}ps \textbackslash{}sn \textbackslash{}se.) The alternate hierarchy's ordering can be outlined as follows:


{\small\color{exlabel}(284)}\par
\noindent
\begin{Verbatim}[fontsize=\footnotesize,breaklines,breakanywhere]
   \lx lexeme
      \sn sense number
         \se subentry
            \ps part-of-speech
            \ps part-of-speech
         \se subentry
            \ps part-of-speech
            \ps part-of-speech
      \sn sense number
         \se subentry
            \ps part-of-speech
            \ps part-of-speech
         \se subentry
            \ps part-of-speech
            \ps part-of-speech

\end{Verbatim}

The main purpose of this alternate hierarchy is to allow the user to group or base subentries on a given sense. A simplified example of this would be:


{\small\color{exlabel}(285)}\par
\noindent
\begin{Verbatim}[fontsize=\footnotesize,breaklines,breakanywhere]
 \lx adá
 \ge throw ; discard
 
 \sn 1
 \de throw
 
 \se mongadá
 \ps vt 
 \ge throw away ; discard
 \de throw something away; to discard

 \se iadá
 \ps vt
 \de throw something away

 \se paadá
 \mr po-adá 
 \ps vCAUS
 \de drop

 \se puadá-adá
 \mr pu-adá-adá
 \ps v
 \de throw away (thoughtlessly or irresponsibly)

 \sn 2
 \de divorce

 \se muadá
 \mr m-pu-adá
 \ps vREC
 \de divorce smb (lit. throw each other)

 \se puadá
 \mr pu-adá
 \ps vREC
 \de be divorced; separated (used for divorce and separation of man and wife)

 \sn 3
 \de defecate

 \se paadá
 \ps vCAUS
 \de defecate

\end{Verbatim}

This example would typically print as follows (shown without the character styles here):


{\small\color{exlabel}(286)}\par
\noindent
\begin{Verbatim}[fontsize=\footnotesize,breaklines,breakanywhere]
   adá   throw, discard.
      1) throw.
            mongadá  vt. throw something away; to discard. 
            iadá   vt. throw something away.
            paadá  Morf: po-adá. vCAUS. drop.
            puadá-adá   v. throw away (thoughtlessly or irresponsibly).
      2) divorce.
            muadá   Morf: m-pu-adá. vREC. divorce smb (lit. throw 
                each other).
            puadá   Morf: pu-adá. vREC. be divorced; separated (used 
                for divorce and separation of man and wife). 
      3) defecate.
            paadá   vCAUS. defecate.

\end{Verbatim}

Note that by judiciously using the definition (\textbackslash{}de) field you can give labels to the sense groupings.


This method of organization was urgently requested by people working in languages that were highly polymorphemic, because this hierarchy allows them to group complex polymorphemic lexemes together based on a sense (rather than stand alone -- as subentries do in the standard MDF hierarchy).


\begin{quote}\small\itshape Currently MDF can only support one method of organization per lexical database. Therefore you must choose which hierarchical method will best meet your needs and use only that one method for all lexical entries in you dictionary. Mixing methods in Shoebox and MDF can produce inaccurate interlinear choices and browse views, as well as odd formatting in the output.\end{quote}

To help with the decision, note that there are advantages and disadvantages to using this alternate method:


\begin{itemize}
  \item[] •    One main advantage to this structure is that you can see all the related polymorphemic forms that carry a particular sense.
  \item[] •    One disadvantage of this structure is that you do not readily see the actual range of senses that a particular polymorphemic form can express.
\end{itemize}

For important information on  how MDF handles sections in a lexical entry as well as related information on the standard hierarchy, please see:


{\color{cflink} \hyperref[key:Sections_in_a_Lexical_Entry]{Sections\_in\_a\_Lexical\_Entry}}


For information on the order in which fields are formatted for printing, see:


{\color{cflink} \hyperref[key:Order_of_Fields]{Order\_of\_Fields}}


To continue on with the general discussion, jump to:


{\color{cflink} \hyperref[key:Using_Subentries_or_Lexical_Entries]{Using\_Subentries\_or\_Lexical\_Entries}}


Also see the document ``MDF-ch2 revised.doc'', pages 55-63 for a detailed discussion of the two different hierarchies.


\chapter{Formatting and Printing Your Dictionary and Finderlists}\label{key:Formatting_and_Printing}

So your lexicon is ready, and you want to print it. This assumes you've set up the database to be an MDF type database . (An MDF type database not only uses the MDF marker set but has checked the box in Database-Properties-Options labeled ``Multi-Dictionary Formatter'').


\section{Specifying the Options}

To print  your lexicon, choose ``Multi-Dictionary Formatter'' from the File-Export menu. Here you will be able to select the intended audience (either English or national), the format (diglot, triglot, or a reversal finderlist), specify text for the footer for each of these choices, select a filter (to limit the formatted document to a smaller set of records), and specify whether you want Word opened automatically (for Windows users).


This dialog box also has an Options button which will allow you to set the basic parameters the standard MDF print feature uses to create the formatted dictionary or finderlists. These options include selecting specific fields to include or exclude and some basic page formatting.


Clicking Export gives you a chance to specify a filename and location for the new document. Clicking OK again starts the whole process. MDF goes through several steps to format the document (detailed below). The final product is a Rich Text Format (RTF) file.


\section{Open the RTF file in Word}

If you checked ``Automatically open document in word  processor'', Toolbox causes the word processor to automatically open the exported document; otherwise you must open the RTF file yourself.


\section{Finish the process of exporting from Toolbox}

To make the final formatting touches, choose the menu option ``Finish exporting from Shoebox'' from under the Tools menu in Word. This runs the FinishExportingFromShoebox macro in the document template. It updates style formatting and does a few other special formatting operations. It also saves the exported file as a Word DOC file.


\begin{quote}\small\itshape Note: Word processor programs other than Microsoft Word may not be able to use the Word macro, and may even require a different template than is supplied with the standard MDF printing process.\end{quote}

\begin{quote}\small\itshape Note: Probably the most significant thing the Finish exporting from Shoebox macro does is to load all the pictures.\end{quote}

If you need to change any style formatting, click on the text that needs changing, choose Format, Style, and Modify. Make the changes needed, and remember to check the ``Add  to Template'' box before you choose OK so the style changes will be available the next time you export.


\section{Print it}

If the document looks fine, choose File Print.


\begin{quote}\small\itshape Note: Errors in the data should be corrected in Toolbox (and export again). If you correct errors in the Word document, they will still be there in the data when you export the next time.\end{quote}

\_\_\_\_\_\_\_\_\_\_\_\_\_\_\_


\begin{quote}\small\itshape This concludes the general discussion of how to use MDF.\end{quote}

For more information, see:


{\color{cflink} \hyperref[key:References]{References}}


To return to the top of the database, simply close all unneeded windows, and then use the Database-First Record menu option, or use the First Record button on the tool bar (or search for `aa').


To see the description of any particular field, simply right-click on its marker in either of these lists (or anywhere else in this Help file):


{\color{cflink} \hyperref[key:Order_of_Fields]{Order\_of\_Fields}}


{\color{cflink} \hyperref[key:Summary_of_Fields]{Summary\_of\_Fields}}


\chapter{Free-form Fields}\label{key:Free-form_Fields}

In the standard MDF printing, the following fields are ``free-form'' fields, meaning that MDF will not add periods, etc. to them; they are basically formatted as you entered them. Therefore, when using these fields, you are encouraged to use capitalization and proper punctuation. Free-form fields may be as long as you need, e.g. several sentences.


{\small\color{exlabel}(287)}\par
\noindent
\begin{Verbatim}[fontsize=\footnotesize,breaklines,breakanywhere]
 \rf  reference                * \pc  picture
\end{Verbatim}

{\small\color{exlabel}(288)}\par
\noindent
\begin{Verbatim}[fontsize=\footnotesize,breaklines,breakanywhere]
 \xv  vernacular example         \nt  notes, general
\end{Verbatim}

{\small\color{exlabel}(289)}\par
\noindent
\begin{Verbatim}[fontsize=\footnotesize,breaklines,breakanywhere]
 \xe  English example            \np  notes on phonology
\end{Verbatim}

{\small\color{exlabel}(290)}\par
\noindent
\begin{Verbatim}[fontsize=\footnotesize,breaklines,breakanywhere]
 \xn  national example           \ng  notes on grammar
\end{Verbatim}

{\small\color{exlabel}(291)}\par
\noindent
\begin{Verbatim}[fontsize=\footnotesize,breaklines,breakanywhere]
 \xr  regional example           \nd  notes on discourse
\end{Verbatim}

{\small\color{exlabel}(292)}\par
\noindent
\begin{Verbatim}[fontsize=\footnotesize,breaklines,breakanywhere]
 \ve  variant form (English)     \na  notes on anthro.
\end{Verbatim}

{\small\color{exlabel}(293)}\par
\noindent
\begin{Verbatim}[fontsize=\footnotesize,breaklines,breakanywhere]
 \vn  variant form (National)    \ns  notes on socioling.
\end{Verbatim}

{\small\color{exlabel}(294)}\par
\noindent
\begin{Verbatim}[fontsize=\footnotesize,breaklines,breakanywhere]
 \vr  variant form (Regional)    \nq  questions
\end{Verbatim}

{\small\color{exlabel}(295)}\par
\noindent
\begin{Verbatim}[fontsize=\footnotesize,breaklines,breakanywhere]
 \ur  usage (regional)
 \or  only (regional) 
 \er  encyclopedic (Regional) 
\end{Verbatim}

{\small\color{exlabel}(296)}\par
\noindent
\begin{Verbatim}[fontsize=\footnotesize,breaklines,breakanywhere]
 \tb  table

\end{Verbatim}

The following fields are also free-form, except that the printing process will add punctuation to the end if you do not.


{\small\color{exlabel}(297)}\par
\noindent
\begin{Verbatim}[fontsize=\footnotesize,breaklines,breakanywhere]
 \uv  usage (vernacular)         \ov  only (vernacular)
 \ue  usage (English)            \oe  only (English)
 \un  usage (national)           \on  only (national)
 
 \ev  encyclopedic (vern)     
 \ee  encyclopedic (English)  
 \en  encyclopedic (National) 
 
\end{Verbatim}

\begin{quote}\small\itshape If a field is not explicitly marked in its description as being a ``free-form'' field, then it is not one.\end{quote}

If a field is not a free-form field, you do not need to use any terminating punctuation; MDF adds the needed punctuation to  the end of each non-free-form field.


* The picture reference field is free-form so long as the field does not begin with .G. (which marks it as a graphics link paragraph). See \textbackslash{}pc for more information on the picture reference field.


To continue on with the general discussion, jump to:


{\color{cflink} \hyperref[key:Range_Sets]{Range\_Sets}}


\chapter{References}\label{key:References}

{\small\setlength{\parindent}{0pt}\hangindent=3em\hangafter=1 Coward, David F. and Charles E. Grimes. 1995. Making Dictionaries: a guide to lexicography and MDF. SIL: Waxhaw, NC.\par}

{\small\setlength{\parindent}{0pt}\hangindent=3em\hangafter=1 Grimes, Charles E. 1987. Mapping a culture through networks of meaning. Notes on Linguistics. 39:25-36. SIL: Dallas.\par}

{\small\setlength{\parindent}{0pt}\hangindent=3em\hangafter=1 Louw, Johannes P., and Eugene A. Nida (eds.). 1988. Greek-English lexicon of the New Testament based on semantic domains. United Bible Societies: New York.\par}

{\small\setlength{\parindent}{0pt}\hangindent=3em\hangafter=1 Nivens, Rick. 1990. Beyond key terms: a lexicon useful for translation. Manuscript.\par}

\chapter{The MDF Documentation}\label{key:The_MDF_Documentation}

The original MDF manual,  ``Making Dictionaries: A guide to lexicography and the Multi-Dictionary Formatter'', by Coward and Grimes, was published in 1995 and gave instructions for a set of independent CC tables running in the DOS operating system. While still an extremely valuable document, there are parts which are no longer relevant, and some which have been changed.


The manual itself is only available as a PDF file, though the revision of chapter two is a Word document. They can be obtained from ...


The following are suggestions for reading the manual:


Chapter 1 - Skip (specifics of installation of the original DOS version)


Chapter 2 - see revised version: MDF-ch2-revised.doc. Has a description of all MDF markers. Ignore references to ``older approaches''. The section on gloss vs reversal vs definition is useful.


Chapter 3 - Skip. Specifics of running the original DOS version.


Chapter 4 Theory, which needs to be considered.


Chapter 5 - Read section 5.3 about categories of markers. The rest is mostly about why to use Shoebox instead of Word, and about sorting, which is done differently in Toolbox.


Chapter 6 discusses: principles for choosing the headword, example sentences, homonym vs different senses, semantic categories, dialect information. Note, sometimes gives alternative ways to handle the same information. Check with a dictionary consultant if you have doubts.


Chapter 7 Various possible lexical function relationships. Usefulness of semantic relationships as an approach for eliciting data.


Chapter 8 Lots of good theory on making good entries; good qestions to ask. Special considerations for activities \& events, states \& processes, loans \& etymologies, ritual speech, etc.


Chapter 9  Discussion on part of speech.


Chapter 10  A dictionary is more than a list of words. Discusses preface material, appendices, etc.


Appendices A-E include lists of markers, recommended order of fields, lists of semantic domains and lexical functions, and some commonly used abbreviations.


Appendix F-I - Skip. Specific to the original version. Not relevant now.


\chapter{Old and Changed Markers}\label{key:Old_and_Changed_Markers}

Within a couple years of the publication of the MDF manual in 1995, as the number of users grew and became more widespread, there were recommendations which modified the MDF marker set. These changes were reflected in the modification to the MDF manual,  ``MDF-ch2  revised.doc''.


In this database, some of the long-since un-recommended markers have been removed from prominent display. (The MDF tables have not been changed.) In addition, the database type distributed with the Toolbox Training Materials and the Toolbox New Project Installer no longer includes these markers. These should all be replaced by the paradigm markers.


\textbackslash{}1s    \textbackslash{}2s    \textbackslash{}3s    \textbackslash{}4s    \textbackslash{}1d    \textbackslash{}2d    \textbackslash{}3d    \textbackslash{}4d    \textbackslash{}1p    \textbackslash{}1i    \textbackslash{}1e    \textbackslash{}2p    \textbackslash{}3p    \textbackslash{}4p


For example,


{\small\color{exlabel}(298)}\par
\noindent
\begin{Verbatim}[fontsize=\footnotesize,breaklines,breakanywhere]
   \1s yoban
\end{Verbatim}

should be replaced by


{\small\color{exlabel}(299)}\par
\noindent
\begin{Verbatim}[fontsize=\footnotesize,breaklines,breakanywhere]
   \pdl 1s
\end{Verbatim}

{\small\color{exlabel}(300)}\par
\noindent
\begin{Verbatim}[fontsize=\footnotesize,breaklines,breakanywhere]
   \pdv yoban
\end{Verbatim}

In this example, the label which will be used in printing is placed in the \textbackslash{}pdl (paradigm label) field. The paradigm form itself is placed in the \textbackslash{}pdv (paradigm vernacular form) field. The paradigm approach has much greater flexibility, both in the choice of labels and in allowing glosses in other languages.


In addition, the following are still included in the MDF database type file but could be replaced by the paradigm approach.


\textbackslash{}sg Singular form          \textbackslash{}pl Plural form          \textbackslash{}rd Reduplication form


Similarly, these markers could use the lexical function approach.


\textbackslash{}an Antonym          \textbackslash{}sy Synonym


For examples and more details, see


\begin{description}
\item[{\hyperref[key:lf]{\textcolor{cflink}{\textbackslash{}lf}}}] the lexical function label
\item[{\hyperref[key:pd]{\textcolor{cflink}{\textbackslash{}pd}}}] paradigm set
\item[{\hyperref[key:pdl]{\textcolor{cflink}{\textbackslash{}pdl}}}] the paradigm label
\end{description}


Be sure to also see


{\color{cflink} \hyperref[key:The_MDF_Documentation]{The\_MDF\_Documentation}}


\chapter{Old verb paradigm markers}\label{key:1s-1p-1e-1i-1d-2s-2p-2d-3s-3p-3d-4s-4p-4d}

These fields were for various possible verb paradigm forms. For example


{\small\color{exlabel}(301)}\par
\noindent
\begin{Verbatim}[fontsize=\footnotesize,breaklines,breakanywhere]
   \1s yoban
\end{Verbatim}

These markers are no longer used. See the topic


{\color{cflink} \hyperref[key:Old_and_Changed_Markers]{Old\_and\_Changed\_Markers}}


\chapter{When MDF Fails to Meet Your Requirements}\label{key:When_MDF_Fails_to_Meet_Your_Requirements}

This topic is still under construction. Terse notes on common problems and answers are below. Contact Toolbox@sil.org for more details or other problems.


\section{Typical problems}

appearance of the printout -- modify the template in Word, check ``add to template'' to keep changes for the next export; check Page Setup under Options for columns, margins, odd-even -- these are not included in the template


labels in the wrong language -- check the set of files mdf\_xxx.cct to see if the language you want is available. if not, modify the mdf\_eng.cct and save as an appropriate name, then select that file by Modifying the Audience options.


labels in the wrong script -- modify the template in Word, check ``add to template'' to keep changes for the next export


label needed after the field, not preceding -- contact Toolbox@sil.org. A more flexible system is coming


punctuation is not what is wanted -- punctuation is very difficult to modify. Contact Toolbox@sil.org. A more flexible system is coming.


some fields not included in the printout -- check list of markers to exclude; check Options settings for various ``include...'' options; check Audience selection; check diglot / triglot selection; check Finderlist selection


sorted order of the printout -- check language encoding assignment of \textbackslash{}lx and \textbackslash{}lc in MDF.typ and also in MDF\_RTF.typ, which MDF printing uses for final sorting. (Do Project, Database types in order to check MDF\_RTF.typ.)


fields are printed in the wrong order -- MDF standard printing reorders all fields according to a specific pattern; a version which doesn't reorder is available. contact Toolbox@sil.org for help on this. (Soon to be made widely available.)


marker NOT to be printed, not available -- no problem, create a new marker and use it for whatever notes or questions you are wanting it for. Also, read the NOTE ON ADDING MARKERS at the end.


marker to be printed but NOT FINAL printing, not available -- create the new marker, check ``include non-MDF field'' under options and be sure it is in the list of fields included (Select Fields to be Excluded button in Options). Note, such fields will be printed at the end of the entry. Also, read the NOTE ON ADDING MARKERS at the end.


marker for final printing not available -- be sure you understand what *is* available (read the descriptions of the markers). MDF standard printing puts all non-MDF fields at the end. If that's OK, be sure it's included for printing (``Select Fields to be Excluded'' button in Options).  Also, read the NOTE ON ADDING MARKERS at the end.


need markers for more scripts than are available, multiple scripts for the same field (eg, lexeme) --  if it's only the lexeme (perhaps entered in Roman but you want to print in script) AND you are not using the lexical citation option (\textbackslash{}lc field), then the \textbackslash{}lc field can be used for the script form of the lexeme. Similarly, if you don't have a regional language, you can use the various regional markers to provide another script (Actually, you may have to shuffle the meanings of the various ``languages'' in order to get the order you want.) -- these tricks are not really encouraged, but can be made to work. If these don't meet your needs, contact Toolbox@sil.org; a more flexible system is being developed. Also, read the NOTE ON ADDING MARKERS at the end.


\begin{quote}\small\itshape NOTE ON ADDING MARKERS:\end{quote}

The MDF marker set is a standard recognized by various organizations. Additional markers should not be added without good motivation. On the other hand, it is also recognized that the marker set does not cover all possible reasonable needs. If you are going to publish electronically or if you are going to archive your dictionary through some organization, be sure they understand and approve the marker set you are using. Also, include your database type file with the archive and be sure your new markers have a good description. If you have legacy data, be sure that the archiving or publishing organization has the appropriate font(s) and language encoding file(s).


\chapter{Field Marker Reference}
\section{Record Marker}
\subsection{\textbackslash{}lx  lexeme or headword of the lexical entry}\label{key:lx}

This is the record marker in Toolbox and is the field by which the database is normally sorted. When formatted, MDF resorts the dictionary based on this field (or the \textbackslash{}lc field, if present).


This field contains the lexeme or headword, which is commonly mono-morphemic in a Toolbox lexical database. But such a lexeme form may not be very accessible for vernacular speakers if printed. To provide a more readable form for vernacular speakers, use the \textbackslash{}lc field.


For a discussion of ``lexical citation'', see:


{\color{cflink} \hyperref[key:lc]{\textbackslash{}lc}}


\begin{quote}\small\itshape Since this is the record marker, it cannot be added inside any record. It is discussed here simply  for completeness.\end{quote}

\section{Basic Fields}
\subsection{\textbackslash{}ce  cross-reference (English gloss)}\label{key:ce}

This gives the English gloss for the vernacular lexeme referenced by the preceding \textbackslash{}cf field.


For more information, see:


\begin{description}
\item[{\hyperref[key:cf]{\textcolor{cflink}{\textbackslash{}cf}}}] confer/cross-reference
\end{description}


Identical national and regional language gloss fields are:


{\color{cflink} \hyperref[key:cn]{\textbackslash{}cn} and \hyperref[key:cr]{\textbackslash{}cr}}


{\small <Basic>}\par

\subsection{\textbackslash{}cf  confer/cross-reference}\label{key:cf}

This is a generic reference marker used to link together any two related entries in the lexicon. For example, in Selaru, `-aswasw kaha' means `high water mark' and needs to cross-reference `manahma' (the entry for `rising tide') and visa-versa.


The \textbackslash{}cf field is bundled with the \textbackslash{}ce, \textbackslash{}cn, and \textbackslash{}cr gloss fields, so this example would be encoded as:


{\small\color{exlabel}(302)}\par
\noindent
\begin{Verbatim}[fontsize=\footnotesize,breaklines,breakanywhere]
 \cf manahma
\end{Verbatim}

{\small\color{exlabel}(303)}\par
\noindent
\begin{Verbatim}[fontsize=\footnotesize,breaklines,breakanywhere]
 \ce rising tide
\end{Verbatim}

{\small\color{exlabel}(304)}\par
\noindent
\begin{Verbatim}[fontsize=\footnotesize,breaklines,breakanywhere]
 \cn air naik
\end{Verbatim}

{\small\color{exlabel}(305)}\par
\noindent
\begin{Verbatim}[fontsize=\footnotesize,breaklines,breakanywhere]
 \cr

\end{Verbatim}

and would typically print in a triglot dictionary as:


{\small\color{exlabel}(306)}\par
\noindent
\begin{Verbatim}[fontsize=\footnotesize,breaklines,breakanywhere]
  See: manahma 'rising tide' 'air naik'.
 
\end{Verbatim}

Multiple cross reference bundles are usually concatenated with semicolons `;' as seen in:


{\small\color{exlabel}(307)}\par
\noindent
\begin{Verbatim}[fontsize=\footnotesize,breaklines,breakanywhere]
 \cf manahma
\end{Verbatim}

{\small\color{exlabel}(308)}\par
\noindent
\begin{Verbatim}[fontsize=\footnotesize,breaklines,breakanywhere]
 \ce rising tide
\end{Verbatim}

{\small\color{exlabel}(309)}\par
\noindent
\begin{Verbatim}[fontsize=\footnotesize,breaklines,breakanywhere]
 \cn air naik
\end{Verbatim}

{\small\color{exlabel}(310)}\par
\noindent
\begin{Verbatim}[fontsize=\footnotesize,breaklines,breakanywhere]
 \cr
\end{Verbatim}

{\small\color{exlabel}(311)}\par
\noindent
\begin{Verbatim}[fontsize=\footnotesize,breaklines,breakanywhere]
 \cf mety2
\end{Verbatim}

{\small\color{exlabel}(312)}\par
\noindent
\begin{Verbatim}[fontsize=\footnotesize,breaklines,breakanywhere]
 \ce low tide
\end{Verbatim}

{\small\color{exlabel}(313)}\par
\noindent
\begin{Verbatim}[fontsize=\footnotesize,breaklines,breakanywhere]
 \cn
\end{Verbatim}

{\small\color{exlabel}(314)}\par
\noindent
\begin{Verbatim}[fontsize=\footnotesize,breaklines,breakanywhere]
 \cr
\end{Verbatim}

{\small\color{exlabel}(315)}\par
\noindent
\begin{Verbatim}[fontsize=\footnotesize,breaklines,breakanywhere]
 \cf rean
\end{Verbatim}

{\small\color{exlabel}(316)}\par
\noindent
\begin{Verbatim}[fontsize=\footnotesize,breaklines,breakanywhere]
 \ce high tide
\end{Verbatim}

{\small\color{exlabel}(317)}\par
\noindent
\begin{Verbatim}[fontsize=\footnotesize,breaklines,breakanywhere]
 \cn
\end{Verbatim}

{\small\color{exlabel}(318)}\par
\noindent
\begin{Verbatim}[fontsize=\footnotesize,breaklines,breakanywhere]
 \cr

\end{Verbatim}

This would typically print in a diglot dictionary as:


{\small\color{exlabel}(319)}\par
\noindent
\begin{Verbatim}[fontsize=\footnotesize,breaklines,breakanywhere]
  See: manahma 'rising tide'; mety2 'low tide'; 
  rean 'high tide'.

\end{Verbatim}

(Note that actually the vernacular text is usually be printed as bold, and the `2' for the homonym number of `mety' as subscripted, but these can't be displayed that way in Toolbox.)


One major short coming of using the cross-reference fields (rather than the lexical function fields) is that the semantic relationship between the two related lexemes is not made explicit. But it is often the case that a researcher will know there is some kind of relationship between two lexemes but is unclear as to the nature of that relationship. This is what the \textbackslash{}cf field is good for. And once the relationship is determined, the cross-reference information could then be transferred to an \textbackslash{}lf field bundle.


It is not uncommon for there to be many lexemes that are interrelated. This creates a myriad of cross-references, where each lexeme of a given type is cross-referenced to all other lexemes of that type. Adding new entries to the group can also be very tedious. A simple solution to this is to choose one lexeme as the focal point and have all other related lexemes refer only to that one. That focal lexeme will then contain a listing of all other lexemes in the group, by using either a series of \textbackslash{}lf (lexical function) bundles, to describe the relationships between that focal lexeme and each of the other lexemes, or by using the \textbackslash{}tb (table) field to make an actual list of the other lexemes and their meanings.


For more on grouping cross-references or mapping related lexemes, see:


\begin{description}
\item[{\hyperref[key:lf]{\textcolor{cflink}{\textbackslash{}lf}}}] lexical function label
\item[{\hyperref[key:tb]{\textcolor{cflink}{\textbackslash{}tb}}}] table
\end{description}


For the English, national and regional glossing fields, see:


{\color{cflink} \hyperref[key:ce]{\textbackslash{}ce,} \hyperref[key:cn]{\textbackslash{}cn} and \hyperref[key:cr]{\textbackslash{}cr}}


{\small <Basic>}\par

\subsection{\textbackslash{}cn  cross-reference (national gloss)}\label{key:cn}

This gives the national language gloss for the vernacular lexeme referenced by the preceding \textbackslash{}cf field.


For more information, see:


\begin{description}
\item[{\hyperref[key:cf]{\textcolor{cflink}{\textbackslash{}cf}}}] confer/cross-reference
\end{description}


Identical English and regional language gloss fields are:


{\color{cflink} \hyperref[key:ce]{\textbackslash{}ce} and \hyperref[key:cr]{\textbackslash{}cr}}


{\small <Basic>}\par

\subsection{\textbackslash{}de  definition (English)}\label{key:de}

This field is used to fully express the semantic domains of each sense of a lexeme in English. For related information, see:


\begin{description}
\item[{\hyperref[key:ge]{\textcolor{cflink}{\textbackslash{}ge}}}] gloss (English)
\end{description}


\subsubsection{Hint}

If an entry has many English glosses, none of which are vitally important to differentiate while glossing texts, give only one gloss in the \textbackslash{}ge field (to simplify interlinearizing) and then give all of the glosses in the \textbackslash{}de and \textbackslash{}re fields. When the dictionary is formatted, the \textbackslash{}de field will be used and the \textbackslash{}ge field will be discarded. Also, because there are \textbackslash{}re fields, the \textbackslash{}ge field will be ignored for reversing the database.


For example, the Indonesian morpheme `-nya' means `his' `hers' or `its', but for glossing `3sPOS' may be adequate (if not preferable), but `3sPOS' is awkward for both a dictionary printout and an entry in a reversed English finderlist. So, do something like this:


{\small\color{exlabel}(320)}\par
\noindent
\begin{Verbatim}[fontsize=\footnotesize,breaklines,breakanywhere]
  \ge 3sPOS
\end{Verbatim}

{\small\color{exlabel}(321)}\par
\noindent
\begin{Verbatim}[fontsize=\footnotesize,breaklines,breakanywhere]
  \re his ; hers ; its
\end{Verbatim}

{\small\color{exlabel}(322)}\par
\noindent
\begin{Verbatim}[fontsize=\footnotesize,breaklines,breakanywhere]
  \de his, hers, its

\end{Verbatim}

This will simplify interlinearizing by keeping Toolbox from asking which gloss (`his', `hers', or `its') is appropriate for each occurrence of `-nya' in a text, and yet will allow your lexical database to be more complete in its description of the meaning or semantic domain of the lexeme.


Identical vernacular, national and regional fields are:


{\color{cflink} \hyperref[key:dv]{\textbackslash{}dv,} \hyperref[key:dn]{\textbackslash{}dn,} and \hyperref[key:dr]{\textbackslash{}dr}}


Closely related fields are:


{\color{cflink} \hyperref[key:ue]{\textbackslash{}ue,} \hyperref[key:oe]{\textbackslash{}oe,} and \hyperref[key:ee]{\textbackslash{}ee}}


{\small <Basic>}\par

\subsection{\textbackslash{}dn  definition (national)}\label{key:dn}

This field is used to fully express the semantic domains of each sense of a lexeme in the national language.


For more information, see ``Hint'' under:


\begin{description}
\item[{\hyperref[key:de]{\textcolor{cflink}{\textbackslash{}de}}}] definition English
\end{description}


See also:


\begin{description}
\item[{\hyperref[key:gn]{\textcolor{cflink}{\textbackslash{}gn}}}] gloss national
\end{description}


For the related definition fields, see:


{\color{cflink} \hyperref[key:dv]{\textbackslash{}dv,} \hyperref[key:de]{\textbackslash{}de,} and \hyperref[key:dr]{\textbackslash{}dr}}


Closely related national language fields are:


{\color{cflink} \hyperref[key:un]{\textbackslash{}un,} \hyperref[key:on]{\textbackslash{}on,} and \hyperref[key:en]{\textbackslash{}en}}


{\small <Basic>}\par

\subsection{\textbackslash{}dt  datestamp}\label{key:dt}

A Toolbox field to help you keep track of the last time you edited an entry. There need only be one of these in a record (usually the last field) and is usually inserted automatically by Toolbox. The field is set up under the Toolbox menu option: Database-Properties-Options tab. The field marker must be defined first before it can be selected as the datestamp field.


\begin{quote}\small\itshape This field does not normally print.\end{quote}

{\small <Basic>}\par

\subsection{\textbackslash{}ge  gloss (English)}\label{key:ge}

This is intended for interlinear morpheme-level glossing. It is used for reversing the dictionary if an \textbackslash{}re field is not present (or is present but empty). It is also used as an English definition in the printed dictionary if there is no \textbackslash{}de field (or it is present but empty).


The user may enter data either with each gloss getting its own field:


{\small\color{exlabel}(323)}\par
\noindent
\begin{Verbatim}[fontsize=\footnotesize,breaklines,breakanywhere]
  \ge put_out
\end{Verbatim}

{\small\color{exlabel}(324)}\par
\noindent
\begin{Verbatim}[fontsize=\footnotesize,breaklines,breakanywhere]
  \ge move

\end{Verbatim}

The underline character `\_' in the example glosses above is to force Toolbox to treat the multiple word gloss `put out' as a single gloss. When printing a dictionary, the underlines normally will be converted to spaces.


Similarly, a dot (``period'' or ``full stop'') can be placed between the words instead of the underline:


{\small\color{exlabel}(325)}\par
\noindent
\begin{Verbatim}[fontsize=\footnotesize,breaklines,breakanywhere]
  \ge put.out
\end{Verbatim}

{\small\color{exlabel}(326)}\par
\noindent
\begin{Verbatim}[fontsize=\footnotesize,breaklines,breakanywhere]
  \ge move


\end{Verbatim}

An oldere style allows you to list the glosses strung together, separated by a semicolon with a space on either side of it (so as to be unlike any normal punctuation):


{\small\color{exlabel}(327)}\par
\noindent
\begin{Verbatim}[fontsize=\footnotesize,breaklines,breakanywhere]
  \ge put_out ; move

\end{Verbatim}

Toolbox can recognize either format and when interlinearizing will give the user the choice of both glosses in either case. The sequence ` ; ` is also converted to `, ` by MDF printing when formatting a dictionary or finderlist.


The advantage of keeping your glosses in separate fields is that if you setup Toolbox to sort on the \textbackslash{}ge fields (to get a type of finderlist), it will sort on all of the glosses; but if the glosses are all concatenated, Toolbox will only sort on the first form.


For important related information, see:


\begin{description}
\item[{\hyperref[key:de]{\textcolor{cflink}{\textbackslash{}de}}}] definition (English)
\item[{\hyperref[key:re]{\textcolor{cflink}{\textbackslash{}re}}}] reverse form (English)
\end{description}


{\small <Basic>}\par

\subsection{\textbackslash{}gn  gloss (national)}\label{key:gn}

This is intended for interlinear morpheme-level glossing. It is used for reversing the dictionary if an \textbackslash{}rn field is not present (or is present but empty). It is also used as a national language definition in the printed dictionary if there is no \textbackslash{}dn field (or it is present but empty).


The user may enter data either with each gloss getting its own field:


{\small\color{exlabel}(328)}\par
\noindent
\begin{Verbatim}[fontsize=\footnotesize,breaklines,breakanywhere]
  \gn keluarkan
\end{Verbatim}

{\small\color{exlabel}(329)}\par
\noindent
\begin{Verbatim}[fontsize=\footnotesize,breaklines,breakanywhere]
  \gn pindahkan

\end{Verbatim}

Or with the glosses strung together, separated by a semicolon with a space on either side of it (so as to be unlike any normal punctuation):


{\small\color{exlabel}(330)}\par
\noindent
\begin{Verbatim}[fontsize=\footnotesize,breaklines,breakanywhere]
  \gn keluarkan ; pindahkan

\end{Verbatim}

Toolbox can recognize either format and when interlinearizing will give the user the choice of both glosses in either case.


MDF will convert the ` ; ` sequence to `, ` when formatting the dictionary. (MDF will also automatically convert any underline character `\_' in a gloss field to a space when formatting.)


For important related information, see:


\begin{description}
\item[{\hyperref[key:dn]{\textcolor{cflink}{\textbackslash{}dn}}}] definition (national)
\item[{\hyperref[key:rn]{\textcolor{cflink}{\textbackslash{}rn}}}] reversal form (national)
\end{description}


{\small <Basic>}\par

\subsection{\textbackslash{}nt  notes, etc.}\label{key:nt}

This is a generic dump for all your personal notes about an entry, subentry, or sense. More specific note fields are provided for ``splitters'', i.e. those who desire a finer differentiation to their notes. These are:


\begin{description}
\item[{\hyperref[key:np]{\textcolor{cflink}{\textbackslash{}np}}}] notes
\item[{\hyperref[key:on]{\textcolor{cflink}{on}}}] phonology
\item[{\hyperref[key:ng]{\textcolor{cflink}{\textbackslash{}ng}}}] notes
\item[{\hyperref[key:on]{\textcolor{cflink}{on}}}] grammar
\item[{\hyperref[key:nd]{\textcolor{cflink}{\textbackslash{}nd}}}] notes
\item[{\hyperref[key:on]{\textcolor{cflink}{on}}}] discourse
\item[{\hyperref[key:na]{\textcolor{cflink}{\textbackslash{}na}}}] notes
\item[{\hyperref[key:on]{\textcolor{cflink}{on}}}] anthropology
\item[{\hyperref[key:ns]{\textcolor{cflink}{\textbackslash{}ns}}}] notes
\item[{\hyperref[key:on]{\textcolor{cflink}{on}}}] sociolinguistics
\item[{\hyperref[key:nq]{\textcolor{cflink}{\textbackslash{}nq}}}] questions
\end{description}


If selected for output, MDF adds the label ``Notes: `` to this field and brackets it with ``[ ]''.


\begin{quote}\small\itshape This is a ``free-form'' field.   Punctuation and capitalization should be used as needed.\end{quote}

{\small <Basic>}\par

\subsection{\textbackslash{}pn  part of speech (national)}\label{key:pn}

This field is used to classify the part of speech with labels found in national language dictionaries. Consistent labeling is important. Use Toolbox's Range Set feature for this field.


MDF requires that the \textbackslash{}pn field follow the \textbackslash{}ps field:


{\small\color{exlabel}(331)}\par
\noindent
\begin{Verbatim}[fontsize=\footnotesize,breaklines,breakanywhere]
 \ps n   (noun)
\end{Verbatim}

{\small\color{exlabel}(332)}\par
\noindent
\begin{Verbatim}[fontsize=\footnotesize,breaklines,breakanywhere]
 \pn kb  (the national abbreviate for 'noun')

\end{Verbatim}

\begin{quote}\small\itshape If the  order is  reversed, MDF  will not  function properly.\end{quote}

MDF will format the \textbackslash{}pn field only if you specify that the output is for a national audience. When a national audience is specified, the \textbackslash{}pn field will replace the \textbackslash{}ps field. But if there is no \textbackslash{}pn field or it is empty, the \textbackslash{}ps field will be output for the national audience as for an English audience.


For more information, see:


\begin{description}
\item[{\hyperref[key:ps]{\textcolor{cflink}{\textbackslash{}ps}}}] part of speech
\end{description}


For important information on Range Sets and on MDF's label substitution feature, see:


{\color{cflink} \hyperref[key:Range_Sets]{Range\_Sets}}


{\small <Basic>}\par

\subsection{\textbackslash{}ps  part of speech}\label{key:ps}

This field is used to classify the part of speech for the vernacular form (not the national or English gloss), i.e., ``fat'' may be an adjective in English, but that does not mean the vernacular form can be classified as such. Consistent labeling is important. Use Toolbox's Range Set feature for this field. For more information, see:


{\color{cflink} \hyperref[key:Range_Sets]{Range\_Sets}}


To specify a national part of speech label, see:


\begin{description}
\item[{\hyperref[key:pn]{\textcolor{cflink}{\textbackslash{}pn}}}] part of speech (national)
\end{description}


{\small <Basic>}\par

\subsection{\textbackslash{}re  reversal form (English)}\label{key:re}

This gives the English word or phrase to be use to reverse the dictionary for an English index. If an \textbackslash{}re * is present, the relevant entry, subentry, or sense will be ignored (i.e. not included in the reversed index). Like the \textbackslash{}ge field, the data for this field can be kept each in its own field or concatenated in one field, separated by ` ; `.


It is often the case that there are several translation equivalents for a single vernacular term. `huma' might mean `house', `hut', `shack', `dwelling', `lean-to', etc. Each of these equivalents would be good to have in the reversed index, e.g.:


{\small\color{exlabel}(333)}\par
\noindent
\begin{Verbatim}[fontsize=\footnotesize,breaklines,breakanywhere]
 \re house ; hut ; shack ; dwelling ; lean-to

\end{Verbatim}

They could also be entered in separate \textbackslash{}re fields. For more information on this, see:


\begin{description}
\item[{\hyperref[key:ge]{\textcolor{cflink}{\textbackslash{}ge}}}] gloss (English)
\end{description}


The advantage of keeping them in separate fields is that if you setup Toolbox to sort on the \textbackslash{}re fields (to get a type of finderlist), it will sort on all of them, but if the reversal glosses are all concatenated, Toolbox will only index on the first form.


\begin{quote}\small\itshape This field does not normally print in the dictionary.\end{quote}

{\small <Basic>}\par

\subsection{\textbackslash{}rf  reference to notebooks, texts, etc.}\label{key:rf}

This field keeps the notebook reference for the following example sentence. This will enable you to validate the example at a later date (e.g., if it comes from an early notebook, it may be suspect).


The \textbackslash{}rf, \textbackslash{}xv, \textbackslash{}xn, and \textbackslash{}xe fields are considered ``bundled'' (or grouped) together. If you include multiple examples for a single entry, subentry, or sense, be sure to include the fields grouped together.


{\small\color{exlabel}(334)}\par
\noindent
\begin{Verbatim}[fontsize=\footnotesize,breaklines,breakanywhere]
 \rf nbk1.23.12
\end{Verbatim}

{\small\color{exlabel}(335)}\par
\noindent
\begin{Verbatim}[fontsize=\footnotesize,breaklines,breakanywhere]
 \xv example1
\end{Verbatim}

{\small\color{exlabel}(336)}\par
\noindent
\begin{Verbatim}[fontsize=\footnotesize,breaklines,breakanywhere]
 \xn ...
\end{Verbatim}

{\small\color{exlabel}(337)}\par
\noindent
\begin{Verbatim}[fontsize=\footnotesize,breaklines,breakanywhere]
 \xe ...

\end{Verbatim}

{\small\color{exlabel}(338)}\par
\noindent
\begin{Verbatim}[fontsize=\footnotesize,breaklines,breakanywhere]
 \rf nbk1.54.09
\end{Verbatim}

{\small\color{exlabel}(339)}\par
\noindent
\begin{Verbatim}[fontsize=\footnotesize,breaklines,breakanywhere]
 \xv example2
\end{Verbatim}

{\small\color{exlabel}(340)}\par
\noindent
\begin{Verbatim}[fontsize=\footnotesize,breaklines,breakanywhere]
 \xn ...
\end{Verbatim}

{\small\color{exlabel}(341)}\par
\noindent
\begin{Verbatim}[fontsize=\footnotesize,breaklines,breakanywhere]
 \xe ...   etc.

\end{Verbatim}

Usually these are grouped with the \textbackslash{}rf field at the beginning of each bundle, but if you don't want to use the \textbackslash{}rf field, then the \textbackslash{}xv field will mark the beginning of example sentence bundles.


You may use as many different example sentence bundles as you need for each sense, part of speech, and/or subentry in a record. Within a given section (e.g. sense), multiple examples are printed one after the other.


MDF adds the label ``Ref: `` to this field.


See also:


\begin{description}
\item[{\hyperref[key:xv]{\textcolor{cflink}{\textbackslash{}xv}}}] example sentence (vernacular)
\end{description}


\begin{quote}\small\itshape This is a ``free-form'' field.   Punctuation and capitalization should be used as needed.\end{quote}

{\small <Basic>}\par

\subsection{\textbackslash{}rn  reversal form (national)}\label{key:rn}

This gives the national language word or phrase form to be use to reverse the dictionary for a national language index. If an \textbackslash{}rn * is present the relevant entry, subentry, or sense will be ignored (i.e. not included in the reversed index). Like the \textbackslash{}gn field, the data for this field can be kept each in its own field or concatenated in one field, separated by ` ; `.


For more information, see:


\begin{description}
\item[{\hyperref[key:re]{\textcolor{cflink}{\textbackslash{}re}}}] reverse English
\end{description}


Also see:


\begin{description}
\item[{\hyperref[key:gn]{\textcolor{cflink}{\textbackslash{}gn}}}] gloss-national
\end{description}


\begin{quote}\small\itshape This field does not normally print in the dictionary.\end{quote}

{\small <Basic>}\par

\subsection{\textbackslash{}xe  translation of example (English)}\label{key:xe}

This provides the English translation of the example sentence given in the \textbackslash{}xv field.


For more information about example sentences, see:


\begin{description}
\item[{\hyperref[key:xv]{\textcolor{cflink}{\textbackslash{}xv}}}] example sentence (vernacular)
\item[{\hyperref[key:rf]{\textcolor{cflink}{\textbackslash{}rf}}}] reference to notebooks, texts, etc.
\end{description}


Related fields are:


{\color{cflink} \hyperref[key:xn]{\textbackslash{}xn} and \hyperref[key:xr]{\textbackslash{}xr}}


\begin{quote}\small\itshape This is a ``free-form'' field.   Punctuation and capitalization should be used as needed.\end{quote}

{\small <Basic>}\par

\subsection{\textbackslash{}xn  translation of example (national)}\label{key:xn}

This provides the national translation of the example sentence given in the \textbackslash{}xv field.


For more information about example sentences, see:


\begin{description}
\item[{\hyperref[key:xv]{\textcolor{cflink}{\textbackslash{}xv}}}] example sentence (vernacular)
\item[{\hyperref[key:rf]{\textcolor{cflink}{\textbackslash{}rf}}}] reference to notebooks, texts, etc.
\end{description}


Related fields are:


{\color{cflink} \hyperref[key:xe]{\textbackslash{}xe} and \hyperref[key:xr]{\textbackslash{}xr}}


\begin{quote}\small\itshape This is a ``free-form'' field.   Punctuation and capitalization should be used as needed.\end{quote}

{\small <Basic>}\par

\subsection{\textbackslash{}xv  example sentence (vernacular)}\label{key:xv}

A good rule of thumb is to keep your example sentences relatively short (one line or so). Good examples should indicate the usage of the lexeme without being stilted.


Be aware that sentences taken straight from texts generally do not make good example sentences, because this removes them from their context. Such sentences are often encoded with particles operating on the larger discourse and are likely to be bound anaphorically with preceding participants.


Nouns rarely need example sentences, but good sentences are crucial to differentiating the various senses of verbs and can demonstrate and verify peculiar or rare domains or usages.


For more information about example sentences, see:


\begin{description}
\item[{\hyperref[key:rf]{\textcolor{cflink}{\textbackslash{}rf}}}] reference to notebooks, texts, etc.
\end{description}


Related fields, used to translate or gloss the vernacular example sentence, are:


{\color{cflink} \hyperref[key:xe]{\textbackslash{}xe,} \hyperref[key:xn]{\textbackslash{}xn,} and \hyperref[key:xr]{\textbackslash{}xr}}


\begin{quote}\small\itshape This is a ``free-form'' field.   Punctuation and capitalization should be used as needed.\end{quote}

{\small <Basic>}\par

\section{Reserved Fields}
\subsection{\textbackslash{}hm  homonym number}\label{key:hm}

This is field is used to differentiate homonym entries (lexemes that sound or are spelled the same but have no semantic relationship). This field generally comes directly after the \textbackslash{}lx field and is simply followed by a 1, 2, or 3, etc.:


{\small\color{exlabel}(342)}\par
\noindent
\begin{Verbatim}[fontsize=\footnotesize,breaklines,breakanywhere]
  \lx asw
\end{Verbatim}

{\small\color{exlabel}(343)}\par
\noindent
\begin{Verbatim}[fontsize=\footnotesize,breaklines,breakanywhere]
  \hm 2
\end{Verbatim}

{\small\color{exlabel}(344)}\par
\noindent
\begin{Verbatim}[fontsize=\footnotesize,breaklines,breakanywhere]
  \ps n
  ...
\end{Verbatim}

(note the lack of punctuation in the \textbackslash{}hm field)


Any cross-reference to one of these entries should also include the homonym number, e.g.:


{\small\color{exlabel}(345)}\par
\noindent
\begin{Verbatim}[fontsize=\footnotesize,breaklines,breakanywhere]
  \cf asw2

\end{Verbatim}

When the lexical database is converted to MS-Word format for printing, the homonym number for both the entry and the cross-reference will be subscripted by MDF using the homonym number character style.


For more on character styles, see:


{\color{cflink} \hyperref[key:Character_Style_Codes]{Character\_Style\_Codes}}


{\small <Reserved>}\par

\subsection{\textbackslash{}lc  lexical citation}\label{key:lc}

This should be added only if the lexical entry form is inappropriate for the printed dictionary, and you want to substitute another form for the printed entry, e.g. you might want the entry ``lewat'' printed as ``léwat'' if its stress pattern were not predictable.


Roots are commonly used as the lexeme form in Toolbox lexical databases, but if the language affixes prepositionally, a dictionary printed with root forms as the entry form can be very confusing to native speakers. In this case, you could choose a consistent conjugated form (e.g. 3s-verb form) for a citation form.


MDF will always replace the \textbackslash{}lx field with the \textbackslash{}lc field, if present, and then resort the dictionary according to these fields. (You can choose to restrict MDF to sort the database only by the \textbackslash{}lx field, even for those entries with an \textbackslash{}lc field. This restriction will not effect the \textbackslash{}lc field though, an \textbackslash{}lc field will still print as the entry form for its record. Choosing to restrict MDF to the \textbackslash{}lx field may cause some citation forms to appear out of sequence.)


{\small <Reserved>}\par

\subsection{\textbackslash{}se  subentry (a polymorphemic form or a phrase)}\label{key:se}

This is like the \textbackslash{}lx field except it occurs within the record, marking the word (or phrase) as a form derived from the root. Following this marker would be all the fields that comprise a typical lexical entry. There can be several of these subentries within a record. Subentries can also have multiple senses. A simple example using subentries would be:


{\small\color{exlabel}(346)}\par
\noindent
\begin{Verbatim}[fontsize=\footnotesize,breaklines,breakanywhere]
 \lx bren
\end{Verbatim}

{\small\color{exlabel}(347)}\par
\noindent
\begin{Verbatim}[fontsize=\footnotesize,breaklines,breakanywhere]
 \ps vi
\end{Verbatim}

{\small\color{exlabel}(348)}\par
\noindent
\begin{Verbatim}[fontsize=\footnotesize,breaklines,breakanywhere]
 \ge play
\end{Verbatim}

{\small\color{exlabel}(349)}\par
\noindent
\begin{Verbatim}[fontsize=\footnotesize,breaklines,breakanywhere]
 \ee Implies lack of focus or purpose.

\end{Verbatim}

{\small\color{exlabel}(350)}\par
\noindent
\begin{Verbatim}[fontsize=\footnotesize,breaklines,breakanywhere]
 \se brenak
\end{Verbatim}

{\small\color{exlabel}(351)}\par
\noindent
\begin{Verbatim}[fontsize=\footnotesize,breaklines,breakanywhere]
 \ps vt
\end{Verbatim}

{\small\color{exlabel}(352)}\par
\noindent
\begin{Verbatim}[fontsize=\footnotesize,breaklines,breakanywhere]
 \ge play_s.t.
\end{Verbatim}

{\small\color{exlabel}(353)}\par
\noindent
\begin{Verbatim}[fontsize=\footnotesize,breaklines,breakanywhere]
 \de play a game, or play with s.t.

\end{Verbatim}

{\small\color{exlabel}(354)}\par
\noindent
\begin{Verbatim}[fontsize=\footnotesize,breaklines,breakanywhere]
 \se inabren
\end{Verbatim}

{\small\color{exlabel}(355)}\par
\noindent
\begin{Verbatim}[fontsize=\footnotesize,breaklines,breakanywhere]
 \ps n
\end{Verbatim}

{\small\color{exlabel}(356)}\par
\noindent
\begin{Verbatim}[fontsize=\footnotesize,breaklines,breakanywhere]
 \ge recreation ; entertainment

\end{Verbatim}

{\small\color{exlabel}(357)}\par
\noindent
\begin{Verbatim}[fontsize=\footnotesize,breaklines,breakanywhere]
 \se rabrenak
\end{Verbatim}

{\small\color{exlabel}(358)}\par
\noindent
\begin{Verbatim}[fontsize=\footnotesize,breaklines,breakanywhere]
 \ps n
\end{Verbatim}

{\small\color{exlabel}(359)}\par
\noindent
\begin{Verbatim}[fontsize=\footnotesize,breaklines,breakanywhere]
 \ge toy
\end{Verbatim}

{\small\color{exlabel}(360)}\par
\noindent
\begin{Verbatim}[fontsize=\footnotesize,breaklines,breakanywhere]
 \dt 17/Jun/92

\end{Verbatim}

This would typically print like the following:


{\small\color{exlabel}(361)}\par
\noindent
\begin{Verbatim}[fontsize=\footnotesize,breaklines,breakanywhere]
bren     vi. play. Implies lack of focus or purpose.
     brenak    vt. play a game, or play with s.t.
     inabren    n. recreation, entertainment.
     rabrenak n. toy.

\end{Verbatim}

But note that the subentries in this example are far too simplistic; they lack much of the information that should be provided for these polymorphemic lexemes (definitions, example sentences, cross-references, notes, etc.).


For a more detailed discussion of the issues involved in using subentries, see:


{\color{cflink} \hyperref[key:Using_Subentries_or_Lexical_Entries]{Using\_Subentries\_or\_Lexical\_Entries}}


For information on the topic of sections in a lexical entry, see:


{\color{cflink} \hyperref[key:Sections_in_a_Lexical_Entry]{Sections\_in\_a\_Lexical\_Entry}}


{\small <Reserved>}\par

\subsection{\textbackslash{}sn  sense number}\label{key:sn}

Where an entry has more than one sense, this code gives the number and marks the beginning of each sense, e.g.:


{\small\color{exlabel}(362)}\par
\noindent
\begin{Verbatim}[fontsize=\footnotesize,breaklines,breakanywhere]
 \sn 3             (no punctuation)

\end{Verbatim}

Generally each \textbackslash{}sn section contains a full basic set of field markers (especially example sentences as these help to validate the distinctions between the senses). Depending on the hierarchy you choose, the sense number \textbackslash{}sn is considered either below the \textbackslash{}ps field (in the standard hierarchy) or is superior to all fields but the   \textbackslash{}lx, \textbackslash{}lc, and \textbackslash{}hm fields (in the alternate hierarchy).


For more information on hierarchy, see:


{\color{cflink} \hyperref[key:Sections_in_a_Lexical_Entry]{Sections\_in\_a\_Lexical\_Entry}}


{\color{cflink} \hyperref[key:Alternate_Hierarchy]{Alternate\_Hierarchy}}


For information relating to examples, see:


{\color{cflink} \hyperref[key:rf]{\textbackslash{}rf}}


\begin{quote}\small\itshape Do not forget to include \textbackslash{}sn 1 in records that have more than one sense.\end{quote}

{\small <Reserved>}\par

\section{Optional Fields}
\subsection{\textbackslash{}an  antonym}\label{key:an}

This field is used to reference an antonym of the lexeme.


MDF's standard printing adds the label ``Ant: `` to this field.


Using the \textbackslash{}lf (lexical function) field for this is a better practice. For example,


{\small\color{exlabel}(363)}\par
\noindent
\begin{Verbatim}[fontsize=\footnotesize,breaklines,breakanywhere]
  \an sal

\end{Verbatim}

could be done instead as


{\small\color{exlabel}(364)}\par
\noindent
\begin{Verbatim}[fontsize=\footnotesize,breaklines,breakanywhere]
  \lf Ant.
  \lv wrong, false

\end{Verbatim}

The latter gives more information to the outside reader by providing a gloss. (National and regional glosses are available also.)


For more detailed information on this, see:


\begin{description}
\item[{\hyperref[key:sy]{\textcolor{cflink}{\textbackslash{}sy}}}] synonym
\item[{\hyperref[key:lf]{\textcolor{cflink}{\textbackslash{}lf}}}] lexical function
\end{description}


{\small <Optional>}\par

\subsection{\textbackslash{}bb  bibliographic reference}\label{key:bb}

This is used to record any bibliographic information pertinent to the lexeme. MDF adds the label `Read:' to this field.


{\small <Optional>}\par

\subsection{\textbackslash{}bw  borrowed word}\label{key:bw}

This is for denoting a borrowed word or the source language. There is no standard way such information should be encoded, but generally the following is most common:


{\small\color{exlabel}(365)}\par
\noindent
\begin{Verbatim}[fontsize=\footnotesize,breaklines,breakanywhere]
 \bw Arabic

\end{Verbatim}

Which would typically print as:


{\small\color{exlabel}(366)}\par
\noindent
\begin{Verbatim}[fontsize=\footnotesize,breaklines,breakanywhere]
  From: Arabic.

\end{Verbatim}

{\small <Optional>}\par

\subsection{\textbackslash{}cr  cross-reference (regional gloss)}\label{key:cr}

This gives the regional language gloss for the vernacular lexeme referenced by the preceding \textbackslash{}cf field. MDF formats the \textbackslash{}cr field with quotes inside square brackets.


For more information, see:


\begin{description}
\item[{\hyperref[key:cf]{\textcolor{cflink}{\textbackslash{}cf}}}] confer/cross-reference
\end{description}


Identical English and national language gloss fields are:


{\color{cflink} \hyperref[key:ce]{\textbackslash{}ce} and \hyperref[key:cn]{\textbackslash{}cn}}


{\small <Optional>}\par

\subsection{\textbackslash{}dr  definition (regional)}\label{key:dr}

Often early in the project the definitions we receive are actually in the regional language, not the national language; such definitions could be stored here.


If this field is included in the output, MDF adds a label and brackets ``[ ]'' around it and it is treated as part of the national field.


For more information, see ``Hint'' under:


\begin{description}
\item[{\hyperref[key:de]{\textcolor{cflink}{\textbackslash{}de}}}] definition (English)
\end{description}


See also:


\begin{description}
\item[{\hyperref[key:gr]{\textcolor{cflink}{\textbackslash{}gr}}}] gloss (regional)
\end{description}


For the related definition fields, see:


{\color{cflink} \hyperref[key:dv]{\textbackslash{}dv,} \hyperref[key:de]{\textbackslash{}de,} and \hyperref[key:dn]{\textbackslash{}dn}}


Closely related regional language fields are:


{\color{cflink} \hyperref[key:ur]{\textbackslash{}ur,} \hyperref[key:or]{\textbackslash{}or,} and \hyperref[key:er]{\textbackslash{}er}}


{\small <Optional>}\par

\subsection{\textbackslash{}dv  definition (vernacular)}\label{key:dv}

This is the hardest of all fields to fill in, because it requires the researcher to explain, in the vernacular, what the salient concepts are (i.e. the domain) that this ``unit of meaning'' captures -- from a native speaker's perspective. Obviously not easy and something usually left to much later.


For the related definition fields, see:


{\color{cflink} \hyperref[key:de]{\textbackslash{}de,} \hyperref[key:dn]{\textbackslash{}dn,} and \hyperref[key:dr]{\textbackslash{}dr}}


This field is also used for creating a monolingual dictionary. For related monolingual fields, see:


{\color{cflink} \hyperref[key:gv]{\textbackslash{}gv,} \hyperref[key:uv]{\textbackslash{}uv,} \hyperref[key:ov]{\textbackslash{}ov,} and \hyperref[key:ev]{\textbackslash{}ev}}


{\small <Optional>}\par

\subsection{\textbackslash{}ec  etymology-comment}\label{key:ec}

Any comments the researcher needs to add concerning the etymology of the lexeme can be given here.


For more information, see:


{\color{cflink} \hyperref[key:et]{\textbackslash{}et}}


\begin{quote}\small\itshape This field does not normally print.\end{quote}

\begin{quote}\small\itshape This is a ``free-form'' field.   Punctuation and capitalization should be used as needed.\end{quote}

{\small <Optional>}\par

\subsection{\textbackslash{}ee  encyclopedic information (English)}\label{key:ee}

A field researcher has an incredible opportunity to assimilate intuitive knowledge of the language and culture of the people with whom he/she works. This type of information is invaluable to others who have no access to the language area (or at least no extended contact with the language community).


This knowledge is often never codified. Now Toolbox provides a simple means by which a researcher can ``put on paper'' this kind of information. Once added, the lexical database becomes more than a ``dictionary'' but a ``knowledge-base'' of the language.


This field crosses over with the \textbackslash{}de, \textbackslash{}ue, and \textbackslash{}oe fields, but is intended for more verbose explanations of the lexeme (headword, subentry or sense). Basically, the researcher should use this field to encode any additional information needed by a non-native speaker to understand and use this lexeme properly. For example, the lexeme ``hatw'' (in Selaru) can be described with the following fields:


{\small\color{exlabel}(367)}\par
\noindent
\begin{Verbatim}[fontsize=\footnotesize,breaklines,breakanywhere]
  \de a fired, earthen, baking form used for making fv:skyerker
\end{Verbatim}

{\small\color{exlabel}(368)}\par
\noindent
\begin{Verbatim}[fontsize=\footnotesize,breaklines,breakanywhere]
  \ee Often today this is replaced by a cast iron griddle, much 
      like a waffle iron, which is frequently found in the shape 
      of four hearts.

\end{Verbatim}

MDF does not add any label to this field, but simply formats it as entered.


\subsubsection{Tip}

Since MDF does not format or add any label to this field, you can actually use this field to include additional information about a lexical entry which is not handled by any of the other fields that MDF supports. You can even add your own label to this information. It's not pretty, but it works. For example, if you wanted to create a new field to keep track of the weather (a hypothetical example to be sure), you could enter your observations thus:


{\small\color{exlabel}(369)}\par
\noindent
\begin{Verbatim}[fontsize=\footnotesize,breaklines,breakanywhere]
  \ee |fl{Weather: }The sun is hot with scattered showers expected.

\end{Verbatim}

The special character code `` |fl\{ \}'' tells MDF that the data contained in the curly braces needs to be formatted as a label. The rest of the field is treated just like a normal encyclopedic field.


\begin{quote}\small\itshape Note: Using the encyclopedic fields in this manner should not be a common practice, as this will make it impossible for Shoebox or any other program to know exactly what kind of data is contained in your encyclopedic fields. But in a pinch, it does allow for this kind of flexibility.\end{quote}

Identical national and regional language fields are:


{\color{cflink} \hyperref[key:en]{\textbackslash{}en} and \hyperref[key:er]{\textbackslash{}er}}


Related fields include:


{\color{cflink} \hyperref[key:de]{\textbackslash{}de,} \hyperref[key:ue]{\textbackslash{}ue,} and \hyperref[key:oe]{\textbackslash{}oe}}


For more on special character formatting codes, see:


{\color{cflink} \hyperref[key:Character_Style_Codes]{Character\_Style\_Codes}}


\begin{quote}\small\itshape This is a ``free-form'' field.   Punctuation and capitalization should be used as needed.\end{quote}

{\small <Optional>}\par

\subsection{\textbackslash{}eg  etymology-gloss}\label{key:eg}

The published gloss for the etymological reference is given here.


For more information, see:


{\color{cflink} \hyperref[key:et]{\textbackslash{}et}}


{\small <Optional>}\par

\subsection{\textbackslash{}en  encyclopedic information (national)}\label{key:en}

The national language equivalent to the \textbackslash{}ee field. This field should cover information that provides a more complete knowledge-base on the lexeme.


For more information, see:


{\color{cflink} \hyperref[key:ee]{\textbackslash{}ee}}


Closely related fields are:


{\color{cflink} \hyperref[key:dn]{\textbackslash{}dn,} \hyperref[key:un]{\textbackslash{}un,} and \hyperref[key:on]{\textbackslash{}on}}


\begin{quote}\small\itshape This is a ``free-form'' field.   Punctuation and capitalization should be used as needed.\end{quote}

{\small <Optional>}\par

\subsection{\textbackslash{}er  encyclopedic information (regional)}\label{key:er}

The regional language equivalent to the \textbackslash{}ee field. This field should cover information that provides a more complete knowledge-base on the lexeme.


MDF adds the brackets ``[ ]'' around this field. If included in the output, it is treated as part of the national field.


For more information, see:


{\color{cflink} \hyperref[key:ee]{\textbackslash{}ee}}


Closely related fields are:


{\color{cflink} \hyperref[key:dr]{\textbackslash{}dr,} \hyperref[key:ur]{\textbackslash{}ur,} and \hyperref[key:or]{\textbackslash{}or}}


\begin{quote}\small\itshape This is a ``free-form'' field.   Punctuation and capitalization should be used as needed.\end{quote}

{\small <Optional>}\par

\subsection{\textbackslash{}es  etymology-source}\label{key:es}

The reference or source abbreviation for etymology of the lexeme is given here.


For more information, see:


{\color{cflink} \hyperref[key:et]{\textbackslash{}et}}


\begin{quote}\small\itshape This field does not normally print.\end{quote}

{\small <Optional>}\par

\subsection{\textbackslash{}et  etymology}\label{key:et}

The etymology for the lexeme is put here, e.g.:


{\small\color{exlabel}(370)}\par
\noindent
\begin{Verbatim}[fontsize=\footnotesize,breaklines,breakanywhere]
 \et *babuy

\end{Verbatim}

This field is bundled with the following fields:


\begin{description}
\item[{\hyperref[key:eg]{\textcolor{cflink}{\textbackslash{}eg}}}] etymology-gloss
\item[{\hyperref[key:es]{\textcolor{cflink}{\textbackslash{}es}}}] etymology-source [doesn't print normally]
\item[{\hyperref[key:ec]{\textcolor{cflink}{\textbackslash{}ec}}}] etymology-comment [doesn't print normally]
\end{description}


A full example would be:


{\small\color{exlabel}(371)}\par
\noindent
\begin{Verbatim}[fontsize=\footnotesize,breaklines,breakanywhere]
 \et *ha[ñ]seng
\end{Verbatim}

{\small\color{exlabel}(372)}\par
\noindent
\begin{Verbatim}[fontsize=\footnotesize,breaklines,breakanywhere]
 \eg breathe loudly
\end{Verbatim}

{\small\color{exlabel}(373)}\par
\noindent
\begin{Verbatim}[fontsize=\footnotesize,breaklines,breakanywhere]
 \es PANDYMC
\end{Verbatim}

{\small\color{exlabel}(374)}\par
\noindent
\begin{Verbatim}[fontsize=\footnotesize,breaklines,breakanywhere]
 \ec *s > /h/? Not a likely cognate.

\end{Verbatim}

Which would print as either:


{\small\color{exlabel}(375)}\par
\noindent
\begin{Verbatim}[fontsize=\footnotesize,breaklines,breakanywhere]
 *ha[ñ]seng 'breathe loudly'.

\end{Verbatim}

Or:


{\small\color{exlabel}(376)}\par
\noindent
\begin{Verbatim}[fontsize=\footnotesize,breaklines,breakanywhere]
 *ha[ñ]seng 'breathe loudly' PANDYMC (*s > /h/? Not 
  a likely cognate.)

\end{Verbatim}

Depending on whether you have changed the default print settings to include the \textbackslash{}es and \textbackslash{}ec fields or not.


\begin{quote}\small\itshape You must explicitly include the asterisk (*) if you want it printed. The print tables will not add it.\end{quote}

{\small <Optional>}\par

\subsection{\textbackslash{}ev  encyclopedic information (vernacular)}\label{key:ev}

This field contains the vernacular description of any pertinent encyclopedic information related to the lexeme or headword, subentry, or sense. This is intended for use in a monolingual dictionary, but can be used in diglot and triglot dictionaries as well.


For more information, see:


{\color{cflink} \hyperref[key:ee]{\textbackslash{}ee}}


Related fields are:


{\color{cflink} \hyperref[key:gv]{\textbackslash{}gv,} \hyperref[key:dv]{\textbackslash{}dv,} \hyperref[key:ov]{\textbackslash{}ov,} \hyperref[key:uv]{\textbackslash{}uv}}


{\small <Optional>}\par

\subsection{\textbackslash{}gr  gloss (regional)}\label{key:gr}

This helps clarify the national language gloss in the prominent regional language of the area. This is also used as a regional definition in the printed dictionary if there is no \textbackslash{}dr field present (or it is present but empty).


The user may enter data either with each gloss getting its own field:


{\small\color{exlabel}(377)}\par
\noindent
\begin{Verbatim}[fontsize=\footnotesize,breaklines,breakanywhere]
  \gr kasi_pinda
\end{Verbatim}

{\small\color{exlabel}(378)}\par
\noindent
\begin{Verbatim}[fontsize=\footnotesize,breaklines,breakanywhere]
  \gr kasi_pulang

\end{Verbatim}

Or with the glosses strung together, separated by a semicolon with a space on either side of it (so as to be unlike any normal punctuation):


{\small\color{exlabel}(379)}\par
\noindent
\begin{Verbatim}[fontsize=\footnotesize,breaklines,breakanywhere]
  \ge kasi_pinda ; kasi_pulang

\end{Verbatim}

Toolbox can recognize either format and when interlinearizing will give the user the choice of both glosses in either case.


The underline character `\_' in the example glosses above is to force Toolbox to treat the multiple word gloss `kasi pinda' as a single gloss. MDF will convert underline characters in gloss fields to spaces automatically. The sequence ` ; ` is also converted to `, ` by MDF when formatting a dictionary or finderlist.


If this field is included in the output, MDF adds a label and brackets ``[ ]'' around it and it is treated as part of the national field.


For important related information, see:


\begin{description}
\item[{\hyperref[key:dr]{\textcolor{cflink}{\textbackslash{}dr}}}] definition (regional)
\item[{\hyperref[key:rr]{\textcolor{cflink}{\textbackslash{}rr}}}] reversal form (regional)
\end{description}


{\small <Optional>}\par

\subsection{\textbackslash{}gv  gloss (vernacular)}\label{key:gv}

This field is available for the development of a monolingual dictionary. And while this may seem somewhat meaningless (a vernacular gloss for a vernacular word), it actually can be a useful place to store the simple explanations the researcher is given by a language assistant concerning the meaning of the headword or lexeme. These can then be formulated into a more exact definition and transfer to the \textbackslash{}dv field.  It could also be used for those lexemes which can actually be defined with short glosses. It could also serve as a subset to the information covered in \textbackslash{}dv field.


See:


{\color{cflink} \hyperref[key:dv]{\textbackslash{}dv}}


{\small <Optional>}\par

\subsection{\textbackslash{}is  index  of  semantics}\label{key:is}

``Beyond key terms: a lexicon useful for translation'' (Rick Nivens, manuscript) discusses the use of Louw and Nida's (1988) Greek-English semantic domain categories. While Nivens proposes another way of doing this in his paper, it is also possible to use this \textbackslash{}is field to catalog lexical entries according to these semantic domains. Reversing on this field would then yield semantically related entries (in relation to the New Testament).


One word of caution: this is an etic approach, i.e. Greek semantics will rarely line up exactly with the vernacular domains. This field is to be a tool, and no attempt should be made to ``force-fit'' lexemes into pre-defined domains.


If selected for output, MDF adds the label ``Semantics: `` to this field.


For related fields see:


\begin{description}
\item[{\hyperref[key:sd]{\textcolor{cflink}{\textbackslash{}sd}}}] semantic domain
\item[{\hyperref[key:th]{\textcolor{cflink}{\textbackslash{}th}}}] thesaurus
\end{description}


For references see:


{\color{cflink} \hyperref[key:References]{References}}


\begin{quote}\small\itshape This field does not normally print.\end{quote}

{\small <Optional>}\par

\subsection{\textbackslash{}le  lexical function (English gloss)}\label{key:le}

This is for giving the English gloss of the vernacular lexeme referenced by the lexical function.


For more information, see:


\begin{description}
\item[{\hyperref[key:lf]{\textcolor{cflink}{\textbackslash{}lf}}}] lexical function label
\item[{\hyperref[key:lv]{\textcolor{cflink}{\textbackslash{}lv}}}] vernacular lexeme referenced by the lexical function
\end{description}


Related fields are:


{\color{cflink} \hyperref[key:ln]{\textbackslash{}ln} and \hyperref[key:lr]{\textbackslash{}lr}}


For generic cross-referencing, see:


\begin{description}
\item[{\hyperref[key:cf]{\textcolor{cflink}{\textbackslash{}cf}}}] confer/cross-reference
\end{description}


{\small <Optional>}\par

\subsection{\textbackslash{}lf  lexical function label}\label{key:lf}

For encoding the semantic networks of a language. The \textbackslash{}lf field bundle includes the following fields:


\begin{description}
\item[{\hyperref[key:lf]{\textcolor{cflink}{\textbackslash{}lf}}}] the lexical function
\item[{\hyperref[key:or]{\textcolor{cflink}{or}}}] relationship label
\item[{\hyperref[key:lv]{\textcolor{cflink}{\textbackslash{}lv}}}] vernacular lexeme referenced by that lexical function/relation
\item[{\hyperref[key:le]{\textcolor{cflink}{\textbackslash{}le}}}] English gloss of the vernacular lexeme
\item[{\hyperref[key:ln]{\textcolor{cflink}{\textbackslash{}ln}}}] national gloss of the vernacular lexeme
\item[{\hyperref[key:lr]{\textcolor{cflink}{\textbackslash{}lr}}}] regional gloss of the vernacular lexeme
\end{description}


MDF supports two ways of using the \textbackslash{}lf bundles. A typical example using the original (mixed data) method would be:


{\small\color{exlabel}(380)}\par
\noindent
\begin{Verbatim}[fontsize=\footnotesize,breaklines,breakanywhere]
  \lf Syn
\end{Verbatim}

{\small\color{exlabel}(381)}\par
\noindent
\begin{Verbatim}[fontsize=\footnotesize,breaklines,breakanywhere]
  \lv -dew
\end{Verbatim}

{\small\color{exlabel}(382)}\par
\noindent
\begin{Verbatim}[fontsize=\footnotesize,breaklines,breakanywhere]
  \le chop on end
\end{Verbatim}

{\small\color{exlabel}(383)}\par
\noindent
\begin{Verbatim}[fontsize=\footnotesize,breaklines,breakanywhere]
  \ln
\end{Verbatim}

{\small\color{exlabel}(384)}\par
\noindent
\begin{Verbatim}[fontsize=\footnotesize,breaklines,breakanywhere]
  \lr

\end{Verbatim}

(Note that empty fields are okay.)


This lexical function would be typically print as:


{\small\color{exlabel}(385)}\par
\noindent
\begin{Verbatim}[fontsize=\footnotesize,breaklines,breakanywhere]
   Syn: -dew 'chop on end'.

\end{Verbatim}

If the other glosses were filled in, they would be included in the printout (if a triglot dictionary was requested). Multiple lexical function bundles are concatenated with a semicolon `; `, e.g.:


{\small\color{exlabel}(386)}\par
\noindent
\begin{Verbatim}[fontsize=\footnotesize,breaklines,breakanywhere]
   Syn: -dew 'chop on end'; Syn: -sin 'split'.

\end{Verbatim}

Note: In the original MDF manual and even in the revision of Chapter 2 you will see examples in which the label and the vernacular word were both in the \textbackslash{}lf field. For example, \textbackslash{}lf Syn = -dew. This approach is now strongly discouraged. The user is strongly encouraged to use the \textbackslash{}lv field as shown in the examples in this record.


The recommended MDF method for handling lexical functions (using the \textbackslash{}lv field) has some distinct advantages over the old way:


1) Since the lexical function label sits by itself in the \textbackslash{}lf field, you can  use the Toolbox Range Set feature to maintain consistent labeling.


2) Using \textbackslash{}lf and \textbackslash{}lv fields with the sort and jump features of Toolbox is much easier.


3) Browsing on a database sorted by the \textbackslash{}lv and \textbackslash{}lf fields and displaying the \textbackslash{}lv, \textbackslash{}lf, and \textbackslash{}le fields will give a virtual \textbackslash{}lv reversal view of your data -- allowing you to see all of the lexical functions that any given lexeme has been categorized under.


4) New Toolbox features like integrity checks, range-set building, and word lists will not work properly on the \textbackslash{}lf field if it contains mixed data.


5) If the vernacular language is written in a non-Roman script, the lexeme and lexical function label need to be in different fields since they have different language encodings.


See also:


\begin{description}
\item[{\hyperref[key:lv]{\textcolor{cflink}{\textbackslash{}lv}}}] vernacular lexeme referenced by the lexical function label
\end{description}


For more information on Range Sets, see:


{\color{cflink} \hyperref[key:Range_Sets]{Range\_Sets}}


For an alternative way to create a table of related lexical items, see:


\begin{description}
\item[{\hyperref[key:tb]{\textcolor{cflink}{\textbackslash{}tb}}}] table
\end{description}


For generic cross-referencing, see:


\begin{description}
\item[{\hyperref[key:cf]{\textcolor{cflink}{\textbackslash{}cf}}}] confer/cross-reference
\end{description}


For more detailed information on language analysis through lexical functions, see the MDF field manual (Coward and Grimes, 1995) and the article ``Mapping a culture through networks of meaning'' (Grimes 1987). For referential information on these, see:


{\color{cflink} \hyperref[key:References]{References}}


{\small <Optional>}\par

\subsection{\textbackslash{}ln  lexical function (national gloss)}\label{key:ln}

This is for giving the national gloss of the vernacular lexeme referenced by the lexical function.


For more information, see:


\begin{description}
\item[{\hyperref[key:lf]{\textcolor{cflink}{\textbackslash{}lf}}}] lexical function label
\item[{\hyperref[key:lv]{\textcolor{cflink}{\textbackslash{}lv}}}] vernacular lexeme referenced by the lexical function
\end{description}


Related fields are:


{\color{cflink} \hyperref[key:le]{\textbackslash{}le} and \hyperref[key:lr]{\textbackslash{}lr}}


For generic cross-referencing, see:


\begin{description}
\item[{\hyperref[key:cf]{\textcolor{cflink}{\textbackslash{}cf}}}] confer/cross-reference
\end{description}


{\small <Optional>}\par

\subsection{\textbackslash{}lr  lexical function (regional gloss)}\label{key:lr}

This is for giving the regional gloss of the vernacular lexeme referenced by the lexical function.


For more information, see:


\begin{description}
\item[{\hyperref[key:lf]{\textcolor{cflink}{\textbackslash{}lf}}}] lexical function label
\item[{\hyperref[key:lv]{\textcolor{cflink}{\textbackslash{}lv}}}] vernacular lexeme referenced by the lexical function
\end{description}


Related fields are:


{\color{cflink} \hyperref[key:le]{\textbackslash{}le} and \hyperref[key:ln]{\textbackslash{}ln}}


For generic cross-referencing, see:


\begin{description}
\item[{\hyperref[key:cf]{\textcolor{cflink}{\textbackslash{}cf}}}] confer/cross-reference
\end{description}


{\small <Optional>}\par

\subsection{\textbackslash{}lt  literal meaning}\label{key:lt}

Used to elucidate the destinct meanings of the parts of an idiom or complex phrase in a lexical entry (\textbackslash{}lx) or subentry (\textbackslash{}se).


{\small\color{exlabel}(387)}\par
\noindent
\begin{Verbatim}[fontsize=\footnotesize,breaklines,breakanywhere]
   \lx setengah mati
\end{Verbatim}

{\small\color{exlabel}(388)}\par
\noindent
\begin{Verbatim}[fontsize=\footnotesize,breaklines,breakanywhere]
   \ps idm
\end{Verbatim}

{\small\color{exlabel}(389)}\par
\noindent
\begin{Verbatim}[fontsize=\footnotesize,breaklines,breakanywhere]
   \ge exhausting ; tiring ; frustrating
\end{Verbatim}

{\small\color{exlabel}(390)}\par
\noindent
\begin{Verbatim}[fontsize=\footnotesize,breaklines,breakanywhere]
   \de used to describe a situation, activity, or person as being
       exhausting, tiring, or frustrating
\end{Verbatim}

{\small\color{exlabel}(391)}\par
\noindent
\begin{Verbatim}[fontsize=\footnotesize,breaklines,breakanywhere]
   \lt half dead

\end{Verbatim}

MDF adds the label ``Lit: `` to this field and adds single quote marks around the data.


{\small <Optional>}\par

\subsection{\textbackslash{}lv  vernacular lexeme referenced by the lexical function}\label{key:lv}

This field is used in MDF for encoding the vernacular lexeme in a lexical function network. The \textbackslash{}lv field ``points to'' the vernacular lexeme (morpheme, word, or phrase) that is semantically related to the current headword as mapped or cataloged by the label in the \textbackslash{}lf field. An example of its use would be:


{\small\color{exlabel}(392)}\par
\noindent
\begin{Verbatim}[fontsize=\footnotesize,breaklines,breakanywhere]
   \lx feten            <-- the headword/lexeme
\end{Verbatim}

{\small\color{exlabel}(393)}\par
\noindent
\begin{Verbatim}[fontsize=\footnotesize,breaklines,breakanywhere]
   \ps n
\end{Verbatim}

{\small\color{exlabel}(394)}\par
\noindent
\begin{Verbatim}[fontsize=\footnotesize,breaklines,breakanywhere]
   \ge millet
\end{Verbatim}

{\small\color{exlabel}(395)}\par
\noindent
\begin{Verbatim}[fontsize=\footnotesize,breaklines,breakanywhere]
   ...
\end{Verbatim}

{\small\color{exlabel}(396)}\par
\noindent
\begin{Verbatim}[fontsize=\footnotesize,breaklines,breakanywhere]
   \lf Gen              <-- the lexical function/relationship
\end{Verbatim}

{\small\color{exlabel}(397)}\par
\noindent
\begin{Verbatim}[fontsize=\footnotesize,breaklines,breakanywhere]
   \lv agat             <-- the related lexeme
\end{Verbatim}

{\small\color{exlabel}(398)}\par
\noindent
\begin{Verbatim}[fontsize=\footnotesize,breaklines,breakanywhere]
   \le grain            <-- the English gloss of 'agat'
\end{Verbatim}

{\small\color{exlabel}(399)}\par
\noindent
\begin{Verbatim}[fontsize=\footnotesize,breaklines,breakanywhere]
   ...

\end{Verbatim}

In this example, `agat' (`grain') is in a Generic relationship to `feten' (`millet'). Note that the lexeme `agat' should also appear in the lexicon.


For more information, see:


\begin{description}
\item[{\hyperref[key:lf]{\textcolor{cflink}{\textbackslash{}lf}}}] lexical function label
\end{description}


Glossing fields for this vernacular lexeme are:


{\color{cflink} \hyperref[key:le]{\textbackslash{}le,} \hyperref[key:ln]{\textbackslash{}ln,} and \hyperref[key:lr]{\textbackslash{}lr}}


For generic cross-referencing, see:


\begin{description}
\item[{\hyperref[key:cf]{\textcolor{cflink}{\textbackslash{}cf}}}] cross-reference
\end{description}


{\small <Optional>}\par

\subsection{\textbackslash{}mn  main entry form}\label{key:mn}

This is used to reference a minimal entry back to its main entry. Minimal or ``minor'' entries consist of abbreviated information, e.g.:


{\small\color{exlabel}(400)}\par
\noindent
\begin{Verbatim}[fontsize=\footnotesize,breaklines,breakanywhere]
  \lx tado
  \ps adj
  \ge calm
  \mn teduh

\end{Verbatim}

This would typically print as:


{\small\color{exlabel}(401)}\par
\noindent
\begin{Verbatim}[fontsize=\footnotesize,breaklines,breakanywhere]
   tado  adj. calm; See main entry: teduh.

\end{Verbatim}

In this example, the minor entry ``tado'' has minimal information and contains a reference back to the main semantically related entry ``teduh''. This cross-referencing is done through the \textbackslash{}mn field:


{\small\color{exlabel}(402)}\par
\noindent
\begin{Verbatim}[fontsize=\footnotesize,breaklines,breakanywhere]
  \mn teduh

\end{Verbatim}

The main entry ``teduh'' would be more verbose, containing all of the information normally included in a lexical entry. The ``teduh'' entry would also have the field ``\textbackslash{}va tado'' to reference the variant (dialect) form.


For related information, see:


\begin{description}
\item[{\hyperref[key:va]{\textcolor{cflink}{\textbackslash{}va}}}] variant forms
\end{description}


For a detailed discussion of the issues involved in using minor (or minimal) entries, see:


{\color{cflink} \hyperref[key:Using_Subentries_or_Lexical_Entries]{Using\_Subentries\_or\_Lexical\_Entries}}


{\small <Optional>}\par

\subsection{\textbackslash{}mr  morphemic representation}\label{key:mr}

This can be used to show the underlying morphemic structure for complex lexemes. MDF gives this field the label `Morph:'


In Buru, the lexeme `agat' (`grain') is a complex morpheme. Its internal structure should be encoded as:


{\small\color{exlabel}(403)}\par
\noindent
\begin{Verbatim}[fontsize=\footnotesize,breaklines,breakanywhere]
 \mr aga-t

\end{Verbatim}

This will typically print as:


{\small\color{exlabel}(404)}\par
\noindent
\begin{Verbatim}[fontsize=\footnotesize,breaklines,breakanywhere]
 Morph: aga-t.

\end{Verbatim}

{\small <Optional>}\par

\subsection{\textbackslash{}na  notes  on   anthropology}\label{key:na}

This field is for any ethnographic note that is pertinent to the lexeme that you wish to keep separate.


If selected for output, MDF adds the label ``Anth: `` to this field and brackets it with ``[ ]''.


For generic notes, use:


{\color{cflink} \hyperref[key:nt]{\textbackslash{}nt}}


For closely related fields, designed to print (not just for notes), see:


\begin{description}
\item[{\hyperref[key:ee]{\textcolor{cflink}{\textbackslash{}ee,}}}]
\item[{\hyperref[key:en]{\textcolor{cflink}{\textbackslash{}en,}}}]
\item[{\hyperref[key:er]{\textcolor{cflink}{\textbackslash{}er}}}] encyclopedic field bundle
\end{description}


Earlier versions of MDF recognized two types of ethnographic fields, this one (\textbackslash{}na), and \textbackslash{}en (which has been discontinued as an ethnographic field).


\begin{quote}\small\itshape This is a ``free-form'' field.   Punctuation and capitalization should be used as needed.\end{quote}

{\small <Optional>}\par

\subsection{\textbackslash{}nd  notes on discourse}\label{key:nd}

This is a place for your notes on discourse/text analysis, should you wish to keep them separate.


If selected for output, MDF adds the label ``Disc: `` to this field and brackets it with ``[ ]''.


For generic notes, see:


{\color{cflink} \hyperref[key:nt]{\textbackslash{}nt}}


\begin{quote}\small\itshape This is a ``free-form'' field.   Punctuation and capitalization should be used as needed.\end{quote}

{\small <Optional>}\par

\subsection{\textbackslash{}ng  notes on grammar}\label{key:ng}

This is a place for your grammar notes, should you wish to keep them separate.


If selected for output, MDF adds the label ``Gram: `` to this field and brackets it with ``[ ]''.


For generic notes, see:


{\color{cflink} \hyperref[key:nt]{\textbackslash{}nt}}


\begin{quote}\small\itshape This is a ``free-form'' field.   Punctuation and capitalization should be used as needed.\end{quote}

{\small <Optional>}\par

\subsection{\textbackslash{}np  notes on phonology}\label{key:np}

This is a place for your phonology notes, should you wish to keep them separate.


If selected for output, MDF adds the label ``Phon: `` to this field and brackets it with ``[ ]''.


For generic notes, see:


{\color{cflink} \hyperref[key:nt]{\textbackslash{}nt}}


\begin{quote}\small\itshape This is a ``free-form'' field.   Punctuation and capitalization should be used as needed.\end{quote}

{\small <Optional>}\par

\subsection{\textbackslash{}nq  questions}\label{key:nq}

This is a place for your questions, should you wish to keep them separate.


If selected for output, MDF adds the label ``Ques: `` to this field and brackets it with ``[ ]''.


For generic notes, see:


{\color{cflink} \hyperref[key:nt]{\textbackslash{}nt}}


\begin{quote}\small\itshape This is a ``free-form'' field.   Punctuation and capitalization should be used as needed.\end{quote}

{\small <Optional>}\par

\subsection{\textbackslash{}ns  notes on sociolinguistics}\label{key:ns}

This is used for encoding sociolinguistic notes. Dialect information (i.e. which villages use this lexeme form, etc.) should be explained in the following fields:


\begin{description}
\item[{\hyperref[key:ue]{\textcolor{cflink}{\textbackslash{}ue,}}}]
\item[{\hyperref[key:un]{\textcolor{cflink}{\textbackslash{}un,}}}]
\item[{\hyperref[key:ur]{\textcolor{cflink}{\textbackslash{}ur}}}] (Usage)
\item[{\hyperref[key:ee]{\textcolor{cflink}{\textbackslash{}ee,}}}]
\item[{\hyperref[key:en]{\textcolor{cflink}{\textbackslash{}en,}}}]
\item[{\hyperref[key:er]{\textcolor{cflink}{\textbackslash{}er}}}] (Encyclopedic)
\item[{\hyperref[key:oe]{\textcolor{cflink}{\textbackslash{}oe,}}}]
\item[{\hyperref[key:on]{\textcolor{cflink}{\textbackslash{}on,}}}]
\item[{\hyperref[key:or]{\textcolor{cflink}{\textbackslash{}or}}}] (Only --restrictions)
\end{description}


If selected for output, MDF adds the label ``Socio: `` to this field and brackets it with ``[ ]''.


For generic notes, see:


{\color{cflink} \hyperref[key:nt]{\textbackslash{}nt}}


\begin{quote}\small\itshape This is a ``free-form'' field.   Punctuation and capitalization should be used as needed.\end{quote}

{\small <Optional>}\par

\subsection{\textbackslash{}oe  only [restrictions] (English)}\label{key:oe}

This is for encoding any semantic and/or grammatical restrictions pertinent to the lexeme.


In Buru, the lexeme `anafina' is glossed as `female'. But this needs to be restricted to only `human' references, which could be encoded as:


{\small\color{exlabel}(405)}\par
\noindent
\begin{Verbatim}[fontsize=\footnotesize,breaklines,breakanywhere]
  \oe Human

\end{Verbatim}

and would print as:


{\small\color{exlabel}(406)}\par
\noindent
\begin{Verbatim}[fontsize=\footnotesize,breaklines,breakanywhere]
   Restrict: Human.

\end{Verbatim}

Since this is a free-form field, you can be as verbose as needed, but in many cases, such as this one, a simple code may suffice.


Closely related English fields are:


{\color{cflink} \hyperref[key:de]{\textbackslash{}de,} \hyperref[key:ue]{\textbackslash{}ue,} and \hyperref[key:ee]{\textbackslash{}ee}}


Identical national and regional language fields are:


{\color{cflink} \hyperref[key:on]{\textbackslash{}on} and \hyperref[key:or]{\textbackslash{}or}}


\begin{quote}\small\itshape This is a ``free-form'' field. Punctuation and capitalization should be used as needed.\end{quote}

{\small <Optional>}\par

\subsection{\textbackslash{}on  only [restrictions] (national)}\label{key:on}

The national language equivalent to the \textbackslash{}oe field. This field is for clarifying semantic and grammatical restrictions pertinent to the lexeme. See the \textbackslash{}oe field for more information.


MDF initially adds the label ``NatRestrict: `` to this field. This is later changed to whatever national language label is specified for this field in the national audience CC table. For more information on this, see:


{\color{cflink} \hyperref[key:Printed_Field_Labels]{Printed\_Field\_Labels}}


For more information, see:


{\color{cflink} \hyperref[key:oe]{\textbackslash{}oe}}


Closely related national fields are:


{\color{cflink} \hyperref[key:dn]{\textbackslash{}dn,} \hyperref[key:un]{\textbackslash{}un,} and \hyperref[key:en]{\textbackslash{}en}}


\begin{quote}\small\itshape This is a ``free-form'' field. Punctuation and capitalization should be used as needed.\end{quote}

{\small <Optional>}\par

\subsection{\textbackslash{}or  only [restrictions] (regional)}\label{key:or}

The regional language equivalent to the \textbackslash{}oe field. This field is for clarifying semantic and grammatical restrictions pertinent to the lexeme.


If included in the output, MDF adds brackets ``[ ]'' around this field, and it is treated as part of the national field.


For more information, see:


{\color{cflink} \hyperref[key:oe]{\textbackslash{}oe}}


Closely related regional fields are:


{\color{cflink} \hyperref[key:dr]{\textbackslash{}dr,} \hyperref[key:ur]{\textbackslash{}ur,} and \hyperref[key:er]{\textbackslash{}er}}


\begin{quote}\small\itshape This is a ``free-form'' field. Punctuation and capitalization should be used as needed.\end{quote}

{\small <Optional>}\par

\subsection{\textbackslash{}ov  only   [restrictions]   (vernacular)}\label{key:ov}

This field contains the vernacular description of any semantic and/or grammatical restrictions pertinent to the lexeme or headword. This is intended for use in a monolingual dictionary, but can be used in diglot and triglot dictionaries as well.


MDF initially adds the label ``VerRestrict: `` to this field. This is not changed by either the English or national audience CC table. But you may add a rule to change this if you wish. For more information on how to do this, see:


{\color{cflink} \hyperref[key:Printed_Field_Labels]{Printed\_Field\_Labels}}


For more information, see:


{\color{cflink} \hyperref[key:oe]{\textbackslash{}oe}}


Related vernacular fields are:


{\color{cflink} \hyperref[key:gv]{\textbackslash{}gv,} \hyperref[key:dv]{\textbackslash{}dv,} \hyperref[key:ev]{\textbackslash{}ev,} \hyperref[key:uv]{\textbackslash{}uv}}


{\small <Optional>}\par

\subsection{\textbackslash{}pc  picture}\label{key:pc}

This contains either the book and page number reference of a relevant picture (you've either sketched in a notebook or found in a picture book), or a graphics link for a PCX file of a picture you want to include in your printed dictionary.


If the field does not begin with `.G.' as in a graphics link, then the field is treated as ``free-form''. The field is output near the end of an entry paragraph, and the print tables put parentheses ( ) around the whole field (to set it off from other types of information).


If the \textbackslash{}pc field begins with a graphics link mark (.G.), then it is NOT a ``free-form'' field and should follow the basic form below:


{\small\color{exlabel}(407)}\par
\noindent
\begin{Verbatim}[fontsize=\footnotesize,breaklines,breakanywhere]
  \pc .G.sago.pcx;1.5";1";PCX

\end{Verbatim}

The .G. marks this as a graphics link; next follows the filename of the picture; next is the width (here 1.5'') and then the height (1'') and finally the graphics format (PCX). Each bit of information is separated by a semicolon.


When the dictionary is formatted, the graphics information is moved to the beginning of the entry, subentry, or sense in which the \textbackslash{}pc field is found. This will cause the text to flow around the picture. Sizes much larger than 1.5''x1.5'' are not recommended. In double column format the picture is placed flush right in the column; in single column format the picture is flush right to the right margin.


\begin{quote}\small\itshape This is a ``free-form'' field, if it is not a Graphics Link paragraph. Punctuation and capitalization should be used as needed.\end{quote}

{\small <Optional>}\par

\subsection{\textbackslash{}pd  paradigm set}\label{key:pd}

This is used for specifying the noun or verb class, gender, or other paradigm set that the lexeme or headword is associated with. These classes are generally given labels or numbers to differentiate them. It is assumed that the classes are explained elsewhere (like the introduction to the dictionary) and all that is needed here is the class label or number, e.g.:


{\small\color{exlabel}(408)}\par
\noindent
\begin{Verbatim}[fontsize=\footnotesize,breaklines,breakanywhere]
 \pd 3 

\end{Verbatim}

Mnemonic codes are actually better than numbers (since what the number represents is very difficult to remember), e.g.:


{\small\color{exlabel}(409)}\par
\noindent
\begin{Verbatim}[fontsize=\footnotesize,breaklines,breakanywhere]
 \pd met:vrt

\end{Verbatim}

might refer to a ``verb root which metathesizes with its subject prefix''. This is much easier to decode than ``3''. For consistency be sure to use Toolbox's Range Set feature on this field. For more information, see:


{\color{cflink} \hyperref[key:Range_Sets]{Range\_Sets}}


This would typically print as:


{\small\color{exlabel}(410)}\par
\noindent
\begin{Verbatim}[fontsize=\footnotesize,breaklines,breakanywhere]
   Prdm: met:vrt.

\end{Verbatim}

To give the actual vernacular form for a paradigm (especially needed where the paradigm is incomplete or is irregular in form), MDF provides the following fields:


\begin{description}
\item[{\hyperref[key:pdl]{\textcolor{cflink}{\textbackslash{}pdl}}}] paradigm label
\item[{\hyperref[key:pdv]{\textcolor{cflink}{\textbackslash{}pdv}}}] paradigm vernacular form
\item[{\hyperref[key:pde]{\textcolor{cflink}{\textbackslash{}pde}}}] paradigm form-English gloss
\item[{\hyperref[key:pdn]{\textcolor{cflink}{\textbackslash{}pdn}}}] paradigm form-national gloss
\item[{\hyperref[key:pdr]{\textcolor{cflink}{\textbackslash{}pdr}}}] paradigm form-regional gloss
\end{description}


{\small <Optional>}\par

\subsection{\textbackslash{}pde  paradigm form (English gloss)}\label{key:pde}

This is used for glossing the vernacular paradigm form in English, e.g.:


{\small\color{exlabel}(411)}\par
\noindent
\begin{Verbatim}[fontsize=\footnotesize,breaklines,breakanywhere]
  \pdl Pl
\end{Verbatim}

{\small\color{exlabel}(412)}\par
\noindent
\begin{Verbatim}[fontsize=\footnotesize,breaklines,breakanywhere]
  \pdv asure
\end{Verbatim}

{\small\color{exlabel}(413)}\par
\noindent
\begin{Verbatim}[fontsize=\footnotesize,breaklines,breakanywhere]
  \pde dogs
\end{Verbatim}

{\small\color{exlabel}(414)}\par
\noindent
\begin{Verbatim}[fontsize=\footnotesize,breaklines,breakanywhere]
  \pdn anjing-anjing

\end{Verbatim}

This will typically print as:


{\small\color{exlabel}(415)}\par
\noindent
\begin{Verbatim}[fontsize=\footnotesize,breaklines,breakanywhere]
    Pl: asure 'dogs' 'anjing-anjing'.

\end{Verbatim}

where the paradigm label is italic, the vernacular form is in vernacular font, and the two glosses are in their appropriate language fonts.


For a more detailed explanation of this set of fields, see:


\begin{description}
\item[{\hyperref[key:pdl]{\textcolor{cflink}{\textbackslash{}pdl}}}] paradigm label
\item[{\hyperref[key:pdv]{\textcolor{cflink}{\textbackslash{}pdv}}}] paradigm vernacular form
\end{description}


For the related glossing fields, see:


\begin{description}
\item[{\hyperref[key:pdn]{\textcolor{cflink}{\textbackslash{}pdn}}}] paradigm form-national gloss
\item[{\hyperref[key:pdr]{\textcolor{cflink}{\textbackslash{}pdr}}}] paradigm form-regional gloss
\end{description}


For more information on the paradigm fields, see:


\begin{description}
\item[{\hyperref[key:pd]{\textcolor{cflink}{\textbackslash{}pd}}}] paradigm set
\end{description}


{\small <Optional>}\par

\subsection{\textbackslash{}pdl  paradigm label}\label{key:pdl}

The paradigm field (\textbackslash{}pd) field is used to define the general paradigm set a headword  or lexeme is associated with. But it is also useful to give the actual forms for a paradigm, especially where a paradigm set is incomplete or irregular. For example, the third-singular genitive marker normally possesses nouns with the form `-na', except with certain nouns, where its form is irregular. Any noun that takes an irregular 3sGen form should be marked:


{\small\color{exlabel}(416)}\par
\noindent
\begin{Verbatim}[fontsize=\footnotesize,breaklines,breakanywhere]
 \lx ow
\end{Verbatim}

{\small\color{exlabel}(417)}\par
\noindent
\begin{Verbatim}[fontsize=\footnotesize,breaklines,breakanywhere]
 \ps n
\end{Verbatim}

{\small\color{exlabel}(418)}\par
\noindent
\begin{Verbatim}[fontsize=\footnotesize,breaklines,breakanywhere]
 \ge jackfruit
\end{Verbatim}

{\small\color{exlabel}(419)}\par
\noindent
\begin{Verbatim}[fontsize=\footnotesize,breaklines,breakanywhere]
 \gn nangka
\end{Verbatim}

{\small\color{exlabel}(420)}\par
\noindent
\begin{Verbatim}[fontsize=\footnotesize,breaklines,breakanywhere]
 \pd n-fd
\end{Verbatim}

{\small\color{exlabel}(421)}\par
\noindent
\begin{Verbatim}[fontsize=\footnotesize,breaklines,breakanywhere]
 \pdl 3sGEN
\end{Verbatim}

{\small\color{exlabel}(422)}\par
\noindent
\begin{Verbatim}[fontsize=\footnotesize,breaklines,breakanywhere]
 \pdv owan
\end{Verbatim}

{\small\color{exlabel}(423)}\par
\noindent
\begin{Verbatim}[fontsize=\footnotesize,breaklines,breakanywhere]
 \pde his/her jackfruit
\end{Verbatim}

{\small\color{exlabel}(424)}\par
\noindent
\begin{Verbatim}[fontsize=\footnotesize,breaklines,breakanywhere]
 \pdn nangkanya

\end{Verbatim}

These paradigm label fields would typically print as:


{\small\color{exlabel}(425)}\par
\noindent
\begin{Verbatim}[fontsize=\footnotesize,breaklines,breakanywhere]
    3sGEN: owan 'his/her jackfruit' 'nangkanya'.

\end{Verbatim}

(where the ``3sGEN'' is formatted as italic and the vernacular form and the glossing are formatted with the appropriate language fonts)


\begin{quote}\small\itshape Note: The paradigm label (\textbackslash{}pdl) fields can be used with or without the paradigm set (\textbackslash{}pd) field.\end{quote}

{\small\color{exlabel}(426)}\par
\noindent
\begin{Verbatim}[fontsize=\footnotesize,breaklines,breakanywhere]
You can use any label you need to mark a paradigm. Use Toolbox's Range Set 
feature to insure consistent labeling. For more information, see: 
\end{Verbatim}

{\color{cflink} \hyperref[key:Range_Sets]{Range\_Sets}}


\begin{quote}\small\itshape Note: All the fields in this bundle are made up of markers using three-letter codes.\end{quote}

The related fields to this set are:


\begin{description}
\item[{\hyperref[key:pdv]{\textcolor{cflink}{\textbackslash{}pdv}}}] paradigm vernacular form
\item[{\hyperref[key:pde]{\textcolor{cflink}{\textbackslash{}pde}}}] paradigm form-English gloss
\item[{\hyperref[key:pdn]{\textcolor{cflink}{\textbackslash{}pdn}}}] paradigm form-national gloss
\item[{\hyperref[key:pdr]{\textcolor{cflink}{\textbackslash{}pdr}}}] paradigm form-regional gloss
\end{description}


For more information on the paradigm fields, see:


\begin{description}
\item[{\hyperref[key:pd]{\textcolor{cflink}{\textbackslash{}pd}}}] paradigm set
\end{description}


{\small <Optional>}\par

\subsection{\textbackslash{}pdn  paradigm form (national gloss)}\label{key:pdn}

This is used for glossing the vernacular paradigm form in the national language. For the related glossing fields, see:


\begin{description}
\item[{\hyperref[key:pde]{\textcolor{cflink}{\textbackslash{}pde}}}] paradigm form-English gloss
\item[{\hyperref[key:pdr]{\textcolor{cflink}{\textbackslash{}pdr}}}] paradigm form-regional gloss
\end{description}


For a more detailed explanation of this set of fields, see:


\begin{description}
\item[{\hyperref[key:pdl]{\textcolor{cflink}{\textbackslash{}pdl}}}] paradigm label
\item[{\hyperref[key:pdv]{\textcolor{cflink}{\textbackslash{}pdv}}}] paradigm vernacular form
\end{description}


For more information on the paradigm fields, see:


\begin{description}
\item[{\hyperref[key:pd]{\textcolor{cflink}{\textbackslash{}pd}}}] paradigm set
\end{description}


{\small <Optional>}\par

\subsection{\textbackslash{}pdr  paradigm form (regional gloss)}\label{key:pdr}

This is used for glossing the vernacular paradigm form in the regional language. For the related glossing fields, see:


\begin{description}
\item[{\hyperref[key:pde]{\textcolor{cflink}{\textbackslash{}pde}}}] paradigm form-English gloss
\item[{\hyperref[key:pdn]{\textcolor{cflink}{\textbackslash{}pdn}}}] paradigm form-national gloss
\end{description}


For a more detailed explanation of this set of fields, see:


\begin{description}
\item[{\hyperref[key:pdl]{\textcolor{cflink}{\textbackslash{}pdl}}}] paradigm label
\item[{\hyperref[key:pdv]{\textcolor{cflink}{\textbackslash{}pdv}}}] paradigm vernacular form
\end{description}


For more information on the paradigm fields, see:


\begin{description}
\item[{\hyperref[key:pd]{\textcolor{cflink}{\textbackslash{}pd}}}] paradigm set
\end{description}


{\small <Optional>}\par

\subsection{\textbackslash{}pdv  paradigm vernacular form}\label{key:pdv}

This field is used to give the vernacular paradigm form specified by the label in the \textbackslash{}pdl field, e.g.:


{\small\color{exlabel}(427)}\par
\noindent
\begin{Verbatim}[fontsize=\footnotesize,breaklines,breakanywhere]
 \pdl 1s
\end{Verbatim}

{\small\color{exlabel}(428)}\par
\noindent
\begin{Verbatim}[fontsize=\footnotesize,breaklines,breakanywhere]
 \pdv koban  

\end{Verbatim}

where ``koban'' is the first-singular verb form of the verb root ``-oban''. MDF will format the `1s' label as italic, and the contents of the \textbackslash{}pdv field as vernacular text.


For a more detailed explanation, see:


\begin{description}
\item[{\hyperref[key:pdl]{\textcolor{cflink}{\textbackslash{}pdl}}}] paradigm label
\end{description}


For the related glossing fields, see:


\begin{description}
\item[{\hyperref[key:pde]{\textcolor{cflink}{\textbackslash{}pde}}}] paradigm form-English gloss
\item[{\hyperref[key:pdn]{\textcolor{cflink}{\textbackslash{}pdn}}}] paradigm form-national gloss
\item[{\hyperref[key:pdr]{\textcolor{cflink}{\textbackslash{}pdr}}}] paradigm form-regional gloss
\end{description}


For more information on the paradigm fields, see:


\begin{description}
\item[{\hyperref[key:pd]{\textcolor{cflink}{\textbackslash{}pd}}}] paradigm set
\end{description}


{\small <Optional>}\par

\subsection{\textbackslash{}ph  phonetic/phonemic form}\label{key:ph}

This can be used as needed to retain the phonetic information that is lost when an orthographic spelling is used for an entry.


The print process can place square brackets around the data. For example:


{\small\color{exlabel}(429)}\par
\noindent
\begin{Verbatim}[fontsize=\footnotesize,breaklines,breakanywhere]
 \ph apa

\end{Verbatim}

would typically print as:


{\small\color{exlabel}(430)}\par
\noindent
\begin{Verbatim}[fontsize=\footnotesize,breaklines,breakanywhere]
[apa]

\end{Verbatim}

The field is formatted with the character style ``Phonetic form (pronunciation)'', so that you can specify a unique font for this field in the document template for MS-Word.


{\small <Optional>}\par

\subsection{\textbackslash{}pl  plural form}\label{key:pl}

This is a special field used to give the plural form of the lexeme, e.g. in Selaru, the plural form of the lexical entry `asw' (`dog') is entered as:


{\small\color{exlabel}(431)}\par
\noindent
\begin{Verbatim}[fontsize=\footnotesize,breaklines,breakanywhere]
  \pl asure

\end{Verbatim}

MDF adds a ``Pl: `` label to this form. The data is formatted as vernacular text.


For more flexible labeling and to allow for glossing, MDF provides the following set of paradigm fields:


\begin{description}
\item[{\hyperref[key:pdl]{\textcolor{cflink}{\textbackslash{}pdl}}}] paradigm label
\item[{\hyperref[key:pdv]{\textcolor{cflink}{\textbackslash{}pdv}}}] paradigm vernacular form
\item[{\hyperref[key:pde]{\textcolor{cflink}{\textbackslash{}pde}}}] paradigm form-English gloss
\item[{\hyperref[key:pdn]{\textcolor{cflink}{\textbackslash{}pdn}}}] paradigm form-national gloss
\item[{\hyperref[key:pdr]{\textcolor{cflink}{\textbackslash{}pdr}}}] paradigm form-regional gloss
\end{description}


For related fields, see:


\begin{description}
\item[{\hyperref[key:sg]{\textcolor{cflink}{\textbackslash{}sg}}}] singular form
\item[{\hyperref[key:pd]{\textcolor{cflink}{\textbackslash{}pd}}}] paradigm set
\end{description}


{\small <Optional>}\par

\subsection{\textbackslash{}rd  reduplication form(s)}\label{key:rd}

A repository for reduplication forms for later study (conceivably these forms will be moved to subentries, etc. as more is learned about them). For example, the reduplicated form for `yoban' (`hit') in Selaru is:


{\small\color{exlabel}(432)}\par
\noindent
\begin{Verbatim}[fontsize=\footnotesize,breaklines,breakanywhere]
  \rd ioboban

\end{Verbatim}

MDF adds a ``Redup: `` label to this form. The data is formatted as vernacular text.


For information on subentries, see:


\begin{description}
\item[{\hyperref[key:se]{\textcolor{cflink}{\textbackslash{}se}}}] subentry (a polymorphemic form \hyperref[key:or]{\textcolor{cflink}{or}} a phrase)
\end{description}


For information on paradigm fields, see:


\begin{description}
\item[{\hyperref[key:pd]{\textcolor{cflink}{\textbackslash{}pd}}}] paradigm set
\end{description}


For more flexible labeling and to allow for glossing, MDF now provides the following set of paradigm fields:


\begin{description}
\item[{\hyperref[key:pdl]{\textcolor{cflink}{\textbackslash{}pdl}}}] paradigm label
\item[{\hyperref[key:pdv]{\textcolor{cflink}{\textbackslash{}pdv}}}] paradigm vernacular form
\item[{\hyperref[key:pde]{\textcolor{cflink}{\textbackslash{}pde}}}] paradigm form-English gloss
\item[{\hyperref[key:pdn]{\textcolor{cflink}{\textbackslash{}pdn}}}] paradigm form-national gloss
\item[{\hyperref[key:pdr]{\textcolor{cflink}{\textbackslash{}pdr}}}] paradigm form-regional gloss
\end{description}


{\small <Optional>}\par

\subsection{\textbackslash{}rr  reversal form (regional)}\label{key:rr}

This gives a regional language form which could be used to reverse the dictionary (to make a regional language list), or to be included in a national language list marked as the regional language (to aid recall).


\begin{quote}\small\itshape The ability to reverse on a regional form is not currently supported in MDF. This field does not normally print in the dictionary.\end{quote}

{\small <Optional>}\par

\subsection{\textbackslash{}sc  scientific name}\label{key:sc}

Providing a scientific name for a lexeme can be very useful if accurate. The gloss ``a type of tree'' is nearly worthless (but often the best we can do at the time). Having the scientific name allows us to eventually find the appropriate English gloss for such a species. Getting the scientific name requires access to high quality books (color pictures help).


The data given is automatically underlined and italicized by MDF; no formatting is needed.


{\small <Optional>}\par

\subsection{\textbackslash{}sd  semantic domain}\label{key:sd}

This is the English version of \textbackslash{}th and probably the one to use first. Here you try to catalog and differentiate the semantic compartments of an entry, being careful to not let the English force or mask the vernacular relations. Moving to the vernacular terms (given in \textbackslash{}th field) as early as possible is best.


If selected for output, MDF adds the label ``SD: `` to this field.


For related fields see:


\begin{description}
\item[{\hyperref[key:th]{\textcolor{cflink}{\textbackslash{}th}}}] thesaurus
\item[{\hyperref[key:is]{\textcolor{cflink}{\textbackslash{}is}}}] index of semantic
\end{description}


\begin{quote}\small\itshape This field does not normally print.\end{quote}

{\small <Optional>}\par

\subsection{\textbackslash{}sg  singular form}\label{key:sg}

This is a special field used to give the singular form of the lexeme, e.g. in Selaru, the singular form of the lexical entry `asw' (`dog') is entered as:


{\small\color{exlabel}(433)}\par
\noindent
\begin{Verbatim}[fontsize=\footnotesize,breaklines,breakanywhere]
  \sg askwe

\end{Verbatim}

MDF adds a ``Sg: `` label to this form. The data is formatted as vernacular text.


For more flexible labeling and to allow for glossing, MDF provides the following set of paradigm fields:


\begin{description}
\item[{\hyperref[key:pdl]{\textcolor{cflink}{\textbackslash{}pdl}}}] paradigm label
\item[{\hyperref[key:pdv]{\textcolor{cflink}{\textbackslash{}pdv}}}] paradigm vernacular form
\item[{\hyperref[key:pde]{\textcolor{cflink}{\textbackslash{}pde}}}] paradigm form-English gloss
\item[{\hyperref[key:pdn]{\textcolor{cflink}{\textbackslash{}pdn}}}] paradigm form-national gloss
\item[{\hyperref[key:pdr]{\textcolor{cflink}{\textbackslash{}pdr}}}] paradigm form-regional gloss
\end{description}


For related fields, see:


\begin{description}
\item[{\hyperref[key:pl]{\textcolor{cflink}{\textbackslash{}pl}}}] plural form
\item[{\hyperref[key:pd]{\textcolor{cflink}{\textbackslash{}pd}}}] paradigm set
\end{description}


{\small <Optional>}\par

\subsection{\textbackslash{}so  source of data}\label{key:so}

This is a place to indicate the name and village of the informant who gave you the data in the current entry. There is no standard way such information should be encoded.


If selected for output, MDF adds the label ``Source: `` to this field and brackets it with ``[ ]''.


\begin{quote}\small\itshape This field does not normally print.\end{quote}

{\small <Optional>}\par

\subsection{\textbackslash{}st  status}\label{key:st}

This is used to indicate how complete or thoroughly checked an entry is, e.g.:


{\small\color{exlabel}(434)}\par
\noindent
\begin{Verbatim}[fontsize=\footnotesize,breaklines,breakanywhere]
   \st OK             
\end{Verbatim}

{\small\color{exlabel}(435)}\par
\noindent
\begin{Verbatim}[fontsize=\footnotesize,breaklines,breakanywhere]
   \st no-print        
\end{Verbatim}

{\small\color{exlabel}(436)}\par
\noindent
\begin{Verbatim}[fontsize=\footnotesize,breaklines,breakanywhere]
   \st check

\end{Verbatim}

Later you could filter the database to select only ``check'' records, export this filtered database through MDF, and print the formatted output from Word. This makes it easy for a language assistant to systematically check these entries.


If selected for output, the label ``Status: `` is usually added to this field and the whole field is bracketed with ``[ ]''.


\begin{quote}\small\itshape This field does not normally print.\end{quote}

{\small <Optional>}\par

\subsection{\textbackslash{}sy  synonym}\label{key:sy}

This and the \textbackslash{}an (antonym) field are helpful for those who want to keep track of such information without using the \textbackslash{}lf structure.


But you are encouraged to use the \textbackslash{}lf fields to handle ``synonym'' and ``antonym'' cross-referencing (instead of the \textbackslash{}sy and \textbackslash{}an fields). This is because the \textbackslash{}lf fields allow for glossing of the reference, whereas \textbackslash{}sy and \textbackslash{}an do not. Glossing has the advantage of giving the outside reader an idea of the meaning of a referenced lexeme without actually having to go and look it up directly. For example, a synonym of `-haw' (`to pound with a pestle') is `-tutu' (`to pound with a rock'). This could be encoded in the `-haw' entry as:


{\small\color{exlabel}(437)}\par
\noindent
\begin{Verbatim}[fontsize=\footnotesize,breaklines,breakanywhere]
  \sy -tutu

\end{Verbatim}

And would print as:


{\small\color{exlabel}(438)}\par
\noindent
\begin{Verbatim}[fontsize=\footnotesize,breaklines,breakanywhere]
   Syn: -tutu.

\end{Verbatim}

But this tells the reader nothing really. Whereas if this were encoded as:


{\small\color{exlabel}(439)}\par
\noindent
\begin{Verbatim}[fontsize=\footnotesize,breaklines,breakanywhere]
  \lf Syn
\end{Verbatim}

{\small\color{exlabel}(440)}\par
\noindent
\begin{Verbatim}[fontsize=\footnotesize,breaklines,breakanywhere]
  \lv -tutu
\end{Verbatim}

{\small\color{exlabel}(441)}\par
\noindent
\begin{Verbatim}[fontsize=\footnotesize,breaklines,breakanywhere]
  \le to pound with a rock

\end{Verbatim}

This would print as:


{\small\color{exlabel}(442)}\par
\noindent
\begin{Verbatim}[fontsize=\footnotesize,breaklines,breakanywhere]
   Syn: -tutu 'to pound with a rock'.

\end{Verbatim}

Which, to the outsider using your dictionary, is far more helpful.


For related information, see:


{\color{cflink} \hyperref[key:lf]{\textbackslash{}lf} and \hyperref[key:cf]{\textbackslash{}cf}}


See also:


{\color{cflink} \hyperref[key:an]{\textbackslash{}an}}


{\small <Optional>}\par

\subsection{\textbackslash{}tb  table}\label{key:tb}

This marks the following text as unformatted. Line-breaks and hard-coded tab characters will be retained. (Multiple spaces are also be retained, but this is only useful if you define the ``Table'' font in MS-Word to be a fixed width font.)


A table, or list, of ``cutting verbs'' might look like this:


{\small\color{exlabel}(443)}\par
\noindent
\begin{Verbatim}[fontsize=\footnotesize,breaklines,breakanywhere]
 \tb The summary of all types of cutting verbs:
     |{tab}fv:-akrina:|{tab}split in two lengthwise
     |{tab}fv:-boras:|{tab}cut s.t. in small pieces with a knife
     |{tab}fv:-dew:|{tab}chop s.t. into smaller pieces while standing it on end
     |{tab}fv:-het:|{tab}chop or hack with a machete
     |{tab}fv:-kety:|{tab}slice open and clean an animal
     |{tab}fv:-lary:|{tab}slice (like chilies, etc.)
     |{tab}fv:-lilit:|{tab}shave or carve
     |{tab}fv:-mair:|{tab}to adz wood
     |{tab}fv:-simat:|{tab}pop out or cut out coconut meat

\end{Verbatim}

In Toolbox, the table field is marked as a ``No Word Wrap'' field (in the Marker-Properties dialog box). This allows you direct control over line-breaks and spacing. Note the |\{tab\} bar-code in the above example. This tells Toolbox to insert a Tab at each of these places, when the file is formatted by MDF or exported to an RTF file. (Currently Toolbox does not support typing the Tab character directly.)


Inevitably, your tables will require some ``tweaking'' in MS-Word before you print the dictionary. To do this, first convert the file to MS-Word format using MDF, and then search for the ``Table'' style. After some tweaking, these could then be converted into a Word table (with the Word menu command: Table-Convert Text to Table).


For a more powerful and in-depth approach to mapping the relations of lexical items, see:


\begin{description}
\item[{\hyperref[key:lf]{\textcolor{cflink}{\textbackslash{}lf}}}] lexical functions
\end{description}


For more on special codes supported by Toolbox, see:


{\color{cflink} \hyperref[key:Punctuation_and_Special_Codes]{Punctuation\_and\_Special\_Codes}}


\begin{quote}\small\itshape This is a ``free-form'' field.   Punctuation and capitalization should be used as needed.\end{quote}

{\small <Optional>}\par

\subsection{\textbackslash{}th  thesaurus}\label{key:th}

This is a field for developing a vernacular-based thesaurus. It is to be labeled with the vernacular term governing the semantic domain of the entry. Reversing on this field (within Toolbox) would yield a vernacular thesaurus.


If selected for output, MDF adds the label ``Thes: `` to this field.


For related fields see:


\begin{description}
\item[{\hyperref[key:sd]{\textcolor{cflink}{\textbackslash{}sd}}}] semantic domain
\item[{\hyperref[key:is]{\textcolor{cflink}{\textbackslash{}is}}}] index of semantics
\end{description}


\begin{quote}\small\itshape This field does not normally print.\end{quote}

{\small <Optional>}\par

\subsection{\textbackslash{}ue  usage (English)}\label{key:ue}

This field should cover such information as common usage, or restrictions in usage, (such as taboos), or any other information that is needed so a non-native speaker can use this lexeme properly. For example, the Selaru lexical entry for ``wai'' contains the field:


{\small\color{exlabel}(444)}\par
\noindent
\begin{Verbatim}[fontsize=\footnotesize,breaklines,breakanywhere]
 \ue A kin term of address which is used for
     same-sex siblings or for marriable non-kin;
     but may not be said to siblings of the
     opposite sex.

\end{Verbatim}

MDF adds the label ``Usage: `` to this field.


Identical national and regional language fields are:


{\color{cflink} \hyperref[key:un]{\textbackslash{}un} and \hyperref[key:ur]{\textbackslash{}ur}}


Closely related fields are:


{\color{cflink} \hyperref[key:de]{\textbackslash{}de,} \hyperref[key:ee]{\textbackslash{}ee,} and \hyperref[key:oe]{\textbackslash{}oe}}


\begin{quote}\small\itshape This is a ``free-form'' field.   Punctuation and capitalization should be used as needed.\end{quote}

{\small <Optional>}\par

\subsection{\textbackslash{}un  usage (national)}\label{key:un}

The national language equivalent to the \textbackslash{}ue field. This field should cover information such as common usage, or restrictions in usage. For more information, see:


\begin{description}
\item[{\hyperref[key:ue]{\textcolor{cflink}{\textbackslash{}ue}}}] usage (English)
\end{description}


MDF initially adds the label ``NatUsage: `` to this field. This is later changed to whatever national language label is specified for this field in the national audience CC table. For more information on this, see:


{\color{cflink} \hyperref[key:Printed_Field_Labels]{Printed\_Field\_Labels}}


Identical English and regional language fields are:


{\color{cflink} \hyperref[key:ue]{\textbackslash{}ue} and \hyperref[key:ur]{\textbackslash{}ur}}


Closely related fields are:


{\color{cflink} \hyperref[key:dn]{\textbackslash{}dn,} \hyperref[key:en]{\textbackslash{}en,} and \hyperref[key:on]{\textbackslash{}on}}


\begin{quote}\small\itshape This is a ``free-form'' field.   Punctuation and capitalization should be used as needed.\end{quote}

{\small <Optional>}\par

\subsection{\textbackslash{}ur  usage (regional)}\label{key:ur}

The regional language equivalent to the \textbackslash{}ue field. This field should cover information such as common usage, or restrictions in usage. For more information, see:


\begin{description}
\item[{\hyperref[key:ue]{\textcolor{cflink}{\textbackslash{}ue}}}] usage (English)
\end{description}


MDF adds the brackets ``[ ]'' around this field. If included in the output, it is treated as part of the national field.


Identical English and national language fields are:


{\color{cflink} \hyperref[key:ue]{\textbackslash{}ue} and \hyperref[key:un]{\textbackslash{}un}}


Closely related fields are:


{\color{cflink} \hyperref[key:dr]{\textbackslash{}dr,} \hyperref[key:er]{\textbackslash{}er,} and \hyperref[key:or]{\textbackslash{}or}}


\begin{quote}\small\itshape This is a ``free-form'' field.   Punctuation and capitalization should be used as needed.\end{quote}

{\small <Optional>}\par

\subsection{\textbackslash{}uv usage (vernacular)}\label{key:uv}

This field contains the vernacular description of common usage, or restrictions in usage, (such as taboos), or any other information that is needed to describe the lexeme fully. This is intended for use in a monolingual dictionary, but can be used in diglot and triglot dictionaries as well. For more information, see:


\begin{description}
\item[{\hyperref[key:ue]{\textcolor{cflink}{\textbackslash{}ue}}}] usage (English)
\end{description}


MDF initially adds the label ``VerUsage: `` to this field.  This is not changed by either the English or national audience CC table. But you may add a rule to change this if you wish. For more information on how to do this, see:


{\color{cflink} \hyperref[key:Printed_Field_Labels]{Printed\_Field\_Labels}}


Identical English, national and regional language fields are:


{\color{cflink} \hyperref[key:ue]{\textbackslash{}ue,} \hyperref[key:un]{\textbackslash{}un} and \hyperref[key:ur]{\textbackslash{}ur}}


Closely related fields are:


{\color{cflink} \hyperref[key:ev]{\textbackslash{}ev,} \hyperref[key:ov]{\textbackslash{}ov,} \hyperref[key:dv]{\textbackslash{}dv,} and \hyperref[key:gv]{\textbackslash{}gv}}


{\small <Optional>}\par

\subsection{\textbackslash{}va  variant forms}\label{key:va}

This is where variant forms of the lexical entry or subentry can be noted (be they from another dialect or minor alternation in the focus dialect, as in ``do not'' and ``don't''). These variant forms can (but do not have to) refer to minor or minimal entries found elsewhere in the dictionary.


The \textbackslash{}va field heads a bundle of comment fields (\textbackslash{}ve, \textbackslash{}vn, \textbackslash{}vr). These comment fields can be used to specify the dialect name or area that uses the variant form given in the \textbackslash{}va field. But because they are comment fields, you may enter any comment information that you want to appear in the dictionary with the variant form.


MDF adds the label ``Variant:'' to the beginning of the first \textbackslash{}va field in any given section of a lexical entry. The comment fields are added to the variant form with parentheses, but no additional label. Multiple variant bundles are allowed.


The related comment fields are:


{\color{cflink} \hyperref[key:ve]{\textbackslash{}ve,} \hyperref[key:vn]{\textbackslash{}vn,} and \hyperref[key:vr]{\textbackslash{}vr}}


For information concerning minor entries, see:


\begin{description}
\item[{\hyperref[key:mn]{\textcolor{cflink}{\textbackslash{}mn}}}] main entry form
\end{description}


For a detailed discussion of the issues involved in using minor (or minimal) entries, see:


{\color{cflink} \hyperref[key:Using_Subentries_or_Lexical_Entries]{Using\_Subentries\_or\_Lexical\_Entries}}


{\small <Optional>}\par

\subsection{\textbackslash{}ve  variant comment (English)}\label{key:ve}

This is bundled with the \textbackslash{}va field and is where English comments can be given for the variant form. For more information, see:


\begin{description}
\item[{\hyperref[key:va]{\textcolor{cflink}{\textbackslash{}va}}}] variant forms
\end{description}


MDF adds parentheses ``( )'' around this field.


The related comment fields are:


{\color{cflink} \hyperref[key:vn]{\textbackslash{}vn} and \hyperref[key:vr]{\textbackslash{}vr}}


\begin{quote}\small\itshape This is a ``free-form'' field. Punctuation and capitalization should be used as needed.\end{quote}

{\small <Optional>}\par

\subsection{\textbackslash{}vn  variant comment (national)}\label{key:vn}

This is bundled with the  \textbackslash{}va field and is where national language comments can be given for the variant form. For more information, see:


\begin{description}
\item[{\hyperref[key:va]{\textcolor{cflink}{\textbackslash{}va}}}] variant forms
\end{description}


MDF adds parentheses ``( )'' around this field.


The related comment fields are:


{\color{cflink} \hyperref[key:ve]{\textbackslash{}ve} and \hyperref[key:vr]{\textbackslash{}vr}}


\begin{quote}\small\itshape This is a ``free-form'' field.   Punctuation and capitalization should be used as needed.\end{quote}

{\small <Optional>}\par

\subsection{\textbackslash{}vr  variant comment (regional)}\label{key:vr}

This is bundled with the \textbackslash{}va field and is where regional language comments can be given for the variant form. For more information, see:


\begin{description}
\item[{\hyperref[key:va]{\textcolor{cflink}{\textbackslash{}va}}}] variant forms
\end{description}


If this field is included in the output, MDF adds parentheses ``( )'' around it and it is treated as part of the national field.


The related comment fields are:


{\color{cflink} \hyperref[key:ve]{\textbackslash{}ve} and \hyperref[key:vn]{\textbackslash{}vn}}


\begin{quote}\small\itshape This is a ``free-form'' field.   Punctuation and capitalization should be used as needed.\end{quote}

{\small <Optional>}\par

\subsection{\textbackslash{}we  word-level gloss (English)}\label{key:we}

This gives the gloss to be used in word-level interlinear glossing.


Related fields are:


{\color{cflink} \hyperref[key:wn]{\textbackslash{}wn} and \hyperref[key:wr]{\textbackslash{}wr}}


\begin{quote}\small\itshape This field does not normally print.\end{quote}

{\small <Optional>}\par

\subsection{\textbackslash{}wn  word-level gloss (national)}\label{key:wn}

This gives the gloss to be used in word-level interlinear glossing.


Related fields are:


{\color{cflink} \hyperref[key:we]{\textbackslash{}we} and \hyperref[key:wr]{\textbackslash{}wr}}


\begin{quote}\small\itshape This field does not normally print.\end{quote}

{\small <Optional>}\par

\subsection{\textbackslash{}wr  word-level gloss (regional)}\label{key:wr}

This gives the gloss to be used in word-level interlinear glossing.


Related fields are:


{\color{cflink} \hyperref[key:we]{\textbackslash{}we} and \hyperref[key:wn]{\textbackslash{}wn}}


\begin{quote}\small\itshape This field does not normally print.\end{quote}

{\small <Optional>}\par

\subsection{\textbackslash{}xr  translation of example (regional)}\label{key:xr}

This provides the regional translation of the example sentence given in the \textbackslash{}xv field.


MDF adds the brackets ``[ ]'' around this field. If included in the output, it is treated as part of the national field.


For more information about example sentences, see:


\begin{description}
\item[{\hyperref[key:xv]{\textcolor{cflink}{\textbackslash{}xv}}}] example sentence (vernacular)
\item[{\hyperref[key:rf]{\textcolor{cflink}{\textbackslash{}rf}}}] reference to notebooks, texts, etc.
\end{description}


Related fields are:


{\color{cflink} \hyperref[key:xe]{\textbackslash{}xe} and \hyperref[key:xn]{\textbackslash{}xn}}


\begin{quote}\small\itshape This is a ``free-form'' field.   Punctuation and capitalization should be used as needed.\end{quote}

{\small <Optional>}\par

\section{Discontinued}
\subsection{\textbackslash{}xg  (discontinued field)}\label{key:xg}

This field has been discontinued. Originally, MDF reserved the marker \textbackslash{}xg for interlinear glossing, but its function was never fully developed. For this reason, it has been removed from the reserved list of field markers that MDF will recognized.


\begin{quote}\small\itshape If this field is included in a lexical entry, MDF will simply treat it the same as an unknown field.\end{quote}

For information on unknown fields, see:


{\color{cflink} \hyperref[key:Unknown_Fields]{Unknown\_Fields}}


\end{document}
